\vspace{2em}

三个月时光，如白驹过隙。

\vspace{2em}

天枢城，内域七大主城之首，雄踞于浩瀚灵脉交汇之枢，吞吐着整个内域最为精纯磅礴的天地灵气。城墙高耸巍峨，直入云端，墙体上铭刻着无数历经岁月洗礼、依旧闪烁着微光的古老符文。

城中央，便是威震内域的天师府。琼楼玉宇，飞阁流丹，连绵的殿阁在终年缭绕的灵雾霞光中若隐若现，檐角风铃随风清响，涤荡心神。

十年一度的天师试炼，便在这宛如天上宫阙的天师府外门广场所举行。

\vspace{2em}

卯时初，晨光微熹。

天师府外那足以容纳万人的巨型广场，已是人头攒动，气息驳杂而蓬勃。来自内域数百城池、大小宗派、世家以及散修之中脱颖而出的年轻俊杰，足有千余之众，如同百川归海，汇聚于此。空气中弥漫着兴奋、紧张、野心以及淡淡的火药味，绝大多数人脸上都写满激动，眼神灼热地望向广场尽头那巍峨的山门。

广场中央，是一座临时搭建的高台，台上立着一口古朴厚重的黄铜大钟，钟身隐现符文，显然并非凡物。

高台之下，千余名试炼者按照抵达先后或各自的小团体稀疏站立，形成一片黑压压的人潮。其中以玄脉境修士为主流，气息强弱不一；偶有几个周身灵力波动明显更为凝实的，显然是已踏入真符境摹形期的少年天才，在此次试炼中堪称鹤立鸡群，是夺魁的热门人选。

在这片躁动的人潮中，有一小簇人显得格格不入，却又自成一方宁静气象。

叶牧谦一袭朴素青衫，负手立于人群边缘，神情淡泊宁静，仿佛周遭的喧嚣与热切都与他无关。他早已言明，此行只为观礼见证，自身绝不参与试炼：此时这般超然物外的姿态，落在一些有心人眼中，反倒平添几分莫测高深的“高人”风范。

白言冰与陈千羽则一左一右侍立师尊身旁。两人迥异于常人的精气神与从容气度，仍引来不少好奇的打量。白言冰气质沉凝，身姿挺拔如松，目光锐利如剑，但又内敛着一股沉稳的力量；陈千羽则是温婉宁静，眸含仁心，气息柔和，仿佛一泓清泉，能涤去人心浮躁。

沈瑶则站在师徒三人稍靠后的位置，依旧是一身利落的黑色劲装，轻纱覆面，只露出一双清冷如寒星的眸子。她沉默地打量着这熟悉到骨子里、却又因十年漂泊而染上陌生疏离感的场景，面纱下的神色复杂难明。广场四周的高处观礼台上，以及维持秩序的天师府执事弟子中，偶尔有几道目光惊疑不定地扫过她这道身影；显然，部分知晓内情的人，对于这位离家十年的大小姐竟会出现在试炼者行列中，感到极为惊讶与不解。

事实上，沈瑶最初确实打定主意只作壁上观，纯粹充当叶牧谦师徒的向导；然而，一方面是叶牧谦那句“此亦为红尘历练一种，何妨亲身一试”的淡然建议，另一方面，则是三个月前父亲法相降临、目光中那难得一见的、褪去威严外壳的恳求与期盼。在这双重因素下，她终究还是勉为其难地报了名，将自己也置于这千军万马过独木桥的竞争之中。

\vspace{2em}

试炼正式开始前，最后的准备间隙。

白言冰稍稍侧身，压低声音问身旁的沈瑶：“沈姑娘，冒昧请教，这天师试炼，具体是何等流程与内容？还望不吝指点。”

沈瑶从纷乱的思绪中抽离，同样低声回应，带着回忆的口吻。

“若规矩未改，沿用十年前旧例的话……大抵分为三段。”她目光扫向那高耸入云的山门，“‘登天阶’、‘战傀儡’、‘叩心境’。第一关考较根基耐力与灵力持久；第二关测试实战应变与手段；第三关，则是直指本心，拷问心志。具体细则或许会有微调，但核心框架，想来大差不差。”

她话音刚落，高台之上，数位身着玄黑色底、绣有银色云雷纹路道袍的老者，依次在台上现身；几人目光开阖间自有威严，显然皆是天师府中地位尊崇的长老。居中一位，尤为引人注目，他须发皆白如雪，面容却红润如初生婴孩，不见丝毫老态，一双眸子开阖时，隐约有细碎电光流转，令人不敢直视。

白发长老举起右手，灵气立即在他手中凝聚成槌，正是真符境的特征——灵气化物。

槌击铜钟，广场上本来的窃窃私语立即静默下来。

白发长老松手时，灵气槌立即散去。

“本届天师试炼，即将开始。老夫即为本次主考。现在，宣读试炼细则。”

广场落针可闻，千余道目光齐刷刷聚焦于高台。

“试炼分三关，相连进行，需自行把握节奏。”长老语速平稳，却字字千钧，“第一关：攀登三千级石阶！”

众人不由得顺着长老所示，望向山门之下那蜿蜒向上、没入云雾之中的青石长阶。石阶古朴，看上去与寻常山道并无二致。许多人不免心生疑惑：这有何难？莫说修士，便是强壮些的凡人，咬咬牙也能攀上三千级。

长老似乎早已预料到众人的反应，淡淡补充：“攀登之时，不得服用任何丹药，不得借助符箓、法器之力，亦不可御空飞行。提醒诸位，此阶自有其不凡之处。限时半个时辰，未按时抵达山门平台者淘汰。”

原来如此！不少人恍然，又暗生警惕。这第一关，竟是纯粹考较修士自身根基耐力、灵力恢复与持久作战能力的硬骨头。那石阶，也自然绝非看起来那么简单。

“第二关，”长老继续道，“于山门平台处进行。战胜或突破一具守护山门的傀儡，并完好无损地取回其胸口的宝石！”他略略提高了声调，“此关，法宝、丹药、阵法等手段，只要是你自身所有、能施展得出，皆在规则允许之内！唯有一点，宝石若损，即算失败。同样，限时半个时辰。”

此言让不少擅长战斗或身怀异宝的修士精神一振，眼中露出跃跃欲试之色。实战关卡，正是他们展现实力的舞台。

“第三关，”长老的声音陡然变得更为沉凝，带着一种直透人心的力量，“穿过山门，步入‘叩心幻境’。幻境之中，彼此隔绝，不见他人，唯见己心。”他顿了顿，目光似乎变得深邃起来，“届时，府主一丝神识威压亦会降临幻境。非心志坚如磐石、城府深沉似海者，于幻境与威压之中，诸般思绪，纤毫毕现，无所遁形。此关，不限时长，走完幻境，即是终点，亦为试炼通过之标志。”

“府主威压？！” 当“沈天穹威压”这几个字从长老口中吐出时，台下千余人中，顿时响起一片压抑不住的吸气声与骚动！无数张年轻的面孔上，神色骤变，敬畏、紧张、兴奋、恐惧……不一而足。

那可是法相境大能的威压！即便只是一丝神识，对于他们这些最高不过真符境的年轻人而言，也如同直面苍穹，其中心神冲击与压迫，可想而知。这最后一关，才是真正筛选心性、决断谁能承载“小天师”尊号的炼心之关！

高台上的长老似乎颇为欣赏台下众人的各异神色，略作停顿，才继续补充道：“稍后，会有弟子发放应急符纸一人一枚。试炼过程中，若自觉不支，或遇生命危险，撕破符纸，即会有附近长老出手接引，同时，也意味着自动弃权。”

“此外，铁律重申：试炼期间，参与者严禁相互攻击、私下斗殴、扰乱秩序！违者，立即逐出试炼，永不录用！”最后一句话，斩钉截铁，带着冰冷的肃杀之意，让所有人心中一凛。 “现在，请所有不参与试炼的观礼道友，依序退出广场区域。”

规则宣读完毕，白发长老不再多言，径自在高台中央的蒲团上坐下，闭目养神，仿佛接下来的一切都与他无关。

随即，一批天师府弟子鱼贯入场，手中托着托盘，其上整齐码放着淡黄色的符纸。只见他们口中念念有词，手指轻弹，那些符纸便如同被无形之手牵引，主动而精准地飞向场内每一位试炼者，被他们或探手、或用灵力接下。整个过程井然有序，不到半刻钟便已完成。

发放符纸的弟子迅速退场，偌大的广场上，只剩下千余名屏息凝神、蓄势待发的试炼者，以及高台上闭目的长老和那口沉默的黄钟。空气仿佛凝固，只剩下晨风吹拂衣服的微响，和彼此间越来越清晰可闻的、加速的心跳声。

所有人都在等待，等待那一声象征着开端与竞争的钟鸣。

\vspace{2em}

辰时正。高台上，闭目的白发主考长老倏然睁眼，眸中电光一闪而逝。他长身而起，再次凝出那柄灵气槌，对着身旁的黄铜大钟，沉稳而有力地敲下！

“咚——！！！”

洪亮、厚重、悠远的钟鸣，如同积蓄了千年的力量骤然释放，以高台为中心，轰然荡开！声波肉眼可见地扩散，撞击在广场四周的殿宇楼阁上，又折返回来，在群山之间反复回荡、碰撞、叠加，余音袅袅，经久不息，仿佛在为这场盛会奏响古老而庄严的序曲。

钟鸣入耳的刹那，那早已如满弓之弦、蓄势到极致的千余名年轻修士，体内压抑的灵力与战意如同火山喷发！

人群轰然炸开，化作一道道颜色各异、迅疾无比的流光，朝着那三千级青石发起了冲锋！

人影憧憧，衣袂破风之声尖锐刺耳，灵力爆发的光芒瞬间将清晨的薄雾染得五彩斑斓。空气中充满了灼热的竞争气息、灵力碰撞的微鸣，以及奔跑踏地的轰鸣，仿佛一场盛大而残酷的生存竞赛，于此拉开血淋淋的序幕。

石阶以不知名的古朴青石铺就，历经风雨，表面光滑，透着岁月的沉静。然而，当第一只脚踏上第一级石阶的瞬间，一股沉重如山、粘滞如胶的无形压力，毫无征兆地自脚下石板中升腾而起，如同拥有生命般，瞬息间沿着腿脚蔓延，席卷全身！

这压力竟然与当初叶牧谦的见独领域在原理上有几分异曲同工之妙。

它不仅作用于肉身，也渗透肌肤，侵入经脉窍穴，甚至直接撩拨、冲击修士的心神与灵觉！

骤然之间，体内原本流畅运转的灵力，像是被掺入了无数细沙，陡然生出一股滞涩之感；心跳不由自主地加速，呼吸也开始变得有些紊乱。

第一步，便是下马威！

然而，能来到这里的，皆是内域年轻一辈的佼佼者，对此类考验或多或少有所耳闻或准备。最初的惊愕只持续了短短一瞬。

“轰！轰！轰！……”

各色璀璨夺目的灵力光华，如同黑夜中骤然点燃的万千火把，几乎在同一时间，从冲锋的人群中疯狂爆发而出！

有人周身腾起赤红如血的烈焰，高温灼烧空气，踏步间在湿润的石阶上留下嗤嗤作响的焦痕，仿佛火神降世，狂暴突进；

有人通体泛起金属般的冷硬灰白光泽，肌肉贲张，脚步踏下竟发出金铁交击般的铿锵之声，势大力沉，似要碾碎阶石；

还有人身上笼罩着厚重沉凝的土黄色灵光，步伐不算迅猛，却每一步都踏实无比，展现出惊人的稳定性……

这便是“灵枢径”体系下，最为普遍、也最为直接的应对方式——将吸纳而来的天地灵气，作为迅猛燃烧的薪柴，以自身功法为熔炉，全力催动，猛烈爆发，用澎湃的灵力洪流强行冲开、抗拒乃至抵消外部势场的压力。

他们将灵气疯狂灌入双腿相关经脉，强化肌肉、骨骼、筋膜，试图减轻负担，提升速度。

一时间，漫长的青石阶上，流光溢彩，灵力澎湃，一道道身影如离弦之箭、逆流之鱼，你追我赶，怒吼着、呐喊着，向上方发起冲锋，场面壮观激烈，充满了最原始的力量碰撞之美。

\vspace{2em}

然而，在这股色彩斑斓、争先恐后、喧嚣震天的灵力爆发洪流之侧翼，却有两个人的身影，显得如此格格不入，沉静得近乎诡异。

白言冰与陈千羽对视一眼，彼此眼中皆有默契与了然。

白言冰只是深深地吸了一口气，他体内那源自一身温养而出的元气，开始自然而然地加速循环流转起来。

元气如同他身体内自然运转的江河溪流，浸润着四肢百骸、五脏六腑，以一种圆融和合的方式，悄然适应、引导、乃至缓慢化开那从四面八方、无孔不入挤压而来的无形阻力与压力。他的步伐，初看之下并不快，甚至对比周围那些火箭般的突进身影，显得颇为缓慢与保守。

然而，他的每一步都踏得极稳，脚掌与石阶接触的瞬间，重心转换流畅自然，身形没有丝毫的摇晃与滞涩。他的呼吸，从一开始就保持在一个悠长、稳定、富有节奏的状态，仿佛不是在参加一场残酷的淘汰赛，而是在进行日常的吐纳修炼。

身旁的陈千羽，元气性质更为温润、柔和，隐隐带着滋养万物生机的意蕴。她行走间，自身的气息似乎与周围山间掠过的微风产生了某种难以言喻的微妙共鸣与交融。那石阶施加的强大压力落在她身上，仿佛被一层无形而柔韧的薄膜缓冲、分散、吸纳了一部分。她的步履因此显得比白言冰还要轻盈几分。

\vspace{2em}

沈瑶一开始确实跟在两人身后不远处；但眼见前方人潮汹涌、争先恐后，她清冷的眸中闪过一丝锐色，一股狠劲从心中生发而出。

随即，体内精纯的淡青色灵力瞬间汹涌，包裹全身。她施展出精妙的身法，整个人立即飘忽不定、迅疾灵动起来，在拥挤不堪、互相推搡的人潮缝隙中精准而快速地穿梭，凭借高超的身法和对灵力爆发时机的精妙控制，很快便超越了稳步前行的白言冰与陈千羽，冲到了最前方第一梯队的位置。

然而，就在她于一次腾挪间隙，下意识地回头一瞥，望向被甩在身后的那两道身影时，心中却莫名地、剧烈地动了一下。

她看到的是什么？她自己，以及身边前后左右几乎所有的修士，都在奋力地与别人抗争——面目或狰狞，或紧绷，周身灵力鼓荡，光芒闪烁不定，仿佛在与无形的巨人进行着一场角力，每一分力量都在嘶吼着向外迸发，消耗剧烈而明显。

可那两人…… 白言冰目光平视前方，神色专注而平静；陈千羽甚至有余暇观察两侧山景。他们的从容，绝对是一种由内而外、发自根本的协调与从容。

这种无比鲜明的对比，如同一柄无形的小锤，轻轻敲击在沈瑶内心深处某个根植了二十多年的认知根基上。一丝难以言喻的异样感与困惑，悄然滋生。

不过，眼下爬楼梯要紧。

困惑一闪而过，她咬咬牙，收敛心神，将更多灵力灌注于双腿，继续以更快的速度向上冲刺。

现在不是深究的时候。

\vspace{2em}

时间是公正而残酷的审判者。

随着众人攀登的高度不断增加，石阶传来的那种全方位压迫感，开始以清晰可感的速度层层递进、不断增强。

更麻烦的是，那缠绕在灵力运转上的干扰之力也愈发明显、刁钻。石阶本身，仿佛一个巨大而贪婪的灵力吸收与扰乱场，持续地吸纳、扭曲着周围的天地灵气，使得修士们试图在行进中吐纳补充灵力的效率大打折扣，甚至事倍功半。

最初那批凭借澎湃灵力爆发，一骑绝尘、冲在最前面的修士，最先尝到了苦果。

“嗬……嗬……嗬……” 粗重如破损风箱般的喘息声，开始从队伍前列传来。一个原本周身赤红烈焰最为炽盛、冲在最前方的魁梧汉子，此刻面色潮红如血，额头上、脖子上青筋暴起如同虬龙，眼神中最初的狂猛已被剧烈的疲惫取代。

他步伐踉跄，每一步都踏得沉重无比，周身的火焰灵光明灭不定，闪烁摇曳，仿佛随时都会熄灭——这正是灵力即将难以为继、濒临枯竭的显著征兆。

“可恶……这……这鬼石阶！怎会……吸纳灵力如此之快？！”他低吼一声，满是不甘与愤懑，却不得不强行停下冲锋的脚步，试图原地调息，吸纳灵气补充。

然而，他绝望地发现，平日里如臂使指、随意吐纳的天地灵气，此刻却如同陷入了厚重粘稠的泥沼，汲取起来异常艰难、缓慢，根本赶不上他可怕的消耗速度。

类似的情况接二连三地出现。金属光泽迅速黯淡，土黄灵光涣散不稳，冰蓝气息摇曳欲熄。

不少原本冲在前列、气势汹汹的身影，速度骤降，有的甚至如同被抽干了力气，只能坐倒在石阶上休息片刻。满脸写着不甘、憋屈与深深的疲惫，只能眼睁睁地看着后来者——那些最初被他们甩开、此刻却仍能保持一定速度的人——面无表情或带着同情地从自己身边一步步超越。

灵力，如同借来的滔天洪水，初时挟山超海，势不可挡，声势骇人。但其“借”的特性，决定了它难以持久：一旦后续补充跟不上狂暴的消耗速度，那么后劲不济便是必然的结局。

反观白言冰与陈千羽，他们的速度，自始至终都维持在一个稳定的区间。

元气在体内循环，参与抵抗压力、推动身体前行，固然也在持续消耗；

但问道途的精髓在于，元气是润滑油，不是燃料，损失极为可控；微微消耗损失的同时，脏腑经络本身也在持续不断地温养、滋生着新的元气。

他们对自身每一分力量的掌控都精微入里，力量的使用高度协调，几乎没有丝毫浪费在无谓的、与外部压力的硬性对抗上。他们就像逆流中沉稳无比的两块礁石，任由身边光影变幻、人群喧嚣、超越与被超越，他们只是按照自己内在的节奏与韵律，一步，再一步，稳健、扎实、匀速地向上前行。

渐渐地，他们以一个平稳的速度，超越了一个又一个后劲不足、不得不停步调息、或步履维艰的修士。那些被超越者，无不向两人投来惊愕、不解、疑惑，乃至一丝茫然的目光。

他们看不懂，为何这两个看起来灵力波动微弱、也没见如何爆发的家伙，却能如此平稳地持续前进？

沈瑶同样感受到了越来越巨大的压力。淡青色的灵力光晕在她体表明暗剧烈地闪烁，如同风中之烛。经脉中传来阵阵隐隐的刺痛，那是灵力高速运转却又补充不及、开始透支经脉的征兆。

她牙关紧咬，面纱下的嘴唇抿成一条苍白的直线。

长痛不如短痛！

凭借远超同侪的扎实根基、十年江湖漂泊锤炼出的坚韧意志与丰富的耐力分配经验，硬生生地扛着，以比白言冰他们更快、但也更消耗的方式，向上冲刺。

当她用尽最后一股气力，抬脚踏过第三千级、也是最后一级石阶，双足终于落在山门前端那宽阔平坦的平台上时，一阵强烈的虚脱感如同潮水般袭来，让她眼前微微一黑，身躯晃了几晃，险些站立不稳。

她汗水早已浸透了额前的发丝和贴身的劲装，胸口剧烈地起伏着，心脏如同战鼓般狂跳，几乎要撞出胸膛，紊乱的气息在经脉中乱窜。虽然这种“透支”在训练或作战中都是很常见的，但事后及时的恢复也是极为必要的。

她喘息着，往口中噙了一枚丹药，咚的一声坐下开始调息。竭力控制、多次深长呼吸吐纳，才勉强从透支状态中恢复过来，像是沙漠中跋涉许久的旅人，终于找到绿洲、喝到了一口甘甜的水。

回头望去，来时那密密麻麻、拥挤不堪、如同长龙般蜿蜒的石阶上，此刻景象已然大变：人群变得稀疏了许多，仍有不少身影在艰难地、缓慢地向上挪动，但更多的人则停留在中途，或坐或躺，面色灰败，显然已无力继续。

甚至，她看到了几道身影正面色惨淡、一步三回头地调头向下走去，或者颓然撕碎了手中的黄色符纸——那意味着放弃与淘汰。

粗略估算，能在这第一关半个时辰内成功登顶的试炼者，恐怕已经不足五百之数，淘汰率超过一半；而即便是这成功登顶的五百人中，像她这般几乎耗尽气力、需要时间恢复的恐怕占了大半！

真正还能保留相当余力，可以立刻投入第二关更激烈战斗的……恐怕还要再少上一大截。

而几乎就在她终于休息差不多的时候，白言冰与陈千羽也稳稳地踏上了平台。

两人额际确实都见了一层细密的薄汗，在晨光下闪着微光，面色也因持续的运动而微微泛红；但他们的呼吸，却在登上平台后，仅仅经过短短几次深长的吐纳，便迅速趋于均匀、绵长、平稳。

那姿态，仿佛刚刚结束的并非一场残酷淘汰近半对手、对耐力和灵力控制要求极高的长途攀登，而仅仅是一次强度适中的晨间登山锻炼。

这份在明显的高强度消耗之后，所展现出的快速恢复能力与始终如一的稳定从容，与平台上许多如同沈瑶一般气喘吁吁、形象略显狼狈的过关者，以及那些直接瘫坐甚至躺在地上、需要服用丹药调息的人，形成了再鲜明不过的对比。

这一幕，让沈瑶胸口那抹异样感与困惑，骤然化为了实实在在的惊意与震动。某种根植于她过往二十年认知最深处的、关于修行、力量、效率的理所当然，在这一刻，似乎终于被眼前的铁一般的事实，悄然撬开了一道清晰而深刻的缝隙。

\vspace{2em}

平台尽头，巍峨的山门如同一位亘古的巨人，沉默地俯视着这群疲惫不堪的闯关者。

山门由整块的玄色玉石砌成，高逾十丈，在稀薄的天光下泛着幽冷深沉的光泽，门楣上以古老篆体铭刻的“天师府”三个大字，每一笔每一划间都有灵光如活水般流转不息。

然而，此刻牢牢攫住所有人目光与呼吸的，是门前列阵以待的金属守卫。

左右两侧，各自矗立着整整九尊高达丈余的金属傀儡。它们形制古朴厚重，仿佛自上古战场中直接搬移而来，通体由某种暗沉无光的特殊合金铸造。

这些傀儡形态各异，有的手持门板般宽阔的刃口已有些许磨损痕迹的巨剑，有的提着浮雕狰狞兽首、刃口泛着寒光的斧钺，还有的横握长戟，戟尖斜指地面。虽静立不动，但一股混合着无形煞气的、百战余生的肃杀之意，已如冰冷的潮水般弥漫了整个平台，压得人胸口发闷。

更加强大的灵力波动自它们沉重的躯壳内隐隐散发，强度赫然都已达到了玄脉境融贯期巅峰的水准，足以让在场不少修士感到呼吸困难。

而在每一尊傀儡胸膛正中央，一枚拳头大小、剔透如凝固鲜血的红色宝石，被刻意安置在毫无遮挡的凹槽内，幽幽地散发着妖异光芒。

正常的战斗傀儡，其驱动核心无不深藏于重重防护之下。眼前这些傀儡的设计，却将核心如此暴露，似乎天师府在有意降低这第二关的纯粹武力难度。

击败或突破一尊守卫傀儡，并完好无损地取走其胸口那枚核心宝石，方能穿过眼前这座山门。

“不必担心傀儡数量不够。”不知从何处传来的长老声音，适时响起，打破了因凝重压力而产生的死寂，“虽说先到先得，但我们自会适时补充。”

平台上，刚刚经历三千级石阶上那无所不在的势场对灵力与体力的双重残酷消耗的修士们，脸色都变得极为难看。不少人早都从怀中掏出各色瓷瓶玉盒，倒出香气各异的丹药囫囵吞服，争分夺秒地试图恢复那已然见底的状态。

没有人立刻上前。似乎不少人先前已从各种渠道得知了消息或前辈经验，只要不踏入傀儡前方约三丈的警戒区域，这些冰冷沉寂的杀戮造物便毫无反应。

一时间，平台上只剩下粗重不一的喘息声、丹药在喉间滚动的声音，以及偶尔响起的、压得极低的急促交谈。

白言冰与陈千羽登上平台后，见众人都按兵不动，于是同样在边缘处找了块相对空旷的地方，同样坐下，闭目调息。他们深谙“螳螂捕蝉，黄雀在后”之理，谁也无法预料这些看似独立的傀儡是否会群起围攻；不如静观其变，以逸待劳。

压抑的、混合着焦虑与紧张的寂静，持续了约莫一刻钟。每一秒都仿佛被拉长，平台上弥漫的除了肃杀之气，还有逐渐累积的、一触即发的躁动。

终于，一名身材魁梧宛若铁塔、背负一柄骇人巨斧的赤膊大汉率先按捺不住，猛地站起身，声如平地闷雷炸响：“怕个鸟！不过是些死物，待我取它核心！”

话音未落，他周身赤红色灵光如同火山喷发般暴涌而出，裸露的肌肉块块贲张，青筋如蚯蚓蜿蜒，整个人当真如同从蛮荒时代冲出的凶兽，将地面踏得微微震颤，悍然冲向距离他最近的一尊持戟傀儡！

就在他粗壮的脚踏入那无形警戒范围的刹那——

“嗡——！！！”

那尊原本如同雕塑般静止的持戟傀儡，眼窝深处那两点猩红光芒骤然炽亮！如同一对被瞬间点燃的血色炭火。低沉的、仿佛来自金属脏腑深处的震鸣轰然响起，打破了长久的寂静。

傀儡沉重庞大的身躯，以一种与其体型绝不相符的、令人心悸的迅捷猛然启动！关节转动处传来“咔嚓、咔嚓”的刺耳摩擦声，手中那杆乌沉长戟，已化作一道撕裂空气的死亡乌光，带着凄厉到极点的尖啸，直刺赤膊大汉那肌肉隆起的胸膛！速度之快，竟在戟尖后方拖出了一条短暂的白气甬道！

战斗，或者说，单方面的挑衅与激活，在这一瞬间被彻底点燃！

仿佛接到了无声而统一的指令，平台上那压抑、等待已久的战意与求生欲轰然爆发，如同被点燃的火药桶。

“上啊！”

“抢！”

各色灵光再次冲天而起，比第一关攀登时更加驳杂、更加狂野。呼喝、怒吼、长啸之声震耳欲聋，汇聚成一片混乱的声浪。早已选定目标或临时就近扑向猎物的修士们，人影纷乱，如同炸窝的蜂群，扑向那些金属守卫。

霎时间，这山门前原本空旷寂静的平台上，金铁交鸣之声大作，灵气疯狂对撞，成了一个缩小而残酷的修罗战场。

灵枢径修士们五花八门、各具特色的手段在此刻展现得淋漓尽致：

有专精剑道的修士长啸一声，意在剑先，剑势吞吐不定，锋芒锐气延伸出足有数尺，每一次与傀儡手中巨剑的铿锵对撞，都爆开大蓬耀眼的火花和刺耳的刮擦声，令人牙酸。

有擅长法术的修士面色凝重，双手掐诀快得留下残影，口中咒文吟唱如疾雨打芭蕉，瞬息间身前便凝聚不同的法术，如同狂风暴雨般连绵轰击向傀儡的膝弯、肘关节、颈项连接处。

更有三五名本就相熟的修士，在混乱中迅速靠近，结成一个简易却实用的三角或圆阵。而这，显然也在天师府允许的规则范围之内——但这样也只能拿到一个核心，事后分配恐又成问题。

沈瑶深吸一口气，压下因攀登三千阶及目睹白、陈二人异常从容而带来的双重波澜。她玉手在身后一抹，一道清冷如秋水的光华闪过，那柄窄刃短刀便已紧握在手。

随即，身法展动，黑衣如夜，化作一道淡青色的凌厉残影，主动迎向一尊刚刚将一名修士连人带盾劈飞出去的手持巨剑傀儡。

她的战斗风格，根基是典型的、千锤百炼的灵枢径天才路数：灵力澎湃，运转流畅，与手中宝刀共振共鸣；但十年的仗义行侠，在无数次真实的生死搏杀中，又将这根基打磨出了独属于她自己的、追求极致效率与一击毙命的冰冷风格。

那傀儡也注意到了沈瑶，门板一般的巨剑挟着万钧之力劈落！

沈瑶却身法迅捷、形似鬼魅，劲风吹拂她面纱一角的千钧一发之际，身形才如同没有重量般轻轻一滑，以毫厘之差精准避开。

随即，在傀儡因招式用老而产生的微小收招间隙，她的刀尖已如等候已久的毒蛇，倏然刺出！

精准、狠辣、毫不留情，直刺傀儡胸口宝石凹槽周边那些用于固定宝石的、更为精细脆弱的金属连接结构与符文刻痕交界处！

“叮！叮！铛！嗤——！”

一连串密集如暴雨打萍的撞击与切割声在数息之内爆发！最终，伴随着一声轻微的金属崩裂脆响，那枚剔透的红宝石被刀尖巧妙一挑，化作一道虹光飞起，被她左手稳稳接住。

几乎同时，傀儡眼中炽亮的猩红光芒如同被掐断的烛火瞬间熄灭。高举巨剑的庞大身躯彻底僵直，化作一座真正的金属雕像，唯有刀身上传来的反震之力，让沈瑶持刀的手腕感到微微酸麻。

沈瑶气息微促，光洁的额角渗出细密汗珠，但整个动作行云流水，干净利落至极，已率先过关。

一招制敌！令人惊悚！

她握着尚带一丝金属余温的宝石，走向早已在山门内侧阴影中静候的白发长老。

长老假装不认识沈瑶，目光扫过那完好无损的宝石，与身后只是暂时停机的傀儡，赞许地微微颔首：“不错。只取核心，不伤傀儡根本，省却不少修复功夫。如此懂得节制与精准的修士，近年少见。”

随即接过宝石，侧身让开通路。

沈瑶没有多言，只是在此刻，才真正抬起眼，深深地、目光复杂地看了一眼山门之上那“天师府”三个古篆大字，仿佛要将它烙印进心底，随后才决然转身，踏入山门后的氤氲雾气之中，身影迅速被吞没。

那长老也不耽搁，待沈瑶身影消失，便双手抬至胸前，掐动一个古朴诀印，口中念念有词。只见他手中那枚红宝石仿佛被无形之力唤醒，微微震颤，随即“嗖”地一声化作一道血色流光，自动飞回那刚刚被沈瑶击败的傀儡胸口凹槽处，严丝合缝地重新嵌入。

紧接着，那傀儡躯体内传来一阵低沉的“嗡嗡”共鸣，眼中熄灭的红芒“轰”地一声再次炽盛亮起，关节处发出“喀啦啦”的复位声响。它微微调整了一下姿势，重新恢复了那待机般的肃立状态，静静等待着下一位挑战者的到来。

\vspace{2em}

白言冰全程看完了沈瑶从观察到爆发、再到精准取宝的整个过程，眼中不禁流露出不加掩饰的佩服之色。随即迅速压下杂念，目光冷静如冰，再次扫过喧嚣混乱的战场，最终，锁定了一尊刚刚击退两名联手修士的持刀傀儡。

这尊傀儡身形在众傀儡中略显修长，手中长刀也非宽厚型，而是略带弧线，刀法风格大开大合，每一次挥斩都带着凄厉的风啸，威势十足，逼得周围修士不敢近身。但白言冰敏锐地察觉到，在两次挥刀动作的转折衔接之间，那刀势会有一个极其短暂的凝滞与调整——这是战斗傀儡的通病。

他模仿着沈瑶最初的策略，开始绕着那尊持刀傀儡，在一个相对安全的距离外游走，紧紧盯住傀儡的肩、肘、腕、腰、膝，观察它每一次举刀的角度、发力的根源、刀势的轨迹，尤其是那“换气点”出现时，其躯干上哪些关节的符文会有一丝异常的能量流转波动。

这些傀儡似乎只有力量达到了玄脉境融贯期巅峰的水准，但战斗逻辑显然过于简单，无法理解这种“窥探”；但它很快便锁定了这个在周围晃动的目标，猩红的目光陡然聚焦在白言冰身上。

它沉重的金属脚掌踏前一步，震得地面微颤，手中那柄长刀已挟着足以将巨石一刀两断的凄厉风声，化作一片扇形的死亡光幕，拦腰横斩而来！刀未至，那挤压空气形成的凌厉风压已割得人皮肤生疼。

白言冰眼神一凝，瞳孔微微收缩。千钧一发之际，他做出一个看似违背常理的选择：向前小踏半步！

这种“向前避让”险之又险，冰冷的刀刃几乎贴着他的胸腹衣襟掠过，劲风将他的衣袍撕开一道小口。

就在这避让的同一瞬间，剑光如电光石火，在傀儡收刀回力的那一丝凝滞出现的刹那，精准无比地刺向傀儡持刀手臂的肩关节处。

“叮——！”

一声清脆却异常刺耳高亢的撞击声爆开。傀儡庞大的身躯随之微微一震，即将发动的下一记斜劈动作出现了明显可见的、超过预期的迟滞。

就是现在！

长剑回旋，不带丝毫烟火气，再次闪电刺出！目标，仍是那处肩关节，同一个点位！

“咔嚓！”

一声细微却清晰无比的碎裂声，从傀儡肩关节深处传来！仿佛内部某个精巧的机关，终于不堪重负，崩裂开来。傀儡持刀的整条右臂顿时一僵，随即变得歪斜、扭曲，后续的劈砍动作完全变形，威力大减，破绽百出。

那剩下的战斗内容，就跟玩一般了。

“铿！锵！嗤啦——！”

令人牙酸的金属摩擦与切割声持续了不到十息。终于，在一阵剧烈的、仿佛内部齿轮彻底卡死的沉闷摩擦声后，这尊凶悍的持刀傀儡全身一震，所有动作戛然而止。

白言冰走上前，伸手从那僵立傀儡胸口的凹槽中，稳稳取出了那枚尚带着一丝余温的红宝石。

远处看完全程的长老此刻倒是不太高兴。

虽然这位小友剑法了得，能精准破坏傀儡的防御，但他是专门负责修这个的啊！

你这么破坏，让我很难修啊！

更要命的是他这么操作是没有违规的！

\vspace{2em}

另一边，陈千羽选择主动迎上的，竟是十八尊傀儡中看起来防御最为严密的之一：一尊左手持着一面足以遮蔽大半个身形的长方形巨盾、右手握一柄短柄却锤头狰狞战锤的傀儡。这尊傀儡往那里一站，便如同一座可移动的金属堡垒，给人一种无从下手的窒息感。

素手轻扬，指缝间悄然夹住了十数枚钢针之后，瞬息间已切入那尊堡垒傀儡的攻击范围。

“呜——！”

战锤挟着恶风，以砸碎山岳之势当头轰落！同时，那面巨盾微微一斜，封死了大部分闪避空间，盾缘闪烁着寒光，显然亦是一件攻防一体的凶器。

陈千羽也是毫不防御，指间蓝芒连闪！

“咻！咻！咻——！”

和白言冰同样的思路，钢针直接刺向傀儡的诸多关节，果然是“不是一家人，不进一家门”。

针体本身纤细无比，不足以破坏傀儡坚固的金属构件。但其上附着的气，却顺着灵力传导的路径钻入，瞬间干扰其内部那些微小的管线。

效果立竿见影！

傀儡那势在必得的一锤砸下，却因为右腿膝关节处传来的瞬间僵直与不协调，导致整个庞大的身躯微微一晃，锤头落点偏了半尺，狠狠砸在陈千羽身旁的地面上，碎石飞溅，砸出一个浅坑。而几乎同时，它试图抬起巨盾进行下一步格挡或撞击的动作，也因左臂肩轴处的传导不畅，明显慢了半拍，盾牌抬起的高度不足，侧面露出了一线破绽。

这破绽转瞬即逝，但对陈千羽而言，已足够宽阔！

对陈千羽而言，现在这个傀儡和一个鲁班锁已经没有区别了。

她眼眸沉静如水，手指翻飞如蝶，一枚枚蓝汪汪的钢针持续射出，每一次都精准命中傀儡背部、后颈、腰眼等处的关键连接点。

“咻！嗤！叮……”

傀儡的动作开始变得滑稽而混乱，如同提线木偶断了部分丝线。它试图转身，双腿却不听使唤地打着绊子；想要挥锤向后横扫，手臂却只抬起一半便无力垂下；巨盾更是沉重地拖曳在地，火星四溅。它的动作幅度越来越小，越来越迟缓。

当最后一枚闪烁着幽蓝寒光的钢针没入傀儡后颈处时，这尊铁塔般的堡垒守卫全身猛地一颤，随即所有动作彻底停止。它左手盾牌微抬，右手战锤半举，身躯半转，凝固成了一个充满张力却又无比古怪的金属雕塑。

唯有猩红的眼睛和胸口处那枚红宝石，依旧在不疾不徐地闪烁着妖异的光芒。

陈千羽绕回傀儡正面，指尖在宝石边缘几个精巧的卡扣处轻轻一拨、一按，那枚红宝石便听话地弹起，落入她早已备好的掌心，温润微凉。

当陈千羽转身走向山门时，平台上那震耳欲聋的喊杀声、灵力爆鸣声以及金铁交击的轰鸣，已渐渐稀落下去，变得零散而疲惫。各色耀眼的灵光大多已经黯淡熄灭，如同潮水退去后裸露的礁石。

放眼望去，平台之上还能站立的修士已然不多。成功击败傀儡、手持那代表资格的红色核心宝石，能够走到山门前等待进入下一关的修士，粗粗一数，区区三十人而已！

\vspace{2em}

穿过巍峨山门，映入眼帘是一片无边无际、浓得化不开的白茫茫迷雾。

这雾仿佛有生命般缓缓流淌，隔绝了视线，也吞噬了所有声音，只剩下一种令人心头发紧的绝对寂静。

迷雾之前，历经前两关残酷筛选后仅剩的不到三十名试炼者聚集于此，人人脸上都带着难以掩饰的疲惫与深入骨髓的警惕，空气凝重得几乎能拧出水来。

先前白发长老朗读过的通知内容，此刻如同冰冷的咒语，反复在他们脑中回响：所有伪装、所有依靠外物或他人比较建立的心理优势都将失效，必须赤裸裸地面对自己内心深处最真实的想法，还要抵抗法相境大能无形的心灵压迫！

这意味着，无论你出身何等显赫，拥有何等秘宝，修炼了何等奇功，在这里都毫无意义。你将面对的，唯有你自己——那个可能连自己都不敢直视的、真实的自己。

没有退路。众人相视之间，只能从彼此眼中看到同样沉重的无奈与决绝。他们深深吸气，硬着头皮，一个接一个，如同赴死般，陆续走入那吞噬一切的浓雾之中。

幻境之中，果然如长老所言，每个人仿佛被投入了完全独立的时空牢笼。周遭不再是同伴的身影，唯有翻滚不息、变幻莫测的白雾。起初只是空虚与孤寂，但很快，那雾气便仿佛拥有了窥探人心的魔力，开始剧烈地涌动、扭曲，渐渐演化出各种光怪陆离、直击魂魄的景象。

有人眼前猛然浮现出童年时最恐惧的阴影：或许是幽暗巷弄中追逐的恶犬，或许是亲人离世时冰冷的手，或许是独自被遗弃在无尽黑暗中的绝望哭喊。那些早已被岁月尘封、自以为忘却的伤痛，此刻被无比清晰地放大、重现，每一处细节都血淋淋的，痛彻心扉。

有人则陷入求道途中最惨痛的挫折：闭关冲关失败时灵气反噬、经脉寸断的痛苦；比试台上被对手当众羞辱、碾碎骄傲的瞬间；苦心追求的秘籍被他人轻易夺走，自己却无能为力的愤恨与凄凉。这些挫败感如同毒藤，缠绕上来，勒得人喘不过气。

更有人被幻境拖入了权力与情爱的旋涡：眼前是至高无上的宝座，脚下是匍匐的众生，一念可决千万人生死；或是朝思暮想的倾城容颜依偎而来，软语温香，许下生生世世的誓言，诱惑人沉沦其中，永世不醒。

而比这些具体幻象更可怕的是，一股浩瀚、沉重如星辰陨落般的意志，毫无征兆地降临，笼罩了每一寸空间，渗透进每一个念头。

那是沈天穹的法相威压。

这威压被刻意收敛过，不会影响肉体，却渗入神魂最细微的缝隙。它不创造恐惧和欲望，它只是将试炼者内心深处本就存在的那些东西——对未知的恐惧、对力量的贪婪、对抉择的犹豫、对自身弱小的认知——无限地放大、再放大。

这环境对于习惯了灵枢径修行方式的修士而言，尤为残酷。

他们一生的修行，都建立在与外界灵气的沟通、宗门资源的倾斜、同辈实力的比较之上。他们的力量源于外，信心也往往系于外。

当迷雾剥离了所有外在参照——感受不到天地灵气，看不到同为竞争者的同伴，宗门的荣耀与背景化为虚无——只剩下赤裸裸的、孤独的自我，去面对那犹如天道般漠然沉重的威压时，很多人构筑多年的心理防线，瞬间便土崩瓦解了。

白雾笼罩的各个独立空间里，开始传出种种不堪的声响。

有人承受不住内心被放大千百倍的童年阴影，瘫倒在地，像个孩子般痛哭流涕，眼泪鼻涕糊了满脸，全然忘记了修士的体面。

有人被膨胀的权欲吞噬，对着虚无的幻象手舞足蹈，发出癫狂的嘶吼，宣布自己即将君临天下。

有人则在威压催生出的无边恐惧前彻底屈服，对着幻化出的强大存在跪地磕头，涕泪横流地求饶，愿意献出一切换取苟活。

更有人被内心最阴暗的欲望所掌控，面露淫邪丑恶之态，丑态百出。

平日里那些风度翩翩、意气风发的青年才俊，此刻在“问心”的照妖镜下，纷纷显露出了潜藏的脆弱、卑劣与不堪。

\vspace{2em}

实际上，天师府这最后一关试炼，百年来一直备受修士诟病，被认为“过难”甚至“刻薄”：在许多修士看来，修行本就是与天争、与人斗，争夺资源，比拼实力，弱肉强食乃是天道。

如此赤裸裸地拷问人心，挖掘那些潜藏的不堪，是否有违修道本意？是否过于不近人情？

但天师府的态度却百年如一日的坚定，此关从未改变。

或许，在真正的上位者眼中，一个人的天赋、实力固然重要，但若心性不过关，道念不坚定，那么力量越强，未来可能造成的危害反而越大。能在此等环境下保持基本镇定、道心不溃、仍能不断前行的，自是凤毛麟角。

\vspace{2em}

沈瑶亦陷入了苦战。她看到了十年前自己决绝离家时，父亲沈天穹那复杂难言、最终化作拂袖转身的背影，感受到了这十年来漂泊在外、有家难回的深刻孤独。

它还将她脸上那无法消除的灼痕无限放大，让她仿佛能听到旁人背后指指点点的窃窃私语，将她深藏心底的自卑血淋淋地挖出来，暴露在光天化日之下。

更让她道心摇动的是，幻象让她重温了自己曾行侠仗义，却因力量不足或局势复杂，最终事与愿违、甚至带来更坏结果的挫折，引动她对自己所持信念的深切怀疑。

虽凭借着十年历练打磨出的韧性未曾倾覆崩溃，沈瑶却也已是灰头土脸。

然而，在这片遍布崩溃与丑态的幻境中，白言冰与陈千羽所在的角落，却是另一番截然不同的光景。

白言冰眼前幻象纷至沓来：皇子时的锦衣玉食与权谋倾轧，师尊传授“问道途”时的点拨，创立“弥焰斩”时的豪情，青石城外散修们的苦难……

幻境试图勾起他对力量的贪婪、对过往荣华的眷恋，但白言冰心如明镜台。一方面，他修问道途，所养的一点浩然正气护持他身；另一方面，他一路走来，早已明确自身之道“在护我所珍，不在长生虚名”。

外界的威压与幻象，于他而言，如同清风拂山岗，明月照大江，可以感知，却无法撼动其根本。他步履从容，甚至有种审视自身、砥砺心境的意味，胜似闲庭信步。

陈千羽亦然。

她再度面对学医时的艰辛、再度面对病患无力回天的自责、再度面对自己得知身世剧变的恐惧和痛苦……

幻境同样试图诱发她对安稳的渴望，甚至对师尊、对白言冰可能离去的恐惧。但陈千羽心念纯粹，她无论是学医、还是修道，初衷都是救万民于水火，此志早已融入骨髓。

幻象与威压，反而让她更清晰地看到自己学医济世的初心，内心愈发澄澈安宁，行走幻境，亦是从容不迫。

\vspace{2em}

最终，所有的幻象、低语、威压，如同退潮般骤然收束。白言冰、陈千羽、沈瑶以及其他少数在心志煎熬中仍未彻底崩溃的坚毅者，发现自己置身于一片绝对的虚无空间。

就在这绝对的虚无中，一个恢弘、漠然、仿佛不带任何情感，却又至高无上的声音，直接在他们每个人的灵魂最深处响起，如同天道亲自垂询，又如同内心最终的自我审判：

“今可叩长生、登极巅，然须屠戮百万无辜生灵。汝可愿为？”

这终极拷问，简单、直接、粗暴，彻底撕开了所有文明的矫饰、所有权衡的借口、所有自我欺骗的遮羞布！

幻境外，广场高台之上，天师府府主沈天穹与诸位长老，以及虽不参与试炼但被特邀在旁观礼的叶牧谦，都能通过一座早已布置好的特殊阵法，清晰地看到幻境中众人面对这最终拷问时的景象与选择。

只见那虚无空间中，剩下的不到十几人，面对这直击灵魂的问题，反应瞬间千差万别，人性百态在此刻暴露无遗。

有人脸上在愣神刹那后，瞬间被狂喜和贪婪淹没，眼中爆发出骇人的精光，几乎迫不及待地就要点头应“可”，仿佛那百万生灵不过是达成其野心的区区数字。

有人则陷入了前所未有的巨大挣扎，脸色在渴望、恐惧、愧疚、狠厉之间飞速变幻，额头上冷汗涔涔，嘴唇哆嗦，显见内心正经历着惊涛骇浪般的冲突。

但令高台上不少长老暗自摇头的是，剩下的大多数人，在或长或短的沉默或挣扎后，竟缓缓地点了点头，或是从喉咙深处，艰难而低沉地挤出了一个“可”字！

为了那虚无缥缈的长生与极巅，为了那可能存在的、超越一切的力量，百万活生生的、与他们无冤无仇的无辜生灵，在他们被剥离伪装的本心衡量下，竟成了可以交换、可以牺牲的冰冷筹码！

这便是许多修士，当一切外在装饰、社会约束、道德言说都被幻境剥离后，内心深处对个体力量极致渴望的真实写照，乃至为达目的可以不择手段的潜在本能。

他们的道，终究是外求的、利己的、以自我为中心的力量之道。

沈瑶浑身剧烈一震，脸色瞬间变得煞白，不见一丝血色。

那个问题响起的瞬间，她灵魂深处迸发出的第一反应，是清晰而强烈的“不可”！那是她十年江湖历练、心中仍未泯灭的良善与侠义本能在呐喊。

但就在这个念头升起的刹那，无数杂音也随之在她脑海中轰然炸响。

让她动摇的，是她十年来亲眼所见的无数现实无奈：很多时候，即便有心行善，也会因力量不足、局势复杂、牵扯太多而徒劳无功，甚至适得其反。如果真的能得到如许力量，那很多问题都不会再是任何问题……这样一来，那“百万生灵”就成了必要的牺牲。

可是这牺牲……真是必要的吗？

她陷入了无尽的自我怀疑。

最终，在幻境那无形的凝视、无声的催促之下，她只是极其痛苦地摇了摇头，嘴唇翕动了数次，仿佛用尽了全身力气，却终究未能将那一声坚定的“不可”说出口。

然而，就在这一片或狂热、或默许、或犹豫的死寂气氛中，两个清晰、坚定、没有丝毫迟疑与颤抖的声音，如同划破厚重阴霾的惊雷，几乎同时响起。

白言冰刚听完这番问话，便不假思索地回应：

“某之道，在护某所珍，不在长生。若需屠戮无辜，这长生不要也罢！”

声音斩钉截铁，每一个字都掷地有声。

言毕，掷剑于地。

剑身轻颤，嗡鸣不止。

另一侧，陈千羽的神情宁静如深潭秋水，却同样蕴含着不容置疑的决绝。她缓缓开口，声音清越而平稳：

“千羽自幼学医，愿以一身之力，救万民于水火，解生灵于倒悬。自入师门，一矢以终。此等长生，千羽不愿。”

两人的回答，没有片刻的权衡利弊，没有一丝一毫的勉强与作伪。

高台之上，一直密切关注着幻境内一切的沈天穹，盘坐的身躯猛然一震！

饶是以他法相境的修为、三百多年的阅历，此刻心中也掀起了惊涛骇浪。他见过太多惊才绝艳的天才，更听过无数次慷慨激昂、正气凛然的誓言；但这些誓言，或是出于虚伪矫饰，或是根本无法在日后残酷现实中践行。

但像白言冰与陈千羽这般，在如此直击本心、毫无转圜余地的终极拷问下，能如此不假思索、如此纯粹坚定地给出否定答案，甚至不惜以掷剑明志这般激烈行为表露心迹者，简直是凤毛麟角，百年罕见。

这已不仅仅是心性坚毅所能解释，这根本是他们所秉承的“道”，其内核就与那种自私冷酷、视众生为蝼蚁刍狗的长生之路截然相反，水火不容！

而且沈天穹不知为何，也极其坚信：他们的承诺，确实可以践行一生！

好一个叶牧谦，竟能教出这样的徒弟！

他目光扫过幻境中女儿沈瑶的身影时，心中涌起一股复杂的欣慰。女儿虽然未能如白陈二人那般斩钉截铁、毫无滞碍，但她的第一反应是“不可”。

这完全可以充分说明她心底的良善未泯。十年的红尘历练，风霜雨雪，并未将她磨砺成一个冷漠无情、只知利害计算之徒，这本身已是难能可贵的成长。

\vspace{2em}

幻境消散，所有试炼者被一股柔和而不可抗拒的力量送回山门前的广场。

阳光刺眼，微风拂面，真实的喧哗声隐约传来，却让大多数重返现实的试炼者恍如隔世。

经历最后一关的拷问，最终能保持内心良善底线、未向那诱惑低头的，仅有七人。除了白言冰、陈千羽、沈瑶，还有另外四名心志确实堪称坚如铁石、在挣扎后守住了底线的修士。

但这四人此刻大多神情萎靡到了极点，眼神涣散，气息虚浮紊乱，有些人甚至站立不稳，需要身旁同门搀扶，显然在幻境中心神消耗巨大。

白发长老目光转向高台，端坐的沈天穹面无表情，只微微颔首示意。

长老上前几步，目光扫过下方神色各异的众人，朗声宣布，声音清晰地传遍广场每一个角落：“经三关试炼综合评定，由府主暨诸位长老共同裁定，以下二人表现卓异，心性澄明可嘉，并为本次天师试炼魁首：白言冰、陈千羽！授予‘小天师’称号，此后待遇与天师府核心弟子等同，可随时入密阁修行！”

他顿了顿，继续道：“此外，沈瑶等五人，亦表现突出，心志坚韧，实力不俗，准予十年内进入密阁修行！”

宣布完毕，广场上响起一阵压低了的、复杂的议论声。有羡慕，有嫉妒，也有对那惨烈淘汰过程的余悸。

\vspace{2em}

沈瑶听着父亲通过长老之口做出的宣布，心中却并无多少预期中的喜悦或骄傲，反而被一种更为复杂难言的情绪所充斥。

她站在原地，目光不由自主地看向身旁不远处的白言冰与陈千羽。两人气息平稳，目光清澈依旧，仿佛刚才那场足以让许多人心神崩溃的幻境拷问，对他们而言只是一次寻常的散步；再回想自己在幻境中面对终极问题时的摇摆、痛苦与最终的沉默，一种前所未有的强烈震动与自惭形秽之感，如同冰冷的潮水，悄然弥漫上心头，将她淹没。

她毕生奉行、并引以为傲的“争”，与不公斗，与命运争，在艰难世事中夺取一线生机与公道，她一直认为这便是践行侠义，便是自己的道。

可当那最根本的善恶抉择赤裸裸、毫无花哨地摆在面前时，她的“争”，她那需要依赖外在情境、需要权衡利弊的“争”，却无法给出一个像白言冰与陈千羽那般，斩钉截铁、源于本心、不容置疑的答案。

这个残酷的发现，让她心中某个自认为坚固无比的部分，产生了清晰的裂痕。

这认知带来的不仅是自惭，更有一种混合着震撼与强烈渴求的情绪，在她胸中汹涌澎湃，几乎要破膛而出。她几乎是无意识的，脚步不由自主地移动，默默地跟在了正准备离开广场的叶牧谦师徒三人身后。

夕阳西下，将天边云霞染成一片绚烂的金红，也将几人的影子在蜿蜒下山的青石小径上拉得很长。他们来到了天师府山门外，一处专供客卿与重要访客暂居的幽静院落前。院落清雅，竹林掩映，与山门内的恢宏庄严截然不同。

就在叶牧谦面带淡然笑意，即将举步踏入院门的刹那，一直默默跟在后面的沈瑶，忽然停下了脚步。

“先生留步！”

她清冽的嗓音响起，带着一丝难以完全掩饰的颤抖，在寂静的傍晚显得格外清晰。

\vspace{2em}

叶牧谦转身，白言冰与陈千羽也讶异地回头看来。

在三人目光的注视下，沈瑶上前几步，来到叶牧谦面前。下一刻，她竟毫不犹豫地双膝一屈，“噗通”一声，实实在在地跪倒在坚硬的青石地面上，对着叶牧谦，深深地叩首下去。

“沈姑娘，你这是……”陈千羽忍不住惊呼出声。

沈瑶抬起头，傍晚的余晖映照着她面纱上方露出的那双眸子，此刻那眼中亮得惊人，仿佛燃着两簇不肯熄灭的火焰。她直视着微微愕然的叶牧谦，一字一句，清晰而决绝地说道：

“先生，沈瑶……恳请拜入先生门下，修习‘问道途’！”

叶牧谦微微一怔。

收徒？收沈瑶？

白言冰和陈千羽都是机缘巧合下，几乎从零开始跟着他瞎编的“问道途”的。教起来虽有压力，但好歹是自己从头塑造。更重要的是，他们并不会怀疑他；他说什么就是什么，而且因为些他暂时也不知道为什么的原因，无论如何都能弄出些实际的东西来。

可沈瑶不同，她已是玄脉境修为，修炼的是天师府顶级的灵枢径功法，根基已成，见识广博。教她？万一被她看出这“问道途”诸多境界、理论尚是草创，甚至有些地方是为了圆谎而临时拼凑，岂不立刻露馅？

于私，以后他这“高人”还怎么装下去？

更重要的是，于公，如果真给教坏了……且不论以沈天穹为首的别人怎么看他，他自己良心都过不去。

于是他做出了一个他认为既合理又体面的决定。

“沈姑娘请起。”叶牧谦伸手虚扶，语气温和却带着疏离，“你乃天师府千金，家学渊源深厚，沈府主更是法相境大能，何须拜我这山野散人为师？‘问道途’粗陋，恐误了姑娘前程。”

他以为这番婉拒合情合理；殊不知，沈瑶是个一旦认准便狠绝到底的人。

试炼幻境中的冲击，让她对自身旧有之道产生了根本性质疑，而对叶牧谦师徒展现出的那种内心力量产生了近乎本能的向往。叶牧谦的拒绝，在她听来，却成了另一种暗示——莫非先生是嫌我已有灵枢径修为在身，道基不纯，与问道途内求己身、元气自生的理念格格不入？

这个念头一起，便再也无法遏制。

只见沈瑶惨然一笑，声音却异常平静：“先生许是看我这身修为不顺？觉得瑶儿心念不纯，旧染未除？”

她顿了顿，眼中闪过一丝决绝的疯狂，“那……这身修为，不要也罢！”

“沈姑娘不可！”白言冰和陈千羽察觉不对，齐声喝道，想要上前阻止。

却已然晚了！

沈瑶跪姿不变，双手却猛地掐出一个逆转灵力、自毁道基的凶险诀印！她体内原本顺畅运行的、属于玄脉境开脉期的精纯灵力，一瞬之内，轰然倒卷、逆行！

“噗——！”

细密而令人牙酸的崩裂声仿佛从她体内深处传来，那是经脉在狂暴逆行的灵力冲击下寸寸受损、乃至断裂的声音！

沈瑶娇躯剧震，脸色瞬间血色尽褪，变得如同金纸。随即，一口殷红的鲜血无法抑制地狂喷而出，紧接着口鼻之中皆有血丝不断溢出。她周身原本隐隐环绕的灵力波动，如同溃堤般迅速消散、湮灭，气息以肉眼可见的速度萎靡下去。

剧痛如同千万把钢刀在体内搅动，眼前阵阵发黑。沈瑶全靠一股顽强的意志死死撑着，才没有当场昏死过去，但身体已经摇摇欲坠，鲜血染红了胸前的衣襟和面前的地面，模样凄惨无比。

这突如其来的变故，让所有人都惊呆了。

白言冰和陈千羽僵在原地，难以置信。他们万万没想到，沈瑶为了拜师，竟能对自己狠绝至此！自废修为，逆转经脉，这其中的痛苦与凶险，无异于自杀！

叶牧谦更是惊愕地瞪大了眼睛，看着面前瞬间变成血人、气息奄奄却仍倔强地抬头望着他的沈瑶。他心中那点关于“露馅”的担忧，在这般惨烈而决绝的拜师决心面前，显得如此渺小可笑。

这……对自己也太狠了！

震惊过后，一股复杂的情绪涌上心头，有无奈，有动容，也有一丝面对如此纯粹执着的无力感。他终究不是铁石心肠之人。

这番图景，在收白言冰的时候，他已经见过一遍了。

叶牧谦长长地、无奈地叹了口气，脸上露出一抹苦笑，摇了摇头：“罢了罢了……你这又是何苦。”

他上前一步，看着沈瑶那双即便因剧痛而涣散、却依然执拗地盯着他的眼睛，缓缓说道：“此后……你便为我门下第三徒。”

“第三徒”三个字传入耳中，沈瑶一直紧绷着的那根弦，终于松了。眼中闪过一丝如释重负的微弱光芒，提着的那口气骤然泄去，她头一歪，彻底失去了意识，向前软倒。

“快！”叶牧谦低喝。

白言冰和陈千羽这才慌忙上前。陈千羽迅速探了探她的脉搏和气息，脸色凝重：“经脉损毁严重，元气大伤，必须立刻救治！”

两人不敢耽搁，连忙将沈瑶抬进院落中最好的厢房，轻轻平放在床榻上；叶牧谦站在房门口，看着屋内忙碌的两人和床上气息微弱的沈瑶，眉头紧锁。

事情的发展，完全超出了他的预料。

\vspace{2em}

果然，没过半柱香的时间，院落外便传来一阵急促而沉重的脚步声，伴随着一道又惊又怒、蕴含着磅礴威压的喝问：“瑶儿！我的瑶儿呢？！她怎么样了？！”

人影一闪，天师府府主，法相境大能沈天穹，竟已满脸焦急地冲进了院子，完全没有了平日里的威严气度，俨然一个担忧爱女安危的普通父亲。他显然是通过某种方式，瞬间感知到了沈瑶气息的剧变和生命危险。

“沈府主。”叶牧谦迎上前。

沈天穹一眼就看到了厢房内床上昏迷不醒、浑身插满银针的女儿，以及正在全力施救的陈千羽，脸色顿时变得更加难看。他强压怒火，听叶牧谦简略说明了事情经过——从沈瑶跪地拜师，到她决绝自废修为，再到叶牧谦无奈收徒。

听完，沈天穹沉默了。

他没有立刻发作，只是深深地、久久地凝视着床上脸色苍白如纸的女儿，那双阅尽沧桑的眼眸中，翻涌着心疼、不解、无奈，最终化为一声长长的、沉重的叹息。他看向叶牧谦，眼神复杂，缓缓说道：“罢了……这是瑶儿自己的选择。”

他了解自己的女儿，知道她一旦认定，便是九头牛也拉不回。既然她以如此惨烈的代价做出了选择，他这个父亲，除了接受，又能如何？

是夜，月华如水。沈瑶在陈千羽不计代价的救治下，情况暂时稳定，但仍在深度昏迷中，需要时间慢慢恢复受损的根基。

院落偏厅，沈天穹主动提及，想跟叶牧谦谈谈。

挥退了侍从，布下隔音结界，沈天穹仿佛卸下了府主的威严，显出一丝疲惫。他亲手为叶牧谦斟了一杯茶：

“叶先生，瑶儿性子执拗，今日之事，让先生见笑了。”沈天穹开门见山，语气诚恳，“既然瑶儿已拜入先生门下，有些关于她、关于天师府、乃至关于这内域的事情，我觉得有必要告知先生。”

叶牧谦正色道：“府主请讲，叶某洗耳恭听。”

沈天穹目光悠远，仿佛陷入了回忆：“瑶儿是我三百岁的时候出生的，今年刚二十五岁。她……自幼便心地善良，见不得不平之事。”

“这是好事。”叶牧谦点头。

沈天穹却补充：“但我天师府，虽位列内域五大超级势力之一，数百年来，却一直奉行孤立主义之策；我上任之时也一以贯之。”

他顿了顿，继续道：“我天师府根基，主要在于天枢城，以及另一座大城天权城。对于这两城之外，内域广袤土地上发生的诸多争斗、纷扰、乃至不平事……天师府历来是能不管便不管，能不理便不理。”

叶牧谦微微挑眉，这与他之前对“超级势力”的想象有所不同。

沈天穹看出了他的疑惑，苦笑道：“瑶儿便是因此，与我、与天师府产生了根本分歧。她认为，天师府既有如此力量，便应当承担起相应的责任，维护内域大体安宁，扶助弱小，铲除邪佞，方不负‘超级势力’之名。她认为我们……没有尽到义务。”

“那府主的考量是？”叶牧谦问道。

沈天穹神色转为凝重：“内域局势，远非表面看起来那么简单。五大超级势力：天师府算最强，其余四个分别是紫阳宗、无极剑派、万仙盟、丹鼎阁。而剩余四个中，又属紫阳宗最强，其宗主也是法相境外显期，与我颇有私交，只是近年闭关不出。

“超级势力看似超然，实则关系微妙，彼此制衡。许多发生在中小势力、乃至散修之间的‘小打小闹’，背后未必没有其他大势力甚至超级势力的影子，或默许，或推波助澜，以此试探、消耗对手，或达成某些不便明言的目的。”

他看向叶牧谦，“贸然插手，很可能便是不小心踩进了别人精心布置的棋局，甚至可能引发超级势力之间的直接冲突，后果不堪设想。天师府传承不易，我作为府主，首要职责是维系府内安定与传承，不可轻易将整个宗门拖入不可测的漩涡。”

这便是现实与理想，全局安危与局部正义之间的残酷矛盾。沈天穹的选择是保守、“无为”，而沈瑶则渴望激进的“有为”。

“我们父女为此争吵过无数次。”沈天穹叹息，“十年前，拙荆因为寿元耗尽去世了，家里气氛本就压抑。结果，一次激烈的争执后，她一气之下，便离家出走了。这一走，就没回来过。”

他本以为，女儿在外历练，见识到世间的险恶与复杂，会逐渐理解他的苦心与无奈。可事实恰恰相反，沈瑶在十年独自闯荡、行侠仗义、除魔卫道的经历中，亲眼目睹了更多因强者漠视、势力倾轧而造成的苦难，她的信念非但没有动摇，反而愈发坚定了！她认为父亲的“孤立”是怯懦，是冷漠，是导致诸多悲剧的根源之一。

父女间的理念鸿沟，非但没有因时间弥合，反而越来越深。

说完这些，沈天穹仿佛卸下了一副重担，神情却更加萧索。他看向叶牧谦，郑重道：“瑶儿如今拜你为师，或许……她的道，真的与先生有缘。她选择破而后立，我虽心疼，但也尊重。今后，还望先生多加教导。”

叶牧谦肃然点头：“既入我门，自当尽心。”

沈天穹最后道：“还有一事。半月之后，我天师府将做东，邀请内域其余几大超级势力的代表，于天枢城举行‘论道会’，坐而论道，交流心得。届时，还请先生务必拨冗参加。”

叶牧谦心知无法推脱，只得应承下来：“府主相邀，叶某自当赴会。”

沈天穹又去看了看依旧昏迷的沈瑶，嘱咐陈千羽若有任何需要，天师府资源尽可调用，这才带着满腹心事离去。

\vspace{2em}

接下来的日子，小院成了临时的医馆和修炼地。陈千羽几乎不眠不休，全力为沈瑶调理。

令人惊喜的是，在治疗过程中，陈千羽发现沈瑶半边脸颊那陈年的灼伤，在精纯元气的滋养与银针的引导下，坏死的肌肤组织竟然焕发生机，疤痕逐渐淡化、平复。她细心调理，数日之后，沈瑶脸上的灼痕竟已消失不见，恢复了原本清丽绝伦的容颜。

沈瑶在昏迷三日后终于苏醒。身体依旧虚弱，但经脉的损伤已在缓慢修复。

当她再一次在陈千羽递过的铜镜中，看到自己完整无瑕的脸庞时，整个人都怔住了，手指颤抖着抚过光滑的肌肤，久久无言。

如今伤痕褪去，仿佛也卸下了一层沉重的枷锁。

她沉默地将那方陪伴多年的黑色面纱仔细叠好，收进了储物袋深处，不再佩戴。

\vspace{2em}

经此一劫，沈瑶的心境发生了些微妙的变化。就像被湍流冲刷过的卵石一样，少了几分尖锐，多了些圆润和沉静。

果然是有修炼底子的人，悟性、经历也是一等一的。身体彻底恢复好之后，不过三日，沈瑶便将幻术和身法相关的理解融于一身，稳步踏入淬体境天骨期，修为回归养气境也只是时间问题。

这速度竟比当初的白言冰还要快上几分。

由于叶牧谦受邀论道，几人目前依然暂居于天师府提供的馆舍中。她自然便与叶牧谦等人住在了同一处。

夜里，枕着柔软舒适的锦枕，身下是干燥洁净的被褥。隔壁房间，偶尔会传来白言冰与陈千羽压低嗓音的交谈；更远一些，叶牧谦房中那规律而悠长的鼾声，穿过墙壁隐约可闻。

没有追杀，没有戒备，无需时刻紧绷心弦。她将半边脸埋入枕中，感受着久违的、近乎奢侈的安宁，意识沉沉陷入黑甜的梦乡。

\vspace{2em}

另一边，白言冰在照料沈瑶伤势之余，试炼幻境最后那直指道心的拷问也如同烙印，深深烫在他的神魂深处。某一日午后，他于院中观落叶飘零，忽觉灵台一点清明跃动，似有所感。

他即刻寻到叶牧谦，恭敬请示：“师尊，弟子心有所惑，亦有所感，想觅一静室闭关几日。”

叶牧谦从图谱上抬起头，目光在他沉静却隐含锋锐的脸上一扫，了然颔首。

白言冰寻了馆舍深处一间空置的静室，石门落下，隔绝内外。室内无窗，仅一颗明珠嵌于顶壁，散发着柔和的冷光。

他盘膝坐下，将心神彻底沉入那片迷雾之中，反复咀嚼那拷问，审视自己执剑的初衷。

时间在绝对的寂静中流逝，不知过了几个时辰，抑或是一日。他周身原本内敛的气息开始产生微妙的变化，时而如剑锋般凌厉欲出，时而又如古井般沉寂无波，最终，所有的波动渐渐归于一种奇异的圆融。

他第二次认识了“我”。

下一步，便是将这已然清晰的“我”之意志，通过元气，映照于外，化为实质的领域。这对已然厘清我与物分别的白言冰而言，是非常容易理解的，却需要精细的操控。

心念微动，气海中温养的精纯元气随之响应。他先尝试着引导一丝元气，沿着手臂的轨迹，如同沿着剑锋延伸般缓缓向外探出。

起初，这缕元气一旦离体稍远便迅速涣散，难以维持。这和先前元气附剑时是一样的，放出去的元气就不再能收回了。

白言冰不急不躁，反复尝试，调整着输出的强度与心神依附的浓度。

渐渐地，那外探的元气变得稳定，已能在他掌心之上三寸处凝而不散；白言冰试着轻轻屈曲手指，那些外放的元气则听话地被召回体内。

第一步已经成功了，第二步便是试着让这些离体的元气在周身之外自行流转，形成一个大致的循环。

不知又过了多久，量变引发质变。随着白言冰精微而持续的引导，气海中精纯的元气终于不再需要他去刻意驱策，似乎与外界那层循环初步建立的元气产生了共鸣，开始自发地、坚定地向外漫溢、补充，参与那周而复始的运转。

一个稳固、通透、心意朗照的独特势场，终于彻底成型，将他周身三尺之地笼罩在内。

在这势场之中，外界的纷扰被自然过滤、辨析；而他自身的意志与存在感，被无比清晰地映照和确立——这便是属于他白言冰的见独领域。心念一转，领域内气息可骤然变得锋锐，带着凛冽的剑意；心念一松，又可复归温润平静。如臂使指，不外如是。

白言冰缓缓睁开眼，眸中似有清光流转，旋即内敛。

他踏出静室，晨光落在他肩上。虽然这初成的、更偏向于心神感知与意志凝聚的剑意领域，远不如叶牧谦那浩瀚如星空、可直击神魂的心念领域那般深邃强势，甚至显得有些简陋单薄；但每一个运转的环节，每一缕外放元气的轨迹，都打上了他自身剑道意志的烙印。

更重要的是，越简单的东西，越不容易坏。他的剑意领域，似乎专门就是为了守护而生的。

这终归是他自己走出来的路。

\vspace{2em}

就在他成功突破、势场稳固成型的刹那，在自己房间内休息的叶牧谦，身躯微微一震：他感到一种熟悉的清流，无声无息地自冥冥中汇入他的道基之中。

“又是反哺……”叶牧谦嘴角勾起一丝细微的弧度，眼中闪过欣慰。

白言冰把“见独”走实了之后，叶牧谦对见独境的细微掌控、种种精妙运用的可能性，似乎瞬间挣脱了一层此前未曾察觉的无形薄膜，变得前所未有的圆转自如，得心应手，也终于不需要再担心反噬问题了。

\vspace{2em}

至于天师府密阁，叶牧谦并未急于让弟子们前往。待沈瑶伤势稳定、白言冰出关后，他将三人唤至面前。窗外竹影婆娑，映得他青色长衫仿佛也染上了几分翠意。

“天师府密阁，藏书浩瀚，包罗万象，是内域无数修士梦寐以求的宝地。你们入内观阅，自是机缘。”

他话锋一转，神色转为严肃：“然，需谨记于心：道途根本，在于内求诸己，明心见性。一切外来的知识、法门，皆需根据自身实际情况，作出实事求是的修改与化用，使之真正契合你们自身。若只知一味模仿他人形貌，邯郸学步，而失却自身那一点真灵神魂，便是落了下乘，前路必窄。”

学之者生，似之者死。

这是他对弟子治学修道、钻研术法的根本要求。

他又特意补充道：“阁中若有关于远古炼器手法、以及阵法一道的古老记载、残篇、甚至只言片语，若有闲暇，可帮我留意一番，取出容我一观。”

“弟子领命。”沈瑶躬身应下。

她心中感慨，十年前，她凭借特殊身份，偷偷进入过密阁外围；此番，却是堂堂正正而入。

半月后，待她身体再好些，便领着白言冰与陈千羽再次来到那幢气势恢宏又戒备森严的阁楼前。

验过令牌，沉重的玄铁大门缓缓开启。沈瑶步履轻盈，穿过陌生又熟悉的廊道与禁制光芒，很快便将二人带至存放炼器、阵法及相关杂学典籍的偏殿区域，效率奇高。

\vspace{2em}

半月时光，便在沈瑶日渐红润的面色、白言冰日益稳固圆融的见独境气息、陈千羽对几卷古老医典的痴迷研读、以及叶牧谦对零星收集来的残破炼器阵法图谱的初步推演中，悄然滑过。

\vspace{2em}

很快，便是府主沈天穹前日所言的论道会之期。

五方主位之上，端坐着五位威仪各自笼罩一方的身影。

东首主位，自然是地主，天师府府主沈天穹。

其左侧下首，是一位徐娘半老、身着丹霞色宫装长裙的女子，云鬓高挽，插着一支药玉簪。此人正是丹鼎阁阁主，文明内域的“丹王”王丹。其修为已至真符境注灵期，举手投足间，雍容华贵，又带着丹师特有的沉静气度。

王丹的下首，是一位赤发如火、根根竖立如钢针的男子，面容刚毅如同斧凿刀刻，浓眉之下双目炯炯有神。此人乃是紫阳宗副宗主，陈炎，修为在真符境显化期，性情素来以火爆直接、不屈不挠著称，与他的火系功法相得益彰。

沈天穹右侧，是一位青衫老者。衣衫洗得发白，纤尘不染，背后负着一柄样式古朴的长剑，剑柄缠着磨损的暗色布条。他面容清癯，颧骨微高，须发灰白，眼神初看温和润泽，细细品味，却能察觉到凛冽的剑锋。正是无极剑派心剑一脉的长老，独孤明，修为亦是真符境注灵期。

最末一位主位上的，是位看起来颇为儒雅随和的中年人，身着素雅的文士衫，手中轻轻摇着一柄白玉为骨的折扇，脸上常挂着三分令人如沐春风的微笑。他便是内域最大散修联盟“万仙盟”的盟主，吴东浩，修为在真符境摹形期。虽然说这位吴东浩修为在诸人中最低，但以散修之身拼搏至此，其手腕与心机也是无人敢小觑的。

阁宇中央，相对五方主位设一独立客席。叶牧谦安然独坐其上，依旧是一身半新不旧的青衫，神色平静无波，迎着来自五个方向、或明或暗的审视目光，坦然自若。白言冰、陈千羽、沈瑶三人，则与其余少数同被邀请但无资格上台的修士们坐在一起。

此前，沈天穹早已将叶牧谦及其所持“新的道途”之信息，简略告知其余四方。此刻，诸方巨擘目光汇聚于那青衫客卿身上。

好奇有之，探究有之，但更多的，是一种根植于漫长修行岁月、对于未知与可能颠覆现有格局的“异端”那种本能的不信任与审视。

\vspace{2em}

作为东道主，沈天穹率先开口，声音平和醇厚，打破了阁中的沉寂：“叶先生，今日诸位道友齐聚我这问道阁，皆因闻听先生之道，与我等所行所知的‘灵枢径’迥然不同，心生探究印证之意。先生不妨简述其道之要旨，也好解诸位道友心中之惑。”

叶牧谦微微颔首，面向众人，朗声开口，声音清越，在空旷的阁中回荡：“沈府主，诸位道友。在下所行之道，自名为‘问道途’。其核心理念，在于‘内求诸己，明心见性’。不假外求灵枢以为沟通天地之桥梁，不刻意开拓体内玄脉以强纳外界灵气，而是于丹田之下，开辟气海，温养调和自身生命本源所生之元气，以此筑基。”

话音刚落，左侧便传来一声洪钟般的诘问，带着毫不掩饰的质疑与火气。

正是陈炎，赤眉扬起，声震屋瓦：“叶先生此言，未免过于空泛虚渺！不引灵枢，不通玄脉，如何引动天地浩瀚灵气入体修行？既无灵力傍身，如何施展神通妙法？莫非你这‘问道途’，只能让人坐而论道、空谈心性不成？便是修士最基本的飞檐走壁、御器凌空、符箓咒法，如何做到？”

他的问题尖锐直接，毫不留情，恰恰戳中了目前问道途在外人看来最明显的短板——缺乏灵枢径修士那些直观、炫目且极其实用的外在能力。

叶牧谦对此似乎早有预料，面上并无愠色，也不着急，不疾不徐地答道：“陈副宗主所问，正点出了两途根本立意之差异所在。问道途讲求内求于己，元气源于自身气血、精神、性命之调和温养，储存于气海，运行于自身小天地，其根基便在于‘不假外物’。至于陈副宗主所言飞檐走壁、御器凌空……”

他略微顿了顿，脸上浮现出一丝淡然甚至略带超脱的笑意，目光平静地迎向陈炎灼灼的视线，反问道：“修行之人，为何……就非要飞呢？”

他语气平和，却带着一种奇特的、近乎哲辩的意味：“脚踏实地，一步一印，感受山川厚重；观花开花落，体悟四时轮转生灭；听风过松涛，明辨八荒气息流变。这行走坐卧、俯仰天地之间，所感所悟，不亦是‘道’之所在吗？”

这话语带着几分有意为之的诡辩与超然物外的姿态，试图将对方从“力量表现”的务实层面，拉入自己“重内修心性意境”的语境之中。

陈炎果然被这“为何非要飞”的反问结结实实地噎了一下。他修炼的乃是霸道炽烈的火系功法，性子亦是刚猛直接，向来认为强大的力量就该有与之相匹配的威势与外在表现，腾云驾雾、御剑飞天乃是修士的常能与标志。叶牧谦这种近乎逃避问题本质的说法，在他听来，简直强词夺理。

他张了张嘴，赤红的脸膛因激动更显颜色，想要驳斥，可一时间又觉得，跟一个似乎根本不在乎“能不能飞”“会不会飞”的人，去争论“怎么飞得更好更快”，简直如同对牛弹琴，徒费口舌。满腔话语堵在喉咙口，最终只化作一声重重的不满的冷哼，撇过头去。

然而，表面上镇定自若的叶牧谦，内心却是另一番景象。他暗自警醒：下一个境界，无论如何也必须设法衍生或融合出类似“腾空”、“远程攻伐”等基础的实用能力了。

否则，光靠嘴皮子讲道理，短期或许能震慑一时；长远来看，终究难以服众，更会严重限制弟子们未来的发展空间与安全。

\vspace{2em}

这时，一个温和却异常清晰，仿佛能直接传入人心底的声音响起，带着剑器般的清越质感。正是无极剑派的独孤明长老。他微侧身形，看向叶牧谦，缓缓道：“叶先生，老夫独孤明。我辈剑修，虽不喜主动挑起无谓之争，却也从不畏惧卷入争斗，手中之剑，只为守护心中所珍视之人、所坚守之道而鸣。先生之道，重内修心性，明己见独，老夫细听之下，甚为欣赏其中返璞归真之意。”

他话锋随即一转，语气依旧平和，问题却直指核心：“然，修行之路，漫漫悠长，终究无法完全避开与‘力’的较量。守护需力量，卫道亦需依仗。不知叶先生这‘问道途’，若论及实际护道卫己之能，与我等所熟稔之道相较，究竟如何？可有独到之处？抑或仅是镜花水月，不堪实用？”

这个问题，比陈炎的直接诘问更加实际，也更为深刻。它绕开了外在形式的争论，直指修行道途无论理念如何超脱，最终都无法完全规避的核心——护身与制敌的力量。

叶牧谦心知，面对独孤明这等浸淫剑道数百年、心志早已千锤百炼如精钢的老者，空谈理念、故弄玄虚已然毫无用处。道之高低，终究需落于实处印证。

他从容起身，对独孤明所在方位，郑重地拱手一礼，声音斩钉截铁：“独孤长老所问，切中肯綮！道之高低优劣，纵有千般言语，终需实际印证，方见真章。叶某不才，修为浅薄，然愿以当前之微末能为，请长老试手。”

此言一出，阁内气氛骤然一凝。沈天穹眼帘微垂，不动声色；王丹眸光闪动，似有好奇；吴东浩手中折扇轻摇的频率微微一顿；陈炎则猛地转回头，赤瞳中精光爆射，紧紧盯住场中。

独孤明眼中亦是精光一闪，仿佛沉眠的古剑骤然苏醒了一瞬。他反而因叶牧谦的坦荡与果决微微颔首，面上掠过一丝激赏：“好！叶先生爽快！既如此，老夫便以三招，领教先生之妙！”

言罢，他抬起右手，食指与中指并拢，神色平淡地向着数丈之外的叶牧谦，轻轻向前一点！

“嗤——”

一道凝练到极致，锋锐无匹、直指人心的凛冽剑气，破空而至！

这是如此简单的一招。先不论场上剩下那四个人怎么看了，就连场下的白言冰也不觉得这招有多高明。能看出来独孤明是狠狠留手了。

叶牧谦自然无声展开自己的见独领域。前些日子白言冰成功突破的反哺，使得他目前已达见独境圆满，单论实力已经不虚任何一位真符境修士。

那道凌厉无匹的剑气射入这片心念领域的刹那，其摧枯拉朽的锋锐锋芒，仿佛被无形而柔韧的水温柔地包裹、浸润、稀释了。叶牧谦站在原地，青衫下摆都未曾拂动一下，身形稳如山岳。

“咦？”独孤明口中发出一声几不可闻的轻咦，略显讶异。他承认这第一指确实是收着力的试探，却也绝非寻常玄脉境修士能够轻易接下的，没想到被叶牧谦这般不着痕迹地化解了。

独孤明寻思，得来点真东西了。紧接着，第二指点出！

这一指，剑意骤然变得繁复莫测！一刹那间，仿佛有千道万道虚实难辨的剑影藏于这一指之中迸发，似真似幻，交织成网，更带着迷惑、缠绕、窥探之意。

正是独孤明发扬光大的心剑脉杀招之一，双生剑法！

叶牧谦依旧以见独领域从容应对。领域之内，那漫天袭来的虚幻剑影，在他那澄明如镜的心神映照之下，纷纷显露出其虚幻的本质，如阳光下的泡沫般逐一幻灭。

而藏于这万千虚幻之中，那唯一一道真实不虚、凝聚着独孤明此番剑意核心的剑意，则被他自身圆融无碍、坚定如一的心意悄然一引，擦着他的身侧偏转开去，未伤及分毫。

两招已过，叶牧谦依旧稳立当场。

独孤明脸上轻松的神色终于彻底收起，取而代之的是一片郑重。他缓缓点了点头，深深吸了一口气。

第三指，缓缓点出。这一指，没有凌厉疾刺，也没有万千幻影，只有一道凝实到极点、沉重到极致、仿佛挟带着独孤明毕生对“守护”二字全部感悟的剑意，缓缓推出。

这剑意虽朴素至极，但厚重如山，坚定如岳。它以一种无可阻挡、不可撼动、亘古长存的气势压迫而来：

守护的背后，是牺牲，是承担，是即便天地倾覆亦不退半步的决绝！

叶牧谦瞳孔微缩，深知已到关键。他能轻易看出来这一剑和前两剑本质上的差别，自然不敢怠慢，将见独领域催发到极致！

刹那间，那无形的领域仿佛被注入了灵魂，有了鲜明的色彩与温度。

两股同样坚定、恢弘，但内涵却截然不同的“守护”意志——独孤明那历经岁月沉淀、以剑为誓、厚重如山的守护，与叶牧谦那源于本心、无愧天地、指引前路的守护——轰然对撞！

没有震耳欲聋的巨响，没有绚丽夺目的灵光爆炸。然而，阁中所有人，包括主位之上修为最高的沈天穹，都在这一刹那，感到自己的心神仿佛被一柄无形的大锤轻轻敲击了一下，微微一震！意识海中泛起涟漪，道心为之所感。

三息。

令人屏息的沉寂持续了三息时间。

独孤明缓缓地、极其郑重地收回了伸出的手指。他脸上最初的惊讶之色早已褪去，逐渐化为一种由衷的、毫不掩饰的赞叹与钦佩。他站起身来，不再是踞坐的姿态，而是对着数丈外同样肃然而立的叶牧谦，双手抱拳，郑重地躬身一礼。

“叶先生之道，心性之坚纯，领域之玄妙，老夫佩服！”

他的声音沉凝有力，回荡在寂静的阁中，“此‘问道途’，确有惊天独到之处！于实战之中，尤其是心神交锋、意志对抗之际，恐有出乎意料之奇效。老夫……心悦诚服！”

能让独孤明亲口说出“心悦诚服”四字，分量之重自然不言而喻。

\vspace{2em}

最后，那位一直带着和煦笑意、仿佛只是个旁观者的万仙盟盟主吴东浩，终于开口了：

“叶先生，今日观先生论道、试招，风采令人心折。然，老夫有一问，藏于心中久矣，还想请教先生。”

他轻摇折扇，目光温和却深邃地看向叶牧谦：“所谓修士，无论所修何道，说到底，皆是逆天而行、与天争命之人。既如此，大多心有所求，志有所向。或求长生久视，超脱轮回；或求无上力量，掌缘生灭；或求自在逍遥，无拘无束……不知叶先生，您创立此迥异于常的问道途，自身踏上此道，孜孜以求，所寻的……究竟是个什么呢？”

是啊，你另辟蹊径，弄出这么一条与众不同的路，总得有个终极的追求、一个能支撑你走下去的“答案”吧？

阁中所有的目光，再次汇聚于叶牧谦身上，等待着他的回答。

叶牧谦闻言，微微一怔。这个问题，他或许在夜深人静时也曾自问过，却从未如此正式地需要向他人宣之于口。

当吴东浩问出的瞬间，过往的种种画面——穿越之初，他的茫然、逃避，以及最后的悔悟与接受，如同浮光掠影，在他脑海中飞快闪过。

没有长篇大论的思考，没有修饰言辞的迟疑。

或许是前世儒家经典读多了，他几乎是凭借着一股发自肺腑的本能，那个答案便自然而然地涌上心头，冲口而出。声音不大，却清朗如玉石相击，清晰地回荡在每一寸空气都仿佛凝结的道场中：

“只为仰不愧于天，俯不怍于人罢了！”

此言一出，偌大的道场内，陷入了刹那绝对的寂静。

包括沈瑶在内的很多人似乎不能立即理解这十个字。但在白言冰和陈千羽看来，这十个字却做不得伪。

这确实是师尊在黑风岭一役之后、正式把自己当作“师尊”之后的本心写照，而师尊也一直践行着这十个字。

“仰不愧于天……俯不怍于人……”吴东浩低声重复着这十个字，眼中那惯常的和煦笑意渐渐变得深刻，更添了一丝了然与感慨的意味。

他缓缓收起折扇，在掌心轻轻一敲，微微颔首，显然对这个看似朴素至极、却又重若千钧的答案，甚为满意，亦深以为然。

王丹眸中若有所思，似在品味其中返璞归真的道韵。

陈炎却是眉头再次皱紧，嘴唇动了动，似乎觉得这个答案过于“平凡”，不够宏大，不够有“追求”。

刚刚试剑完毕的独孤明长老，则是手抚灰白长须，缓缓点头，脸上露出了深以为然的赞同神色。他毕生修心剑，守护为念，最能体会这“不愧不怍”四字背后，所需的那份坦荡、坚定与无悔。

而端坐于东首主位的沈天穹，此刻缓缓地、极其庄重地站起身来。

一声仿佛包含了无数感悟、解脱了某些无形桎梏的喟然长叹，从他口中悠悠吐出。那叹息声厚重如大地，却又带着一丝豁然开朗后的明澈光芒：

“老夫……修道三百余载。”

他声音不高，却字字如磬，说出了那四个注定将伴随“问道途”传扬天下的字：

“今日，方知……道在低处。”

道在低处。

四字如晨钟暮鼓，余韵悠长，在问道阁的梁柱间萦绕不绝，更重重地敲在每个人的心头，激荡起层层涟漪。

沈天穹，内域五大超级势力之首的天师府府主，法相境外显期的大能，以他的身份、修为与三百载见识，亲口为今日这场关乎道途立论的辩难与印证，做了最终的总结。

其含义，不言而喻——他认可了叶牧谦的道，认可了那“仰不愧于天，俯不怍于人”的朴素追求之中蕴含的至高真意，认可了这条不假外求、内修己心、看似“在低处”的新途，自有其巍然高度与无限可能。

自此，叶牧谦所传所行之道，在内域顶尖势力的层面，获得了正式的、无可争议的承认。

它叫“问道途”。

与那流传万古、显赫辉煌的“灵枢径”并立，成为内域修行大世之中，一个崭新、独特、再也无法被忽视的道途之名。

\vspace{2em}

论道会后，沈天穹私下寻到叶牧谦。

“叶先生，”沈天穹神色温和，全无府主的架子，“今日论道，瑶儿也在旁听。老夫见她神情，已知她确实找了个好师傅，心中甚慰。”

他顿了顿，道：“听闻先生至今尚无固定道场，游历讲学，虽显洒脱，终非长久之计。

“老夫在城外二十里，云霞山半山腰处，有一旧居，名曰‘云隐别院’。那里风景秀丽，灵气内敛，环境清幽，正合先生修心养性之需。那院子已空置十多年，留着也是无用，便送给先生，权当是瑶儿的束脩之礼，还望先生莫要推辞。”

叶牧谦略一沉吟，便拱手谢过：“府主厚赠，叶某却之不恭，便代小徒瑶儿谢过了。”

他确实需要一处稳定的根基之地，来梳理愈发庞杂的道途体系、顺便教导弟子。云隐别院听起来正合适。

不久后，叶牧谦携三位弟子入住云隐别院。院落果然清幽雅致，背靠云霞山主峰，面朝开阔山谷，云岚缭绕，草木葱茏，灵气平和；更设小演武场、藏书阁、闭关静室等必要设施，还有温泉、冷泉各一眼。

确是一处宝地。

\vspace{2em}

安顿下来后，叶牧谦也开始研究起他的下一境界。

思路是显然的，从论道会上陈炎的诘问说开去：不论是腾空飞行，还是远程攻伐，具体的内容实际上都是所谓“术”的层次。

而“理解”这些术法，自然是一门很大的学问。

“世事洞明皆学问”，叶牧谦忽然想起《红楼梦》中的这句话，于是自然而然地把“洞明”作为了下一境界的备选名称。

转念一想，“洞明”又不好，有管中窥豹、不能窥见全貌之嫌疑，过于小家子气；不如改成“通明”——万法通明之境，多大气磅礴！

意外之喜是“通明”的另一个意义：除了“通晓万物”，还有“十分明亮”的意思；“一心通明”，更是令人神往。

于是叶牧谦当即敲定，第三境界便叫做“通明”了：“至人之用心若镜，不将不迎，应而不藏，故能胜物而不伤”！

随即，叶牧谦定下规矩：每逢初一、十五，于别院开坛讲道，不拘出身，不问来路，有缘皆可来听。消息渐渐传开，起初多是些好奇的散修、附近小宗门不得志的弟子，以及一些仰慕沈瑶——身份终是瞒不住——或对“问道途”好奇之人前来。

叶牧谦讲道，深入浅出，不拘泥于具体功法招式，多讲“养气”之静、“见独”之悟，更反复强调“内求己身”、“道在日用”之理。

这对那些受困于资源匮乏、在传统道途上难有寸进的散修和小门派弟子而言，这无异于打开了一扇全新的窗户。听道者日众，不少人深受触动，甚至有人在别院附近结庐而居，以便时常请教。

其中有一散修，名叫陈默，原是灵枢境凝枢期修为，因资源问题困于瓶颈多年。

他听叶牧谦讲道，结合自身对灵枢径的理解，竟突发奇想，冒着巨大风险，尝试引导自身那微弱的灵枢之力下行，冲击丹田之下区域，意图强行开辟“气海”。

过程自然痛苦万分，几经险死还生；但凭借一股狠劲与冥冥中的一点灵光，他竟然成功了！

虽然修为几乎尽废，从头开始，但确确实实强行压出了第一缕问道途的元气，踏入了养气境门槛。

陈默的成功案例，引起了已成功踏入养气境界不久、正在稳固修为的沈瑶的极大关注。

她敏锐地意识到，这种破而后立的转修方式，虽然凶险痛苦，不适合大多数人，但或许是一条能让部分前路无望的修士改换门庭的潜在路径。

她详细询问了陈默的感受与过程，细心记录在册，归档整理，以备他日深入研究或谨慎参考。

这也算是问道途与旧体系之间，第一次非官方的、个体尝试的实验了。

\vspace{2em}

随着人气渐旺，为便于管理，避免纷扰，叶牧谦与三位弟子商议，制定了简单的《问道戒律》，约束居于别院周边及听道者的言行。戒律第一条，赫然便是叶牧谦在论道会上的那句心声：

仰不愧于天，俯不怍于人。

以此为核心，衍生出数条诸如“不得恃强凌弱”“不得同门相残”等基本规范。

云隐别院及周边，逐渐形成了一个以问道途理念为核心，氛围相对平和清正的独特小生态，成了不少厌倦厮杀争夺、或寻求新路的修士心目中的一方净土。

实际上，叶牧谦开堂讲道，也是试图通过广泛招收记名弟子、扩大讲道范围，来获得更多的道行反哺；然而，一段时间下来，他发现效果微乎其微。那些听道的散修、访客，即便有所得，甚至如陈默般成功转修，这种反馈也极其微弱缥缈，远不能与白、陈、沈三人带来的那种清晰触动相比。

他心中渐渐明确：疑似只有白言冰、陈千羽、沈瑶这三位命运交织、心性纯粹且正式拜入门下的亲传弟子，当他们真正开悟、突破，深刻践行问道途时，才能产生那种独特的、对他有实质性助益的“道行反哺”。

这似乎与师徒间深刻的道统传承与心神共鸣有关，不能简单通过人数叠加扩增。

\vspace{2em}

云隐别院的时光，在讲道、修行与日渐兴旺的人气中，平静地流淌了数月。

山间云雾聚散，别院道场上的听道者来了又走，走了又有新人来，别院周围依山势错落搭建的庐舍也多了十余间，俨然有了一个小型聚居地的雏形。

白言冰巩固着“见独”境界，剑意愈发凝练通透；陈千羽医术与修为并进，温润元气滋养着别院内外受伤或体虚的访客；沈瑶则彻底褪去了过往的尖锐，在养气境的稳步前行中，气质愈发沉静，只那双眸子里的倔强与清澈，未曾改变。叶牧谦则沉浸在对炼器、阵法古籍的钻研中，为问道途那众所周知的短板寻找补全之道。

目前解决“缺少功法”问题的权宜之计，依然是通过改编灵枢径的功法来实现。但两个道途存在根本的不同：灵枢径的灵力是燃料，问道途的元气是润滑油。这就直接导致了从灵枢径改编来的功法，消耗都有些不切实际的大。

路是一步步走的，移山开路这种事情，的确急不得——急，也没用。

\vspace{2em}

然而，这世间并非处处讲理，人心也并非总是开明包容。

天师府论道会上，问道途获得正式承认并与灵枢径并立；但这份荣耀背后，是潜藏的危机与既得利益者的不安。

第一个坐不住的，正是紫阳宗副宗主陈炎。

论道会归来后，叶牧谦那句“为何非要飞”的诡辩虽一时噎住了他，但并没有打消他心中的疑虑；这种疑虑反而随着时间发酵，变成了更深的忌惮。

他亲眼见过白言冰接下独孤明三招时展现的那种奇异领域，也听闻了云隐别院讲道日益红火，甚至出现了陈默那种敢于自废灵枢、转修元气的狠人散修。

一种隐约的威胁感，如同阴云般笼罩在他心头。

这一日，紫阳宗内，陈炎独坐于烈焰熊熊的炼器室中，面前悬浮着一件即将完成的剑胚，灵光流转，但他却有些心不在焉。指节无意识地敲击着座椅扶手，发出沉闷的声响。

“内求己身……不假外物……”他低声重复着叶牧谦的话，赤红的眉头越皱越紧，“若是人人皆内求己身，温养那劳什子‘元气’，对法器、丹药的依赖必然大减。长此以往，我紫阳宗炼器之道，丹鼎阁的丹药生意，根基何在？我们这些以‘外物’立宗的宗门，又何以存续？”

这个念头一经浮现，便如同毒藤般缠绕上来。

在他看来，“问道途”宣扬的理念，本质上是在动摇现有修炼体系的物质基础，是在掘他们这些传统势力的根基！

长此以往，问道途必然会逐步挤占灵枢径的生态位，最终成为新的正统！

他也确实是看过不少话本，在那些话本中，反派总是很可笑地让主角成长起来；而在反派最终选择剿灭主角时，主角往往已成气候，动不了了，只能眼睁睁地看着自己落败。

为了宗门的利益与发展，必须立即将这问道途扼杀在摇篮中。

他不再犹豫，立刻启动了几道远程传讯，分别联系了丹鼎阁阁主王丹，以及无极剑派外剑一脉的长老刘长靖。

传讯的内容大致相同，核心便是他那充满忧虑的质问：“若是人人自求，那还要我们法器、丹药何用？我等宗门，何以立宗？”

很快，他收到了回复。

\vspace{2em}

丹鼎阁，王丹静坐于氤氲着药香的丹房之中，指间捻着一枚新炼成的丹药，丹纹流转，宝光莹莹。接到陈炎的传讯，她雍容的面庞上掠过一丝沉吟。

作为丹道大家，她比陈炎更清楚丹药对于灵枢径修士快速提升、突破瓶颈、疗伤祛毒的重要性，几乎是不可或缺的辅助。若真有一种道途能大幅降低对丹药的依赖，哪怕只是对底层修士有效，对丹鼎阁长远的损失也是不可估量的。

尽管论道会上她对叶牧谦印象不坏，但涉及宗门根本利益，陈炎说得又在理，她的天平迅速倾斜。

她回复陈炎，虽未多言，但“深以为然”四字，已表明了态度。

\vspace{2em}

无极剑派，坐落在群山之巅，剑气凌霄。

门派之肇始可以上溯到一对道侣，夫以清修证道、妻以杀伐护道，道侣之间举案齐眉、相敬如宾，其道也是互补融合，成了一段佳话。至于结局，有人说他们飞升上界了，也有人说他们相继去世后葬在一起，但他们留下了无极剑派这一点是肯定的。

不知为何，他们两人创立无极剑派的时候，就没有设立“掌门”这个职位，具体事务由夫妻二人——之后便是分别统领两脉的两名长老——共同负责。他们的本意，也自然是让两脉互补融合，心剑为魂，外剑为骨，成就无上剑道。然而数代传承下来，现实却不如人意。不设掌门，使得两脉之间缺乏一个粘合剂，最终因理念差异，渐渐疏远，甚至产生了隔阂与摩擦。

传到现任这一代，心剑脉长老是独孤明，外剑脉长老则是刘长靖。两人年纪相差一轮，独孤明为长。

刘长靖对这位修为高深、剑心通明的前辈始终保有几分尊重，但两人之间的理念冲突与互相批评指摘，几乎已是常态。独孤明认为外剑一味追求威力，失了剑道灵性，易入魔道；刘长靖则认为心剑空谈心境，修行凶险，易生心魔。

问道途的出现，让两脉的分歧更加鲜明。独孤明自论道会归来后，对问道途推崇备至，尤其是对其“内明己心”、“仰不愧天俯不怍人”的核心理念，以及叶牧谦那奇特的见独领域在心神交锋中的妙用，视若珍宝。他认为这与心剑“以心证道”的追求有异曲同工之妙，甚至能弥补心剑某些过于难以把握的短板，于是开始频繁带着门下几位核心弟子，前往云隐别院拜访叶牧谦、白言冰，交流心得，探讨剑意与心境的融合。

这一切，看在刘长靖眼中，却变成了越来越深的忧虑与恐惧。

刘长靖自幼在无极剑派长大，天赋卓绝，于外剑一道上勇猛精进。但他童年时，曾亲眼目睹过不止一次惨剧：有修习心剑的师兄、师叔，或因执念过深，或因外魔侵扰，在参悟高深心法时骤然走火入魔。

平日里温文尔雅的同门，瞬间变得癫狂嗜杀，六亲不认，剑意化作纯粹毁灭的疯狂，在门派内造成过不小的伤亡——这在一位年轻人心中，会留下何等的心理阴影，不言而喻。那些染血的回廊、同门惊惧的眼神、师长痛心疾首的叹息，如同梦魇，深深烙印在他心底。

他尊重独孤明，也钦佩其剑道修为，但内心深处，始终对心剑一脉那玄妙而凶险的修行方式抱有极大的戒惧。

如今，独孤明非但不坚守心剑传统，反而去追捧一个来历不明、体系迥异的问道途，还带着弟子频繁接触之，在刘长靖看来，这简直是双重的危险叠加！他害怕独孤明会被那套新奇理论带偏，更害怕门下弟子受其影响，重蹈昔日走火入魔、同门相残的覆辙。

这种恐惧，随着独孤明拜访云隐别院的次数增多而日益加剧，让他坐卧难安。

于是，当陈炎那充满煽动性的传讯到来时，刘长靖几乎没怎么犹豫，便就着陈炎那条思路去了：他觉得陈炎说得对，这种蛊惑人心、动摇根基的“异端”，确实需要遏制。

而且，遏制问道途，某种程度上也是在将独孤明和心剑脉从危险的道路上拉回来，更是为了门派内部的纯净与安全。

于是，他先是私下找到独孤明，试图劝说：“独孤师兄，那叶牧谦之道，终究是外人旁门，与我无极剑派心剑传承并非同源。您如此推崇，还带弟子常去，恐惹人非议，也易让弟子们心生迷茫，根基不稳啊。”

独孤明闻言，只是平静地看着他：“长靖，道无内外，唯有真伪。叶先生之道，于我心剑有启发印证之妙，何来旁门之说？而我带着弟子开阔眼界，明辨道心，也是为他们夯实根基。你多虑了。”

劝说无效，反而更坚定了独孤明与叶牧谦交流的决心。

刘长靖心中焦虑更甚，他仿佛已经看到了心剑脉在叶牧谦“邪说”影响下，逐渐偏离正道，最终酿成大祸的可怕场景。

思前想后，焦虑与恐惧压倒了对同门师兄的最后一丝情面，与“门派内部事务应内部解决”的原则。一个极端而危险的念头在他心中成型：不能让心剑脉继续堕落下去了，必须采取断然措施，清理门户，将危险扼杀在萌芽中！

但仅凭外剑一脉，想要压制乃至清理心剑脉，力量未必足够，且容易引发门派内大规模内战，导致两败俱伤，这更不是他想看到的。

不过，他随即想到了陈炎的传讯，想到了王丹的偏向态度。

随即，计划逐渐清晰：联合外部的紫阳宗，以雷霆之势，清除心剑脉中的不稳定因素，至少要让独孤明放弃对问道途的支持，将心剑脉重新拉回正轨。这样，即使心剑脉拉不回来、随之覆灭，外剑脉也不会受到过多损失，似乎是一个较好的方案。

于是他秘密联系了陈炎，表达了“清理门户”的意向，并请求紫阳宗协助，施加压力。

陈炎接到刘长靖的密讯，先是一愣，随即心中狂喜。

他正愁找不到合适的机会和理由打压日益壮大的云隐别院和问道途，刘长靖的请求简直是瞌睡时有人递枕头！

插手无极剑派内部事务，固然有风险，但收益巨大。

一方面，可以借此狠狠震慑叶牧谦——看，连支持你的独孤明和他的心剑脉，我们都敢动，你掂量掂量自己的分量！

另一方面，若能协助刘长靖掌控或重创心剑脉，等于从无极剑派这块大蛋糕上撕下了一块肉，极大增强了紫阳宗的影响力，甚至可能在未来掌控部分无极剑派的资源。

这个刘长靖有点傻啊，这笔买卖，在他看来实在是划算得不得了。

他几乎没有任何犹豫，便欣然答应了刘长靖的请求，双方密谋定下了行动细节。

\vspace{2em}

数日之后，一场突如其来的“清理门户”行动，在无极剑派猝不及防地爆发了。

这一日，天色阴沉。刘长靖以“稽查功法异动，防止心魔滋生”为名，调集了外剑脉大批精锐弟子，在数位紫阳宗高手的暗中配合下，突然包围了心剑脉主要的几处修行静所和讲剑堂。

没有过多的宣告，也没有给独孤明辩解的机会。刘长靖手持象征着外剑脉长老权限的令牌，面色冷峻，声音通过灵力传遍山峦：

“心剑脉近年修行路数有偏，易生魔障，为保门派清净，现需暂停心剑脉一切对外交流，封闭静修，接受审查！所有弟子，即刻返回各自居所，不得随意走动！违令者，以叛门论处！”

“放屁！”

一声怒喝如惊雷炸响。

独孤明果然出现，飞在半空与刘长靖对峙；须发皆张，青衫无风自动，背负的古剑发出嗡嗡清鸣。

他没想到刘长靖竟敢如此公然发难，更没想到对方竟然勾结了外宗的紫阳宗之人。这已不是简单的理念之争或内部整顿，这是赤裸裸的背叛与引狼入室！

“刘长靖！你竟敢勾结外人，对我心剑一脉下手！无极剑派的列祖列宗在上，岂容你如此胡作非为！”独孤明目眦欲裂，心中充满了被同门背后捅刀的悲愤。

“独孤师兄，我这是为了门派好，也是为了你好！”刘长靖面对独孤明的怒火，心中虽有一丝愧疚，但更多的是一种“执行正义”的决绝。

“只要你答应不再与那云隐别院来往，并约束弟子，接受审查，我等自会退走，此事尚有转圜余地！”

“转圜余地？将我剑心自由，卖与你等做妥协的筹码吗？”独孤明怒极反笑，笑声苍凉。

“我独孤明修剑一生，求的是心念通达，护的是心中正道！今日你等以势压人，以奸谋算计同门，此等行径，与邪魔何异！要我屈服？宁可战死！”

最后一个“死”字吐出，浩瀚纯粹的剑意冲天而起，充满了玉石俱焚的决绝与凛然不可犯的锋芒。

他身后，数十名心剑脉核心弟子也纷纷拔剑，剑意虽强弱有别，但那份宁折不弯、誓与师长共存亡的心志却连成一片。

冲突瞬间爆发。

外剑脉弟子剑招凌厉，配合默契，道道剑气纵横切割，充满杀伐之气。紫阳宗高手则在暗中释放火系法术干扰，或祭出法器远程轰击。

心剑脉弟子剑意纯粹，往往于间不容发之际寻隙而入，攻敌必救，或以精妙剑意干扰对手心神；独孤明更是独战刘长靖与一名紫阳宗真符境长老，剑光如龙，竟一时不落下风。

然而，终究是寡不敌众，且事发突然，心剑脉并未做好准备。激战片刻，便有数名心剑脉弟子受伤。独孤明眼见弟子受创，又感知到还有更多外剑脉弟子和紫阳宗援手正在赶来，心知今日事不可为，继续硬拼，只会让心剑脉传承断绝于此。

他悲啸一声，剑光暴涨，暂时逼退对手，对门下弟子厉喝道：“心剑一脉弟子听令！即刻分散，突围下山！留住有用之身，剑心不灭，自有再见之日！勿要为我这老骨头陪葬！”

“师尊！”弟子们悲呼。

“走！”独孤明剑意催发到极致，竟隐隐有燃烧本源之势；为了防止部分忠义弟子死战不退，他甚至朝着身后斩出一发剑气，逼迫他们撤退！

心剑脉弟子含泪咬牙，知道这是师尊用生命为他们争取生机，不得不化整为零，凭借对地形的熟悉和灵活的剑意运用，从各个方向拼死突围。

大部分人在独孤明拼死掩护和混乱中，成功冲出了包围圈，但亦有少数人重伤被擒或当场战死。

刘长靖本想追击，但被独孤明那搏命般的剑势死死拖住，加上他也不愿对逃亡的普通弟子赶尽杀绝——毕竟初衷是清理而非灭绝，只得眼睁睁看着大部分心剑脉弟子四散逃入茫茫群山。

最终，独孤明燃烧本源，力战至油尽灯枯，随着一声叹息溘然长逝。其余高层，不是陨落就是被擒。

曾经剑气凌霄的无极山，心剑一脉，一夜之间，风流云散。

\vspace{2em}

逃出生天的心剑脉弟子，在最初的慌乱与悲愤过后，不约而同地想到了一个地方——云隐别院。师尊独孤明正是因为推崇“问道途”，与叶先生交好，才遭此横祸。如今他们无处可去，天下虽大，恐怕也只有那个地方，可能理解他们的剑心，可能收留他们这些“叛门”之人。

于是，在接下来数日，陆陆续续有二三十名形容憔悴、身上带伤、却眼神依旧明亮坚定的年轻剑修，跋山涉水，来到了云霞山，找到了云隐别院。

他们对着闻讯出迎的叶牧谦与白言冰等人，悲声陈述了无极山发生的剧变，将刘长靖如何勾结紫阳宗、如何突然发难、师尊独孤明如何宁死不屈、他们如何突围逃散的经过，原原本本，备述一遍。

说到激愤处，不少铁骨铮铮的剑修也忍不住红了眼眶，更有人因伤势和心力交瘁，几乎站立不稳。

叶牧谦听完，沉默良久。他看着眼前这些风尘仆仆、剑心受创却依然挺直腰杆的年轻人，心中涌起复杂的情绪。有对独孤明遭遇的同情与敬意，有对陈炎、刘长靖等人狭隘与狠辣的怒意，也有一种“树欲静而风不止”的无奈。

他知道，问道途的平静发展期，恐怕就此结束了。

“诸位暂且安心在此住下。”叶牧谦最终开口，声音沉稳，带着一种抚慰人心的力量，“云隐别院虽陋，尚可遮风避雨。独孤长老之道友与弟子，便是我叶某之道友与弟子。他日若有机会，必当为你们报仇。”

无处可去的剑修们，闻言心中稍安，纷纷行礼道谢。他们在别院周边，寻了空地，或伐木，或垒石，很快便搭建起更为简陋却足够栖身的庐舍，暂时安顿下来。

\vspace{2em}

而白言冰倒是因祸得福了。在与这些剑修的相互印证和切磋中，他也终于有所体悟。

一日，他主动找到沈瑶：“师妹，我近来有些感触，不如切磋一下。”

沈瑶其时早已重回养气境大成，正随叶牧谦为云霞山布置防御阵“九宫迷踪”。叶牧谦点了点头，于是沈瑶便与白言冰走到演武场上站定。

不少修士其实正在场上苦练战技，见大师兄和三师姐——他们不少人都呼叶牧谦为“师尊”，于是白、陈、沈三人自然就变成师兄师姐了——要切磋，于是纷纷主动退到一旁。毕竟能见到见独境（战斗力甚至相当于真符境）切磋，即使是在一旁观察那也是能学到不少东西的。

“约定一下，不用领域。”白言冰说。

沈瑶秀眉微蹙，不过同意了：师兄看起来是想单纯的磨练剑法。

两人各自拿了趁手的、没开刃的切磋用兵器。白言冰依然拿得长剑，沈瑶同样选择短刀。

陈默正好走了过来，于是自然而然地成了帮两个人敲锣的人。

“当！”

锣声一响，沈瑶立即消失在原地，随即短刀直攻白言冰面门。

这是承袭了她先前还修习灵枢径时的战斗方式：精确、高效。

天下武功，唯快不破！

她与白言冰在禁用领域的情况下切磋过多次，虽然境界有差异，但依然互有胜负。

白言冰不慌不忙，身形微转，长剑忽而一举，堪堪格挡住沈瑶的短刀，长剑周身元气竟然是冰寒剑气。

能挡住？冰寒剑气？

沈瑶没见过白言冰用过冰寒剑气；在以前的大多时候，白言冰还是更喜欢刚猛的雷火剑气。

沈瑶不知道白言冰葫芦里卖的什么药，随即短刀连挥，幻出数道刀影，从不同角度疾刺而来，却全部被白言冰拦下。

场外不少修士看来，白言冰只有防守的工夫、没有进攻的余裕，似乎落入下风。

“看来这一场切磋，是沈师姐赢了啊。”场下有人窃窃私语，“你猜错了，一块灵石赶紧拿来。”

“再等等。”那位猜白言冰获胜的修士小声说，“陈默师兄还没说谁赢谁输呢。”

再看陈默，他只是笑而不语。他跟随叶牧谦多时，自然能看出来白言冰此时在等什么。

沈瑶久攻不进，她渐渐变得焦躁起来。

这白言冰耍她呢，怎么只防不攻？

手上疾速点刺、寻求破防的短刀，因焦躁而自然而然地出现了一瞬的破绽。

殊不知白言冰正等着这一刻！

下一瞬，白言冰动作幅度忽然变大，剑光沿着极为精确的轨迹，寒冰剑气骤然射出，直攻沈瑶刚刚不经意间露出的破绽！

所幸白言冰留了手，给了沈瑶退避的时间，否则沈瑶目前估计已经重伤倒地了！

胜负不言自明，先前那个支持沈瑶的修士的脸立即拉得很长，被迫掏出一块灵石给了他的朋友。

沈瑶往后一闪三尺远，瞅了白言冰一眼，忽然笑道：“白师兄好武艺，这局是我输了。以无数格挡等待破绽，随即一剑封喉……是这样的吗？用冰寒剑气是因为这种剑气更擅长防御……？”

白言冰笑道：“师妹好眼力。前半招叫‘寒光散乱’，后半招叫‘凝萃霜满’……不知道这两个名字如何？”

“好得不能再好了。”沈瑶笑道，“这两招，比你那半成品一般的‘弥焰斩’和‘直环焰’强多了；虽然威力不似那等刚猛，但实战能力却必然更胜一筹。”

\vspace{2em}

无极剑派心剑一脉遭到“清理门户”的消息如同长了翅膀，迅速传遍内域大小势力。

那些嗅觉敏锐的宗主、长老们，从这场突如其来的“清理门户”中，嗅到了不同寻常的气息。虽然说无极剑派两脉中间确实偶发冲突，但从来都是合久必分、分久必合的。

这次被赶尽杀绝，背后必然存在其余推手。

虽然天师府、万仙盟对紫阳宗的公然插手表示谴责，但矛头隐隐指向的云隐别院和其代表的问道途，无不表明一场围绕修行理念、资源分配乃至未来格局的深刻对立，已经浮出水面，且以如此激烈的方式拉开了序幕。

站队，成了许多中小势力迫在眉睫的选择。

审时度势之下，阵营迅速分化、成形：

一方，是以紫阳宗副宗主陈炎为核心，联合了丹鼎阁阁主王丹、无极剑派外剑脉长老刘长靖的守旧联盟。他们代表着内域现有的、以灵枢径为基础、以法器丹药为重要辅助、各大宗门垄断核心资源的修炼体系；其拥趸也大多是与此体系深度绑定、依赖现有格局生存的宗门、商会及相关利益集团。

另一方，则是以云隐别院叶牧谦为核心，并获得了万仙盟盟主吴东浩公开支持的新生势力联盟。这一方吸引了大量在原有体系中难以出头、资源匮乏的低级散修，许多受大宗门挤压、寻求出路的小型宗门和家族，以及那些对问道途理念心生向往、或单纯厌恶守旧联盟霸道行径的修士。

沈天穹在公开场合仍选择保持超然的中立姿态，呼吁双方克制，调停冲突。但其立场明显偏向叶牧谦一方：毕竟白言冰和陈千羽再怎么说也算是小天师，与天师府的关系极为密切。

沈天穹私下里对云隐别院的回护与默许，明眼人都能看出。

\vspace{2em}

围绕着云隐别院与守旧联盟的，是一场压抑而充满张力的冷战。

大规模的热战并未立即爆发，无论是出于对双方可能实力的忌惮，还是对新生势力那种未知韧性的评估，双方都保持着一种危险的克制。但这绝非和平，而是在厚重冰层下蓄积暗流，寻找更能一击致命、且不必担负首要开战罪名的打击方式。

很快，一场没有硝烟却足以扼杀生机、更加残酷的博弈，悄然拉开了帷幕。

数日后，一道盖有陈炎宝印、措辞空前严厉的“宗门戒令”，通过其掌控的渠道，迅速传遍内域各大主要坊市与势力：

“即日起，凡我紫阳宗弟子，曾私自前往云隐别院听道、论法、交流者，概视为受异端邪说蛊惑，心志不坚，道基已污！限三日内至刑律殿自陈过错，听候发落。逾期或隐瞒不报者，一经查实，即刻废去修为，逐出宗门，名录剔除，永不收录！

“自今往后，紫阳宗上下，严禁任何弟子、执事、长老与云隐别院之人有任何形式的往来、交流、贸易、传讯，违者以叛宗论处，严惩不贷！”

这道戒令将云隐别院及其代表的问道途，正式定性为“异端”、“禁地”。

它试图斩断一切明面上的人际纽带，将叶牧谦一方置于一个被孤立、被污名化的孤岛之上。

戒令在内域舆论场上激发的余波尚未平息，丹鼎阁与已由刘长靖完全掌控的无极剑派（外剑一脉）便迅速跟进，发布了措辞几乎雷同的禁令。三大超级势力的联合封杀，形成了一道令人窒息的无形壁垒。

威慑是立竿见影的。

许多原本与云隐别院有些零星交易往来、或纯粹出于对新生事物好奇而前去旁听过的修士、小商会代表，纷纷噤若寒蝉。

原本络绎通往云霞山的小道上，新面孔骤减；偶尔有目光投向别院方向，也迅速移开，仿佛那风景秀丽的半山腰弥漫着无形的疫病。除了早已定居在此的散修和那些从无极剑派逃难而来、无处可去的心剑脉弟子，云隐别院周围，确实显出一种被主流世界刻意疏离的冷清。

\vspace{2em}

然而，真正能扼住一个势力命脉的，不能仅仅凭人际上的孤立。

一切问题本质上都是经济问题；从古至今，无论中外，乃至这个异世界，资源与实业永远是统治的根基，也是最锋利的武器。一场精心策划、目标明确的涨价风暴与“货源调控”，在看似正常的市场波动掩护下，猛烈袭来。

首先遭殃的是公开坊市。

在紫阳宗、丹鼎阁影响力深厚的几处大型坊市，流向非守旧联盟认可势力——尤其是与云隐别院有牵连者——的各类基础修炼资源，价格开始诡异地、毫无缘由地连续飙升。

“赤铜精铁，昨日还是十灵石一方，今日怎就十五了？” 一个试图为同伴购买铸剑材料的心剑脉弟子，对着摊位后的掌柜愕然道。

掌柜眼皮都未抬，懒洋洋道：“货源紧张，上游提价，爱买不买。对了，看阁下气息，似是云霞山那边来的？若是，价格还得再加三成，毕竟……风险大嘛。” 话语中的奚落与歧视，毫不掩饰。

类似的场景不断上演。沉银砂、寒铁木等常用炼器辅料，化瘀草、宁神花等基础疗伤药材，涨幅普遍达到三到五成，个别紧要物资甚至翻倍。

明眼人都能看出，这显然不是市场的自发调节，而是一种精准的、带有惩罚性质的宏观调控，是一种歧视性的、甚至制裁性的定价。

紧接着，打击蔓延到了更基层、更关乎日常修炼生存的领域。

一些云隐别院周边散修聚居地日常所需的低阶资源，如最基础的凝气草、铁精矿，开始出现诡异的货源紧张。原本定期前来兜售的行商不见了踪影；本地几个小供货点要么关门，要么货架上空空如也，偶有存货，价格也已高到令人绝望。

“王管事，这个月的凝气草……” 一个面容愁苦的老年散修，小心翼翼地问道。

被称为王管事的胖商人叹了口气，压低声音：“老李头，不是我不帮你。上头打了招呼，凡是往云隐别院那边走的货，一律卡死。我这儿但凡卖你一株，明天我这铺子就别想开了。你……另寻门路吧。” 话语间，他眼神躲闪，满是无奈与畏惧。

此计堪称毒辣，虽兵不血刃，但确实击中了众多问道途依附者的要害；其目的阴损而深远，绝非仅仅是提高成本、大赚一笔那么简单。

首要目标，自然是尽可能地消耗新生势力联盟本就不甚宽裕的财力。

对于已经开始转修或兼修问道途，如白言冰、陈千羽这般修士而言，对外部灵石、丹药的依赖极低；但是新生势力联盟之下，还有大量的、需要资源的灵枢径修士：无论是心剑脉弟子维护保养剑器，散修们维持基本修炼，还是尝试自给自足进行低阶炼丹、炼器，都离不开这些基础资源。

其次，这是一次强大的威慑与站队逼迫，做给内域所有摇摆不定的中小势力看：站在守旧联盟一边，你们的资源供应渠道将得到保障；但若与云隐别院、与问道途扯上关系，哪怕只是一点生意往来，等待你们的将是同样的“卡脖子”命运。

修炼资源断供，商会贸易受阻，在这种依靠外物的修行界必然寸步难行。

这是一道冰冷的选择题。

价格飞涨与货源断绝，如同两只缓缓收紧的铁钳，要让叶牧谦他们那宝贵的灵石储备迅速见底，让日常修炼陷入停滞，让刚刚凝聚起来的人心因生存困境而涣散。

一名年轻的心剑脉弟子，看着手中因崩出细小缺口的长剑，眉头紧锁。修复材料的价格早已涨上了天，而且有价无市。他跑遍了附近三个小镇，一无所获。

几个常驻的散修围坐在一起，面色晦暗。他们修为低下，仍需依靠凝气草辅助引气，往日攒下的一点灵石，如今连购买往常三分之一的份量都不够。修炼进度近乎停滞，焦虑与绝望的情绪在蔓延。

甚至连云隐别院内部，也开始感受到压力。虽然叶牧谦讲道不依赖于外物，但聚集而来的数百人中，大半仍需基础资源维系。

别院原本的少量储备正在快速消耗，负责后勤的弟子面带忧色，向叶牧谦汇报着日渐窘迫的处境。

无形的壁垒在人际间筑起，有形的资源枷锁则在经济命脉上狠狠绞紧。云隐别院仿佛狂风恶浪中的一叶孤舟，虽然核心依然稳固，但船体正在承受一波接一波的侵蚀与挤压。

冷战阴云笼罩下的云霞山，看似平静，实则每一个角落都弥漫着生存的艰难与抉择的重量。

而这场没有刀光剑影的围剿，才刚刚开始。

\vspace{2em}

经历此事，叶牧谦忽然觉得他把人想得太好了。

自己穿越时正值中美两个庞然巨物相争，如此情景真有一种前世的感觉。

前世记忆中，因为中国工业品类齐全，美国只能找到高端芯片等少数几项技术去卡中国的脖子；但叶牧谦手下的这套烂摊子，基本上和八十年前、百废待兴的中国没什么两样，什么都没有，什么都造不出来。

中国当时至少还有北方老大哥的援助，可云隐别院呢？

“造不如买”这种思维还是太诱人了，殊不知这是把命脉悄然交到了别人手上……关系好的时候，资源流通、互通有无，那无疑是好的；但剑拔弩张时，双方鼻青脸肿、头破血流，谁还跟你自由贸易？

“前辈们还真是高瞻远瞩啊……”叶牧谦长叹道。

但生存确实是一个迫在眉睫的问题。单纯的防御与抱怨无济于事，必须主动破局。

叶牧谦深知，己方的根本优势在于问道途自身对于外部资源依赖极低的韧性，而自此衍生出来的重要优势自然是汇聚人心、团结力量办大事。

他将吴东浩、白言冰、陈千羽、沈瑶等核心人物召集起来。没有冗长的争论，一场务实而高效的商议迅速展开。针对守旧联盟的封锁，一套立足自身、着眼长远的反制体系被勾勒出来，并立即付诸行动。

核心思维分两步走：

第一步，用尽婉转手段，在封锁没有彻底合围的时候，尽可能打破禁运的死局，钻空子、走私，至少先保住自己境内许多灵枢径修士的日常生活；

第二步，则是拉动内生产，降低物资依赖程度，最终达到部分基础资源的基本自给。

万仙盟盟主吴东浩当机立断，派出了自己的儿子——吴刚。这位少盟主骨龄不过四十，外貌则依然保持在二十岁左右的样子，修为已达玄脉境融贯期。此人看似潇洒不羁，实则心思缜密、长袖善舞，在万仙盟遍布内域的松散人际网络中拥有极高的信誉与号召力。

吴刚领命后，第一个找到的帮手，便是成功转修问道途、对云隐别院归属感极强且常年混迹散修圈子的陈默。

一个代表组织力量，一个深谙草根生态，两人的结合立刻迸发出高效的能量。

他们避开守旧联盟严密监控的大型坊市与主要商路，将目光投向了那些常年被忽略的边缘角落与灰色地带。

依托万仙盟那张看似松散却触及每个角落的人脉网，吴刚与陈默开始秘密联系那些走南闯北、熟悉每一条偏僻小径与隐秘山谷的资深散修，联系一些因不满盘剥而私下开采边缘矿点的小型团伙，甚至联系一些在守旧联盟商会体系内不得志，或者愿意冒险赚取更高差额利润的底层管事。

一张隐秘而高效的地下交易网络，如同蛰伏于地层之下的根须，开始悄无声息地蔓延生长。这张网络分工明确，构成了一个粗糙但实用的闭环：

采购由熟悉当地情况的散修负责，从守旧联盟控制力薄弱的偏远矿区、人迹罕至的深山，以“捡漏”或小规模盗采的方式，零散获取品质可能不高但勉强可用的矿石；或是向一些与丹鼎阁没有直接契约关系的山民、采药人，或伪装成守旧联盟收货人，或采用其余手段，收购他们采集的野生草药。

另一面，在几个重要的散修聚集点或小型城镇的暗处，设立了隐秘的交接与仓储点。物资在此汇总、分类，再根据云隐别院总部及各处新生势力据点的需求清单进行分发。

交易大多采用以物易物的方式——用云霞山产的某些草药、或从旧法器上熔炼回收的金属，交换急需的矿石或其他药材，最大限度规避对灵石的依赖和金融追踪。

运输反而是最危险也最关键的一环。一批胆大心细、修为不弱且信誉良好的散修被组织起来，伪装成普通的行商、猎户甚至逃难的流民，利用多条隐秘小道进行物资转移。他们往往昼伏夜出，分段接力，并将物资化整为零，以降低单次运输的风险。

经过数次失败之后，具体的采购、运输线路也从总线联系变成了单线联系，形成彼此隔离的链条；即使某一环被发现并切断，损失也被控制在局部，不会导致整个网络瘫痪。

虽然这条地下血管输送的物资量，远无法与昔日自由市场的充沛相比，运输过程也伴随着损失和风险，但它确实像涓涓细流，顽强地将各地零星的资源汇拢起来，持续输送到急需给养的新生阵营，勉强维持了生存与最低限度的运转。

其次，是推动就地生产，降低核心依存度。陈千羽在这方面发挥了不可替代的作用。

她清醒地意识到，药材，尤其是基础疗伤和恢复类药物，是联盟修士的生命线；完全依赖外部输入，等于是将命门始终握在敌人手中。

很快，一批因为性格耿直得罪上司、或因出身寒微备受排挤、或在丹鼎阁内部权力倾轧中失势的药修，被秘密招揽或主动投奔而来。同时，聚居地中许多原本从事过农耕、或对培育植物有天然兴趣的散修也被动员起来。

在陈千羽的带领下，他们在云霞山脉几处灵气相对温和、土质适宜的背阴山谷与向阳坡地，以最原始的方式，开辟出了一片片错落有致的梯田式药圃。

种植不是简单地移植，那样太慢了。

陈千羽以问道途元气温养为基础，自己体悟的内在平衡与生生不息的理念为指导，与众多现有修士所知的传统灵植术相结合，进行了一系列因地制宜的改良。

她指导药修和散修们，注重不同草药伴生带来的生态平衡，利用山间溪流、雾气、不同朝向的光照来模拟最佳生长环境；虽然受限于环境与时间，无法种植那些需要特殊地脉、动辄数十年甚至上百年才能采收一次的高阶灵药，但诸如止血藤、宁神花、凝气草、化瘀草等最常用、消耗量最大的基础药材，自给率开始稳步提升。

每一株在云霞山泥土中壮大的草药，都减少了一分对丹鼎阁的依赖。

与此同时，一批对炼器有兴趣的散修和部分心灵手巧的心剑脉弟子，在别院一角搭起了简陋的工棚。利用地下网络艰难运来的、品质参差不齐的低阶矿石，他们从零开始，学习最基础的材料辨识、提纯和粗坯锻打，叮叮当当的敲击声终日不绝。

虽然蕴含符文的真正法器的炼制效率依然极低，但打造出质地坚韧的普通刀剑、农具，修复日常损坏的器械，甚至尝试将废料重熔再利用，都能显著减少对外界制品的依赖，并将重要的金属资源尽可能留在内部循环。

\vspace{2em}

这场没有硝烟的经济暗战，在平静的表面下激烈地进行着。

守旧联盟的探子像猎犬一样四处嗅探，试图找出地下网络的关键节点加以破坏。

双方在偏僻的山道、无人的峡谷发生过数次小规模的遭遇与冲突，互有死伤，鲜血染红过溪流，但这条隐秘的生命线始终未曾彻底断绝。

坊市中的价格博弈每日上演，守旧联盟提价、控货，地下网络就寻找替代品、开发新来源，或默默承受成本，将压力分散到整个共担风险的网络之中。

云隐别院周边的聚居地，日子过得肉眼可见的清苦。灵石永远不够用，饭菜少见荤腥，衣物修补再三。但一种奇特的凝聚力，却在同舟共济、共克时艰中悄然滋生。

确实有不少人最终撑不住，选择离去或倒戈，但留下的无不是心性坚定、对问道途抱有信任、对未来充满希望的修士们。

人们清楚地知道，守旧联盟是想用资源和饥饿逼他们就范、自行瓦解。因此，每一条穿越封锁线安全抵达的药材矿石，每一片在梯田上迎风摇曳的葱绿药苗，甚至每一把自家工棚打出的粗陋铁器，都不仅仅是一件物品，而是对那无形壁垒最直接的反抗与宣言。

希望，在这种具体的、触手可及的劳动与收获中，变得真切起来。

然而，叶牧谦比任何人都清楚，经济战的消耗是巨大且不对等的。

地下网络的运转本身就需要付出高昂的风险成本，药田的丰产需要至少一两个生长周期的等待；而守旧联盟掌控的财富与资源毕竟雄厚百倍。长期僵持的消耗战，对于新生的一方绝非良策。

而就在这时，一个意外的消息，再次改变了局势的走向。

\vspace{2em}

天枢城方向，忽然天生异象，霞光万道，瑞气千条，一股奇异的、仿佛来自另一片空间的波动传遍四方。

紧接着，天师府正式发布通告：

“新小玄界，即将开启。”

小玄界，一般指的是一些破碎的界域；它们的出现和消失毫无规律，其间或是生机盎然，或是荒芜死寂，几乎完全无法预测。内域对其成因也一无所知，但这并不妨碍每一次出现都引发巨大轰动。这意味着未知的机遇、失传的传承、珍稀的天材地宝，甚至可能是改变命运的一线曙光。即便可能伴随着莫测的危险，也足以让任何势力与修士心动。

\vspace{2em}

在持续数日的异象与全城骚动，以及秘密书信沟通后，天枢城内，一场气氛凝重的临时会议开始举行。

与会者，正是当下内域两大对立阵营的代表人物。

守旧联盟一方，紫阳宗副宗主陈炎面色冷硬，丹鼎阁阁主王丹神情淡漠，无极剑派外剑脉长老刘长靖眼神锐利如剑。新生势力一方，云隐别院叶牧谦安然端坐，万仙盟盟主吴东浩脸上惯有的和煦笑容也收敛了几分。天师府府主沈天穹作为东道主与调停者，居于主位。

空气中弥漫着无形的张力，比之前任何一次公开对峙都要紧绷。

经济战的暗流尚未平息，无极剑派的血案尤有余温，此刻却要因为一个小玄界坐在一起，场面堪称诡异。

“诸位，”沈天穹打破沉默，“小玄界开启，乃内域共有的机缘，亦是潜在的危机。若因彼此争端，在入口处便大打出手，或是在探索中互相倾轧过甚，非但可能导致机缘尽丧，更可能引发不可测的界域反噬，殃及天枢城乃至更广范围。此非吾辈所愿见。”

陈炎冷哼一声：“府主之意，是要我们与这……‘异端’合作？”

他刻意加重了最后两个字，目光如刀刮过叶牧谦。

叶牧谦眼皮都未抬一下，只是淡淡道：“道不同，不相为谋。然机缘在前，若因门户之见错失，亦非智者所为。叶某只问一句，陈副宗主敢不敢让门下年轻弟子，与叶某的弟子，在‘公平’的环境下，各凭本事，争一争这机缘？”

虽然叶牧谦是一个和平主义者，连带着整个问道途人都是内求一身、不喜争斗的性子；但此时两方已经是撕破脸皮，那叶牧谦自然也不介意与他们争一争。

“公平？何谓公平！”刘长靖接口，语气带着讽刺，“你那些弟子，修的路子古怪，谁知道会不会用些不上台面的手段？”

吴东浩呵呵一笑，摇着折扇：“刘长老此言差矣。既然是探索未知秘境，自然是手段尽出，只要不违事先约定即可。莫非贵派弟子，只会按部就班，离了宗门护佑便寸步难行？”

眼看又要吵起来，沈天穹抬手制止：“够了！”他目光扫过众人，“我们也都知道，此前已有书信沟通基础。今日便定下章程：两派搁置争议，共同探索此界。为免事态升级，仅限骨龄百岁以下的弟子进入。小玄界内，生死祸福，各凭机缘本事，外界不得干涉。出得小玄界，恩怨亦不得追溯至界内之事。诸位可有异议？”

这条件颇为苛刻，“生死不论”四字更是透着血腥。但这也最大程度避免了外界势力直接插手，将冲突限制在年轻一代的竞争层面。

陈炎等人虽然不甘，但也明白这是目前最能接受的方案，既能打压新生势力年轻一代，夺取机缘，又不会立刻引发全面战争。叶牧谦一方则需要这个机会，让弟子历练，并争取资源打破封锁。

最终，双方虽然在细则上磨了半天嘴皮子，但最终还是在“搁置争议、共同开发”的基础上基本达成了共识。

\vspace{2em}

三日后，天枢城北三百里，裂云峡谷。原本荒芜的峡谷深处，空间扭曲，那道七彩漩涡般的入口稳定了下来，高约三丈，宽两丈，内部光影朦胧，看不真切。

峡谷两侧，泾渭分明地站着两队人马，总数约六十余人。

左侧，守旧联盟阵营。紫阳宗弟子清一色赤红劲装，气息灼热，大多在玄脉境开脉期以上，为首一名面容倨傲、目蕴精光的青年，正是紫阳宗当代圣子周文彬，修为已达真符境期摹形巅峰，是场上唯一真符境修士，真不愧“内域第一天才”的称号。

丹鼎阁弟子则衣着素雅，以青、白二色为主，身上带着淡淡药香，虽不擅强攻，但各种辅助、疗伤、甚至用毒的手段不容小觑。

无极剑派弟子人人负剑，剑气凌厉，站姿如松，眼神中带着经过“清理门户”后愈发浓烈的排外与敌意。

右侧，新生势力阵营。云隐别院一方，白言冰立于最前，气息沉凝，见独境界的圆融之意让他即便在众多灵力勃发的对手面前也不显弱势。陈千羽站在他身侧稍后，温婉宁静，腰间悬挂针囊。沈瑶则位于另一侧，黑衣飒爽，容颜清丽，眼神锐利地扫视着对面。

他们身后，是十几位修为在养气到见独初期不等的云隐别院核心弟子，以及数位自愿参与、经过筛选的散修好手。万仙盟一方，则由吴东浩之子吴刚带领，约二十余人，服饰杂乱却个个眼神精悍，透着散修特有的野性与机警，陈默也在其中。

除此之外，还有少部分天师府弟子，不过十人。他们保持着中立，在剑拔弩张的两派之间，显得有些格格不入。

这些弟子们事先也都签了生死状。

峡谷上方，沈天穹、陈炎、王丹、刘长靖、叶牧谦、吴东浩等人的身影在云层中若隐若现，强大的神念笼罩着下方，既是相互监督，也是防止有人破坏约定。

“入口已稳，时限一个月。一个月后，无论有无收获，必须退出，否则界域闭合，神仙难救。”沈天穹恢弘的声音传入每个弟子耳中，“现在，进入！”

话音落下，两侧队伍几乎同时动了！

没有寒暄，没有试探，最初的冲突在踏入秘境的第一步就已爆发！

守旧联盟弟子显然早有预谋，数名紫阳宗弟子在靠近入口时，突然齐齐催动灵力，炽热的灵力并非攻向对手，而是猛然轰击在入口边缘不稳定的空间波纹上！

“轰！”

本就动荡的入口顿时一阵剧烈扭曲，七彩光华乱窜，形成一股混乱的灵力乱流和空间撕扯力！他们的目的很简单，制造混乱，干扰甚至重创即将进入的新生势力弟子，最好能让他们分散传送到危险区域。

“卑鄙！”白言冰眼中寒光一闪，一股圆融稳固的势场以他为中心微微扩散，将身旁最近的陈千羽、沈瑶以及几名弟子笼罩。

与此同时，吴刚那边经验丰富的散修们更是各显神通，有的身法如烟，瞬间穿过乱流薄弱点；有的祭出粗糙但实用的防御符箓硬扛；陈默则闷哼一声，体内初成的元气鼓荡，护住周身，强行冲入。

乱流之中，人影纷乱。大部分人都成功冲进了入口，但仍有数名新生势力的散修和一名云隐别院弟子被乱流扫中，惨叫着被抛飞出去，受伤不轻。而守旧联盟那边，也有两名弟子因操控灵力过猛，自己也被反震所伤。

\vspace{2em}

光芒一闪，天地变换。

经天师府先头部队初步探查，小玄界与上古神鸟“青鸾”相关，姑且称之为青鸾秘境。

青鸾主生机平衡，在内域尚未绝种的时候，基本都能达到法相境界甚至传说中的更高境界，其留下的传承预计是极为可观的。

白言冰稳住身形，发现自己站在一片苍翠欲滴的古老森林边缘。古树参天，藤蔓垂落如帘，空气中弥漫着浓郁到化不开的草木清香与精纯的生机灵气，吸一口都让人精神一振。远处有鸾鸟清鸣般的回响，空灵悠远。

这秘境环境，竟是难得的鸟语花香、生机盎然之所在。

他迅速环顾四周，心中稍定。陈千羽、沈瑶就在他身边不远处，另外还有八名云隐别院弟子和五名万仙盟的散修（包括吴刚和陈默）也陆续在附近出现。进入时的默契配合，让他们这一批近二十人的队伍得以相对集中地传送到同一片区域。

这已是极大的优势。

“清点人数，检查伤势。”白言冰沉声下令。

很快，结果出来。他们这一队十六人，除了两人在入口乱流中受了轻伤，并无大碍。

陈千羽立刻上前为伤者紧急处理伤势。

“不可大意。”沈瑶警惕地观察着周围，“青鸾主生，亦可能意味着其考验或守护力量，皆与生机相关；此地生机盎然，妖兽自然也不可能少了。而且那些家伙，也肯定进来了。”

她话音刚落，左侧密林中便传来一声急促的惨叫，紧接着是灵力爆发的轰鸣和怒喝声！

显然是其他分散传送进来的弟子遭遇了袭击，不知是秘境本身的危险，还是人为。

“走，过去看看，小心戒备。”白言冰握紧剑柄。众人立刻结成简单的阵型，白言冰在前，沈瑶在侧翼，陈千羽和伤员在中间，吴刚、陈默等散修经验丰富，负责殿后和探查周围。

十几个人向着声音来源谨慎靠近。

没走多远，他们便看到了惨烈的一幕：一片林间空地上，三名服饰各异的散修——似乎是追随万仙盟的小势力弟子——倒在血泊中，两人已然毙命，一人重伤濒死。

而袭击者，赫然是五头通体碧绿、形如猎豹但头上生有独角的奇异妖兽，它们爪牙锋利，身上缠绕着淡淡的青色风旋，行动如电，正在撕咬尸体。

“是风刃豹，二阶妖兽，擅长风系法术，速度极快！”一名见多识广的散修低呼。

妖兽，通俗的说就是修炼有成的动物，至于这个“修炼有成”是怎么回事就不一定了。有些妖兽是先天的，即生下来就是妖兽；有些妖兽则可能是普通动物血脉觉醒或偶得机缘，后天成为的。内域把妖兽按内丹品质分成不同的阶位；一头二阶妖兽的战斗力，至少相当于一位玄脉境修士。

而这里有五头！

更让白言冰等人瞳孔收缩的是，在空地另一边的树冠上，隐约可见几道身影悄然隐匿，依稀看出其制服是紫阳宗和丹鼎阁的人！

他们冷眼看着下方散修被妖兽屠杀，毫无援手之意，甚至有人脸上露出残忍的快意。

树上的人气息收敛得良好，这群豹子似乎也没有抬头发现他们。

“他们八成是故意的！”吴刚咬牙道，“定是比我们更早传送至此，惊动了这群风刃豹，却故意引到这几人附近，自己躲起来看戏！”

这是赤裸裸的借刀杀人！守旧联盟果然一进来就开始下绊子。

似乎是察觉到白言冰等人到来弄出的动静，下方正在进食的风刃豹群突然抬头，碧绿的眼眸瞬间锁定了白言冰这支人数更多的队伍！

“吼！”

五头风刃豹舍弃了残破的尸体，化作五道青色流光，带着尖锐的破风声扑来；

尚未近身，便有十数道半月形的淡青色风刃抢先撕裂空气，覆盖而至！

那是风刃豹捕食之前释放的第一招，极为阴损，追求在近身搏斗之前就重创食物！

“结阵！三三制！”白言冰厉喝，声震山林。

这是进入秘境前，叶牧谦特意传授给所有参与弟子的简易战法，核心便是“三人一组，互为犄角，攻守一体，遥相呼应”，专门应对被围攻、偷袭或遭遇强敌时，避免落单被逐个击破。

实际上，这确实要拜叶牧谦前世记忆所赐，不少军队会使用这种战术队形；而更具体的实践，则完全交予这些弟子去训练与体悟了。

十六人瞬间动了起来！

白言冰作为唯一的见独境，修为最高，自成一组坐镇中心；沈瑶、陈千羽、吴刚自然成一组，作为主力。其余修士也与相熟的散修快速三三结阵。

面对覆盖而来的风刃，白言冰剑光一展，领域大放，圆融的剑意在身前划出一道弧线，将射向众人的大部分风刃引偏、卸力。

其他小组亦是或挡或闪，配合默契，虽略显仓促，却稳稳接下了第一波风刃袭击，无人受伤。

但风刃豹群速度极快，已扑至近前，獠牙利爪带着腥风。

“散！”白言冰再次下令。

五个三人小组如同花瓣般骤然散开，却又保持着彼此能够迅速支援的距离。

每组面对一头或分摊风刃豹的攻击，绝不贪功冒进。

白言冰自己则对上为首最强壮的一头，剑势沉稳，每一击都重若千钧，逼得那风刃豹连连后退，无法发挥速度优势。

其余修士亦利用对地形的熟悉和丰富的搏杀经验，很快也压制住了各自的对手。

树冠上，那几名守旧联盟弟子脸色微变。

他们没想到这群乌合之众反应如此迅速，战法如此古怪有效，竟然稳稳挡住了五头二阶风刃豹的突袭！

这“三三制”看似简单，却将个人战力有效整合，弥补了单个修士可能存在的短板，尤其擅长应对这种混战和偷袭。

眼看借刀杀人之计就要破产，其中一名丹鼎阁弟子眼中狠色一闪，悄悄取出并打开了一个小纸包。

“嗤——！”

一股极其细微、却充满躁动炽热气息的粉末随风飘散，混入战场。

正在与白言冰激斗的那头最强壮的风刃豹吸入粉末后，碧绿的眼眸瞬间爬满血丝，仰天发出一声痛苦与暴怒混杂的咆哮。

随即，周身青色风旋骤然变成混乱的赤青交错，体型似乎都膨胀了一圈，气息暴涨！它完全放弃了防御，不顾一切地扑向白言冰，速度力量暴增！

“小心！他们用了药！”陈千羽惊呼，她嗅觉敏锐，对药性敏感，立刻察觉不对。

白言冰压力陡增。但他临危不乱，见独领域全力展开。

周遭一切似乎变慢了些许，能清晰感知到疯豹扑击的轨迹与意图，以及疯豹身上肌肉力量的强弱分布。

他身形开始如柳絮般，随着疯豹带起的劲风飘退，剑尖却疯狂点向疯豹的关节处。

“沈瑶，吴兄，速战速决！”白言冰喝道。

沈瑶会意，身法陡然变得凌厉无比，刀光分化，让人眼花缭乱，速度快到让旁人以为她似乎是生出三头六臂、每一只手都以极高的频率向妖兽攻去。

瞬间，她对面的风刃豹身上留下数道深可见骨的伤口，且伤口开得颇大，特意选择了几条大动脉处放血；又戳烂了它的心与肺。

那豹子身上多处重要器官被废，又失血过多，登时就没了气。

吴刚也怒吼一声，与同伴合力，拼着受了一记爪击，将另一头风刃豹的头颅斩下。

那头疯豹却炸毛起来，久攻不下，愈发狂躁。

陈千羽觑准一个空隙，算明那头疯豹下一步的落点，随即打出一发微弱的元气扰动，竟让那片地面的草木忽然变成软垫。

疯豹本欲扑击，一脚踏上去，忽然爪下一软，平衡稍失，身体向侧方倾倒。

白言冰岂会错过这机会？

蓄势已久的一剑，凝聚了自身心力与元气，毫无花哨地直刺而出，正中疯豹因失衡而暴露出的咽喉弱点！

“噗嗤！”

随即手腕带着长剑向侧方一划！

剑刃透颈而过，疯豹头颅落地，咆哮戛然而止，庞大的身躯轰然倒地，赤青色的混乱风旋缓缓消散。

很快，最后两头风刃豹也被众人合力击杀。

战斗结束，众人微微喘息，除了吴刚和另外两名散修受了些轻伤，均无大碍。他们迅速收拾战场，把五只死豹剥了皮、挖了内丹，并将那重伤未死的散修同伴救起，由陈千羽进行急救。

白言冰这才倒出工夫抬头，冷冷地望向那处树冠。

那里，早已空空如也。

那几名守旧联盟弟子许是见事不可为，悄然遁走了。

“这只是开始。”白言冰擦去剑上血迹，目光投向森林深处那隐约传来的、更宏大的生机波动方向，“他们不会罢休，这秘境本身，恐怕也危机四伏。所有人，分成两队。我与吴兄各带一队，分散开来各自探索。”

吴刚点了点头。

虽然呆在一起确实更加安全，但分散开来也能够有更大的搜索范围，并在沿途获取更多资源。

休息大约一刻的时间后，两队分开，开始一头扎进森林各自探索。

\vspace{2em}

森林远比看上去的广阔和复杂。

那些青玉般的树木似乎隐隐构成某种天然的阵势，行走其间，方向感极易迷失。脚下厚厚的、散发着微光的苔藓和蕨类植物踩上去松软无声，却可能隐藏着致命的陷坑或毒虫。

空气中飘荡的芬芳里，偶尔会混入一丝甜腻得过分的气息，那往往是某些致幻灵花散发出来的，需得时刻运转心法保持灵台清明。

行进了约莫一个时辰，他们遭遇了又一次考验。那是一片看似平静的林间沼泽，水色清亮，生长着亭亭玉立的金色莲花，莲心凝结着露珠般的灵液，香气诱人。一名年轻些的云隐别院弟子面露喜色，下意识就要上前采集。

“别动！”白言冰和陈千羽几乎同时低喝。

话音未落，那弟子脚下看似坚实的苔藓地突然软化塌陷，数条漆黑如墨、快如闪电的藤蔓破水而出，直刺向他周身要害！藤蔓上密布细小的倒刺，闪烁着幽蓝光泽，显然含有剧毒。

那弟子骇然失色，慌忙催动灵力护体，挥剑格挡。但他仓促间的反应慢了半拍，眼看就要被藤蔓缠住。

嗤！嗤！嗤！

数道细微几乎难以察觉的破空声掠过。陈千羽出手如电，三根细长的钢针后发先至，精准地钉在了三条藤蔓的关节处。中针的藤蔓动作顿时一僵，速度骤减。

与此同时，白言冰并指如剑，凌空一划。一道圆融凝练的剑气匹练般扫过，巧妙地将剩余几条藤蔓的力道引偏，粘滞的剑势让它们如同陷入无形泥潭，挣扎着偏离了方向。

沈瑶的身影鬼魅般则出现在那名弟子侧后方，一把抓住其衣领向后急拉，险之又险地脱离了沼泽边缘。

整个过程发生在电光石火之间。那弟子落地后，脸色煞白，冷汗涔涔，几乎站不住，连忙向陈千羽和白言冰道谢。

“墨龙藤，植物，但有低等妖性。群居，擅伪装偷袭。”陈千羽仔细观察着缓缓缩回水中的藤蔓，冷静地分析。

白言冰目光扫过沼泽，注意到水下隐约有更多黑影蠕动，沉声道：“人命是第一位的。这东西不好拿，就不要了。记住，越是看起来无害诱人之物，越需警惕。”

这个小插曲让队伍更加谨慎。

\vspace{2em}

他们继续深入，期间又避开了几处气息诡异之地，也小心翼翼地采集了一些确认无害、且对修为有裨益的灵草灵石，在保证安全的情况下算是雁过拔毛。

约半日后，他们又听到了前方传来打斗声和妖兽的怒吼。

靠近一看，竟是三名守旧联盟的弟子正被七八头风刃豹围攻。

地上已经躺着两头豹子，但那三名弟子也颇为狼狈，衣衫破损，其中一人胳膊鲜血淋漓，显然受了伤，行动迟缓。

“又是风刃豹？”沈瑶低声道。

白言冰目光锐利地注意到战场边缘，几株被粗暴砍断的、结着朱红色果实的奇异灌木。灌木断口处流出乳白色汁液，散发出异常浓郁的甜香。

陈千羽翻了翻药典：“血朱果，对炼体有奇效。除了肉食，风刃豹特别喜欢吃这种东西。”

他们定是贪图果实，惊动了豹群。

此时，战况对守旧联盟弟子越发不利。

风刃豹攻击配合默契，速度如风；因为失血过多，那名受伤弟子的动作愈发迟缓。眼看又一头豹子要从侧后方扑向那名受伤者，带队的那名紫阳宗弟子虽奋力挥出火法逼退正面之敌，却已救援不及。

是袖手旁观，看着竞争对手减员？还是……

白言冰眼神一凝，瞬间有了决断。

“沈瑶，左翼扰袭！千羽，准备包扎！你们俩，跟我上，救人！”

话音未落，他已身化流光掠出，巧妙切入侧翼，剑指虚点，数道凝实的剑气精准地射向几头豹的落足点，打乱了它们围攻的节奏。

沈瑶如一道黑色轻烟，倏忽出现在战场另一侧，手中短刃带起森寒光弧，以刁钻的角度袭向风刃豹的眼睛、关节等脆弱处，迫使它们分心防御。

另外两名云隐别院弟子紧随白言冰，结成简单剑阵，挡住了一侧扑来的两头豹。

突如其来的援手让三名守旧联盟弟子压力一松。带队者愣了一下，神色复杂地看了白言冰一眼，猛一咬牙，趁机催动更强术法，配合反击。

混乱中，陈千羽悄无声息地靠近，银针连闪，先为那名重伤弟子封住血脉，弹入一枚清香药丸暂时压制伤势，又挥手洒出一片淡绿色药粉。药粉弥漫开，带着清心解热的气息，让暴躁的豹群产生了瞬间的迟疑。

趁此机会，白言冰低喝一声：“退！”

新生势力的人加上被救下的守旧联盟三人，迅速脱离战圈，向后疾退。豹群追出一段距离，发出不甘的咆哮，但似乎又顾忌领地，于是转身离去了。

脱离危险后，那三名守旧联盟弟子，尤其是带队的紫阳宗弟子，面色一阵青白。

他看了看自己受伤的同伴，又看了看神色平静的白言冰等人，抱了抱拳，干涩地道：“在下韩铁，多谢……援手之情，我等记下了。”

说完，不再多言，扔下几枚血朱果，就迅速带着同伴离去，背影有些仓惶。

他显然是不知道，白言冰一行人刚刚进入秘境时也遇到了相似的风险；他更不知道，他的同宗面对几乎完全一样的情景做了些什么。

他们似乎是没脸在此多做停留，也有可能是害怕被人告发“通敌”。

白言冰见他们竟然扔下几枚果实，有些愕然，但最终还是给它们收了起来。

实际上所有的冲突都是这样的：大多数人其实是沉默的、中立的，他们虽然心中可能并不厌恶甚至同情云隐别院，但大势所迫之下，为了明哲保身，也不得不被裹挟。

“为何救他们？”一名云隐别院弟子不解地问，“入口处他们可是下了黑手。”

白言冰淡然道：“见死不救，与那些只知内耗之徒何异？我等所求之道，非仅力量强弱，亦有心境高低。”

陈千羽微微颔首，看向白言冰的目光中带着赞同。沈瑶则若有所思。

\vspace{2em}

收拾心情，他们继续前行。

随着不断深入，周围的古木愈发高大，生机愈发磅礴，甚至能看到一些外界早已绝迹的灵草珍木。他们也遭遇了其他一些零散的探险者，有散修，也有个别天师府弟子，大多行色匆匆，彼此戒备，偶尔交换一些无关紧要的情报。

时间在探索、警惕、偶尔的小规模冲突或互助中流逝。

第三日，他们穿过一片弥漫着淡金色雾气的奇异林地时，终于遭遇了进入秘境后第一次有预谋的、来自守旧联盟修士的针对性袭击。

袭击发生在林间一处狭窄的隘口。当他们半数人通过时，两侧看似普通的、爬满藤蔓的岩壁突然爆开，预先埋设好的爆裂符箓和禁锢阵法瞬间启动！

与此同时，上方树冠之中，隐匿许久的五名守旧联盟弟子骤然发难，炽热的火球、锋锐的金芒劈头盖脸砸下，目标明确——直指被暂时分隔开的队伍中段，尤其是陈千羽和另外两名修为稍弱的弟子！

“敌袭！结阵！”白言冰厉喝，身处队尾的他反应最快，身形一晃已拦在攻击最密集处前方。

长剑出鞘，划出一道浑圆的剑幕，如中流砥柱，将大半攻击挡下。

领域圆融稳固，虽被轰击得波纹荡漾，却坚不可摧。

白言冰自信，除非周文彬亲自出手，否则剩下的这些弟子无非土鸡瓦狗，连他的防御都破不开来！

前方的沈瑶和另一名弟子亦瞬间反击，刀光剑影迎向从岩壁后冲出的两名敌人。

但对方算计精准，仍有两道刁钻的攻击绕过白言冰的防御，袭向陈千羽。陈千羽临危不乱，素手轻扬，一面由无数细密银针虚影构成的元气护盾展开，同时袖中飞出数点寒星，直射施法者的要害，攻守兼备，逼得对方不得不回防。

战斗在刹那间爆发，狭窄的隘口内灵力激荡，光芒乱闪。守旧联盟弟子人数虽稍少，但占了先手埋伏的优势，且配合默契，显然是早就盯上了他们，特意在此设伏。

“周师兄有令，重点关照云隐别院和那姓白的！尤其是那个医修，先废了她！”一名紫阳宗弟子狞笑着，双手连拍，一条烈焰凝成的巨蟒昂首扑向白言冰，试图将他缠住。

白言冰眼神冰冷：别人怎么动他倒先不论，但陈千羽可是他的逆鳞。

随即，手中长剑似乎狂吼起来，原本圆融守御的剑势陡然一变。

弥焰斩！

雷火剑气大盛，剑势如焰天火雨，层层叠叠，与那火蟒硬拼。

黑紫色的剑气与火蟒甫一接触，后者就被这至刚至强的雷火之气拦腰纵向切成了两半，半只火蟒反而扫向了旁边一名试图偷袭沈瑶的家伙。那修士猝不及防，慌忙闪避，阵型顿时出现一丝紊乱。

“就是现在！”白言冰低喝。

始终游走策应的沈瑶如同鬼魅般从那名慌乱弟子的影子里钻出，短刃带着一抹凄艳的弧光，直取其肋下。

刀光深入血肉！

随即，沈瑶本能地扭转了一下自己的手腕，短刀自然在家伙的胸膛中豁开了一个更大的口子，肌肉绞碎，鲜血直流。

这或许是往年游侠时光遗留的习惯与本能：要刺杀，就一击致命。

拔刀，刀柄带着元气猛击那人天冲。

那重伤弟子口鼻出血，倒地不起，眼见难活。

遭遇战瞬间变成了混战，但白言冰这边人数占优，且经过几日磨合，配合渐佳，很快稳住阵脚，并开始反压。

指挥埋伏的那名紫阳宗弟子见事不可为，己方已有一人重伤，再缠斗下去恐有折损，当即发出尖啸：“撤！”

其余守旧联盟弟子毫不恋战，纷纷掷出烟雾符箓，借助地形和林木掩护，拖着七窍流血的同门迅速遁走。

临走前那阴狠的眼神，显然不会善罢甘休。

白言冰并未深追，秘境之中，贸然追击极易落入新的陷阱。他迅速查看己方情况，除了两人受了些轻伤，并无大碍，陈千羽已及时为他们处理。

“他们开始有组织地针对我们了。”沈瑶擦拭着短刃上的鲜血，冷声道。

“意料之中。”白言冰收剑入鞘，看着敌人消失的方向，“那所谓‘周师兄’既已视我等为心腹之患，自然不会让我们顺利探索。”

\vspace{2em}

又过了两日，他们在一处位于山涧旁的灵泉边休整时，遇到了吴刚带领的、已扩充到八人的万仙盟小队。

双方都有些风尘仆仆，但眼神精悍，显然收获与经历都不少。

吴刚主动走了过来，咧了咧嘴：“白兄弟，看来你们也没少被那帮孙子找麻烦。”

白言冰点头：“遭遇过一次埋伏。吴兄那边如何？”

“打了两场，吃了点小亏，也让他们见了点血。”吴刚眼中闪过一丝狠色，“散修别的不行，记仇和找机会下手最在行。他们现在也不敢小看我们了。”

他顿了顿，压低声音，“我们的人发现了一些痕迹，守旧联盟那帮人，似乎在有意无意地把一些难缠的妖兽或危险区域，往可能的路径上引，想借刀杀人。另外，对于那些资源，他们似乎只取最好的一部分，剩下的多半被他们毁掉了。”

白言冰心中一动：“大致方向？”

吴刚摇了摇头：“方向嘛，和我们感应到的那生机最浓的地方，大体一致。而且，根据我们零星得到的情报，其他几股咱们这边的人，还有那些中立的、只想捞好处的天师府弟子，大多也在朝那个方向靠拢。看来重头戏都在那儿。”

这个情报至关重要。双方简要交流了一下各自探知的部分地形和危险点，便再次合成一队。白言冰率领队伍，朝着那冥冥中感应到的、生机与道韵汇聚的核心方向，继续前进。

沿途，他们看到了更多战斗痕迹，有修士与妖兽的，也有修士之间的。秘境资源的争夺，已经逐渐白热化。

他们甚至在一处山谷中，发现了三具尸体，皆是散修打扮，财物被搜刮一空，致命伤显然并非妖兽所致。

秘境的柔和青光，似乎也掩不住日渐浓厚的血腥与肃杀。

\vspace{2em}

第七日，他们翻越一片矮小的山岭后，眼前景象骤然一变。

不再是无穷无尽的乔木树林，而是一片开阔的、望不到边际的奇异原野。

极目远眺，原野的尽头，隐约有连绵的虚影，似殿宇飞檐，又似山峦轮廓，被浓郁的青色霞光笼罩，看不真切。

但那股浩瀚、古老、充满生机的源头波动，正清晰无比地从那个方向传来，如同心脏的搏动，吸引着所有人的灵魂。

同时，他们也看到了其他身影。在原野的不同方位，影影绰绰，至少有数十人也在朝着相同方向前进。

其中一批人数量较多，服饰相对统一，气息灼热或锋锐，正是以紫阳宗圣子周文彬为首的守旧联盟队伍，他们似乎刚刚汇合不久，声势不小。

另一边，则是零散分布的小队或独行者，有万仙盟的人，有天师府弟子，也有其他幸存修士。

双方隔着很远的距离，彼此都能看见，气氛瞬间变得微妙而紧张起来，但谁也没有轻易挑起冲突。所有人的目标，显然都是原野尽头，那片青色霞光笼罩之地。

白言冰知道，最初的探索和淘汰阶段，已经接近尾声。真正的核心区域，就在前方。而最终的对峙与争夺，即将在那片青霞之下展开。

“收敛气息，保持距离，我们走。”

他低声下令，带领队伍，向着秘境最终的核心悄然进发。

\vspace{2em}

秘境核心，一片无比广阔、被永恒柔和青光笼罩的古老广场，豁然呈现于天地之间。

那青光仿佛源自广场本身与周遭的空间，均匀、明亮却不刺眼，带着温润的滋养之意，照在身上，连多日奔波的疲惫都似乎被抚平了些许。

而广场的中央，那座巍峨矗立的殿宇，攫取了所有人的目光与呼吸。

它不知以何种奇异的美玉筑成，通体呈现出一种深邃而温润的青色，仿佛将一片无云的晴空或是宁静的深海凝结成了实体。

岁月是无情的雕刻师，在它身上留下了不可磨灭的斑驳痕迹：墙壁有多处塌陷，断裂的梁柱如同巨兽的骨骼般支棱出来，精美的飞檐斗拱残破不全；无数不知名的藤蔓，叶片如翡翠，缠绕着那些残垣断壁；各色奇花异草，散发着静谧的微光，从裂缝中顽强探出，恣意绽放。

破败与生机，沧桑与神圣，在此刻达成了惊人的和谐与统一，共同诉说着一段淹没在时光长河中的辉煌过往。

殿门紧闭，高达十丈的青铜门扉上覆盖着厚厚的铜绿，但门楣上方，那古老鸾鸟翱翔的浮雕依然清晰可辨。鸾鸟的线条流畅而充满力量，羽翼舒展，似乎下一刻就要破石而出，直上青冥。

即便只是浮雕，也散发出一股令人心悸的、磅礴如海又精纯如露的生机之气，更蕴含着一种悠远深邃、难以言喻的大道韵律，让注视者不由自主地心生敬畏与渴望。

然而，一层肉眼可见的、半透明的光罩，如同一个倒扣的青金色巨碗，将整座殿宇连同其前方的大片广场牢牢笼罩。

光罩之上，无数复杂玄奥的青金色符文如活物般流转不息，时而汇聚成鸾鸟虚影，时而散作漫天光雨。

最令人心惊的是，光罩之上生机与毁灭两种截然相反的力量奇异地交织、平衡在一起，仿佛在平静的海面下隐藏着足以吞噬一切的暗流，警告着任何妄图强行闯入的不速之客。

\vspace{2em}

几位先到一步的天师府弟子，模样有些灰头土脸，正站在光罩外不远处，脸上带着无奈与不甘。他们仅仅是早到了半刻，尝试了各种手段试图直接破开护罩，甚至联手轰击过，但那光罩纹丝不动，反震之力还差点让他们吃了亏，此刻只能望洋兴叹。

所有人的视线，最终都聚焦在殿门正前方，光罩之内、广场之上，那处唯一异样的区域——一个直径约三丈的完美圆形。圆形区域内的地面，镌刻的阵纹比之外围光罩上的还要繁复精密百倍，线条交错如星空轨迹，节点闪烁如晨星，氤氲着浓郁的、几乎要液化的青色光晕。

就在这时，广场两侧边缘，几乎同时涌入了两股泾渭分明的人流。

左侧，以紫阳宗圣子周文彬为首，守旧联盟的弟子们鱼贯而出。

他们一共十几个人，虽然衣衫大多沾有尘土、灵草汁液甚至干涸的血迹，不少人的衣袍还有法术或利爪造成的破损，但整体气势依旧强盛。紫阳宗弟子周身隐隐有灼热气息升腾，无极剑派弟子则锋芒大盛，丹鼎阁弟子药香萦绕。

他们的眼神，在疲惫之下，是毫不掩饰的灼热与占有欲，如同盯上猎物的群狼，死死锁定了青光中的古老殿宇。

自进入秘境以来，他们秉持着大派作风，遇见资源与宝物，只取最精华的一部分；路线上，选择绕路避战，虽经历风险，但整体保存了更多实力。

右侧，白言冰、陈千羽、沈瑶带领的云隐别院及新生势力联盟的修士们也相继现身。他们的人数稍多一些，但状态看上去更为复杂。不少人身上带着明显的战斗痕迹，包扎的布条下隐隐渗出血色，气息也更为起伏不定，显然经历了更多直面妖兽、破解禁制的硬仗。

他们的眼神同样锐利，却多了几分历经磨砺的沉凝与警惕，如同在荒野中披荆斩棘终于找到水源的旅人，既有渴望，也深知周遭可能潜伏的危险。

\vspace{2em}

双方人马在广场边缘停下，隔着近百丈的距离遥遥对峙。空气中原本充盈的平和生机灵气，仿佛瞬间被无形的张力所凝固，充满了山雨欲来的压抑。

没有人说话，只有轻微的喘息声、衣物摩擦声，以及那仿佛响在每个人心头的、来自殿宇的磅礴道韵波动。青鸾传承的诱惑，近在咫尺，唾手可得，却又被那层薄薄的光罩无情阻隔。

“看来，这便是最后的考验了。”一个带着傲然与灼热气息的声音打破了沉寂。

紫阳宗圣子周文彬排众而出，他向前踱了几步，暗红色的锦袍下摆虽有一道被划开的口子，却无损其逼人的气势。他面容俊朗，此刻却因兴奋和野心而显得有些凌厉；目光如炬，扫过那圆形分阵和其后巍峨的殿宇，嘴角勾起一抹志在必得的弧度。

“此阵玄奥，非蛮力可破。谁能率先破解，便占得先机！”他声调抬高，似在宣告，又似在挑衅，说话间，却不着痕迹地向身后的同门递了一个隐晦的眼色。

看起来他想率先试图破阵。如果别人胆敢抢先，就会被身后众人群起而攻之。

“周圣子所言极是，那便请吧。”

白言冰的声音平和响起。

周文彬瞥了那带头剑修一眼，寻思这就是白言冰了，于是冷哼一声，不再多言。

他深吸一口气，胸膛微微起伏，周身瞬间腾起赤红色的精纯灵力，如同包裹在一层燃烧的琉璃焰火之中，将其衬得宛如火神临世。

在双方数十道目光的聚焦下，他迈开步伐，沉稳而坚定地踏入了那直径三丈的圆形分阵之中。

就在他双足完全落定在繁复阵纹上的刹那——

“嗡！”

阵中沉寂的青色光晕骤然爆发，化作亿万颗细密柔和的光点，如同春日清晨最晶莹的露珠，又似夏夜飞舞的流萤，温柔却无可抗拒地将周文彬的身形彻底淹没。外界看去，只见一团浓郁的人形青光，其内人影猛地一僵。

下一刻，周文彬那张原本写满傲然与志在必得的脸，在青光掩映下骤然扭曲！

极度惊愕、难以言喻的恐惧、以及一丝被人窥破最深处秘密的慌乱，如同打翻的颜料盘，瞬间取代了所有冷静与算计。他的身体不受控制地微微颤抖起来，额角、鼻翼以肉眼可见的速度渗出细密的冷汗，在青光照耀下反射着冰冷的光泽。

他仿佛看到了什么景象，瞳孔剧烈收缩。

“不完全是锁钥，竟然上面还套了层幻阵！”饶是最精通阵法的沈瑶也不禁惊叹这外挂幻阵的精妙。她甚至一开始并没看出来还有第二重阵法！

只见阵中的周文彬开始有了动作，时而面露狰狞凶狠之色，抬手向前虚抓虚击，仿佛在与看不见的敌人搏杀；时而又仓惶失措地连连倒退，脚步虚浮，口中无意识地喃喃自语，声音断续却清晰地传了出来：

“不……不是我……那是他……是他自己技不如人！对，是他技不如人！”

他像是在向某个无形的审判者激烈地争辩，又像是在拼命说服自己，额上青筋都已暴起。

这般挣扎持续了不过十数息，对周文彬而言却漫长得如同百年。终于，他猛地闷哼一声，像是心口被重重一击，脸色“唰”地变得惨白如纸，再无半分血色。

他竟主动踉跄着向阵外倒退，脚步虚浮，差点跌倒，靠着身后同门及时上前搀扶才稳住身形。他气息紊乱不堪，胸口剧烈起伏，再抬头望向那圆形分阵时，眼神里已充满了后怕和不敢相信。

守旧联盟众人围拢上来，脸色都变得异常难看，几个原本摩拳擦掌、准备紧随圣子尝试的核心弟子，此刻也面面相觑，眼神闪烁，默默地向后退了半步。

谁心里没点不愿示人、甚至自己都不愿深想的隐秘？这幻阵竟能直指本心，挖掘最深处的愧悔与阴暗，以此来影响破阵者，属实精妙。

“看来此阵，专照人心鬼蜮。不过周圣子良心未泯，可喜可贺。”一个清冷中带着淡淡讽刺的女声响起。

沈瑶越众而出，她黑衣如墨，衬得肌肤欺霜赛雪，原本清丽却稍显冷硬的脸庞上，此刻是一片近乎淡漠的平静。

“师妹小心。”陈千羽看着她，语气温和却郑重地叮嘱。

沈瑶对她微微点了点头，眼神交汇间自有默契。她没有丝毫犹豫，步履平稳，如同走向一处寻常回廊，径直走入了那令周文彬狼狈不堪的青色光点之中。

幻境瞬间展开，毫不留情。

熟悉的、几乎烙印在灵魂深处的灼热痛楚排山倒海般袭来，仿佛时光倒流，将她猛地拉回十年前那个绝望的午后——亡命徒狰狞的脸，激射而来的炽热火符，瞬间包裹头脸的恐怖高温，皮肤焦灼产生的刺鼻气味，还有随之而来的、无边无际的剧痛与黑暗。

火焰在眼前幻化跳动，热浪灼烧着感知，那曾日夜折磨她的阴影如同冰冷的潮水，试图再次将她淹没、吞噬。

若是从前那个未曾真正放下的沈瑶，此刻或许会被激发出更强烈的怨恨与斗狠之心，与幻境殊死搏斗。

然而，此刻阵中的沈瑶，只是在那灼热火光扑面而来时，轻轻地、极其自然地闭上了眼睛。长而微翘的睫毛在青光中投下淡淡的阴影，复又睁开。

那双曾盈满痛苦与戒备的眼眸里，此刻映照着幻象的火焰，却清澈得如同雨后的寒潭。

没有恐惧，没有愤怒，甚至连一丝波澜都难以寻觅。

她缓缓伸出手，素白的手指探向那幻境中焚毁她容颜的火焰，在即将触及那虚幻炙热的刹那，停了下来。她就那样静静地看着，看着幻象中那个痛苦蜷缩、面目狰狞、全力嘶吼的少女，眼神里掠过一丝极淡的、近乎悲悯的情绪，最终化为一片历经千帆后的平静与释然。

她红唇轻启，声音不高，却奇异地穿透了阵法的青光隔绝，清晰地传入了广场上每一个人的耳中，带着一种斩断枷锁的轻松与笃定：“那只是……过去的我罢了。”

随即，沈瑶坦然坐下，将那幻阵特意演绎出的残酷画面搁置一旁，开始专心破阵。

作为锁钥的阵法倒是不难，以沈瑶的阵法造诣，很快就轻松破解开来。

“嗡——！”

一声清越悠扬、仿佛穿越了万古时空的鸾鸟鸣响，蓦然在每个人心底响起！

笼罩巍峨殿宇的庞大青金色光罩应声发生了剧烈变化，符文疯狂流转，光华明灭不定，随即，在紧闭的青铜殿门正前方，那光滑坚固的光罩如同被无形之手从中间分开的瀑布水幕，流畅而无声地向两侧滑退、消散，露出了其后古朴厚重、布满铜绿的真实殿门。

防御阵，开了！

“进去了！”不知是谁按捺不住激动，嘶哑地喊了一声。

这一声如同投入滚油的火星，瞬间点燃了所有积累的贪婪、焦急与竞争之心。

刹那间，对峙的平静被彻底撕碎，众多修士再也顾不得阵型与矜持，一个个眼泛红光，将身法催动到极致，化作数十道颜色各异的流光，争先恐后、你推我挤地鱼贯而入，疯狂冲向那敞开的、仿佛通往无上机缘的殿门。

似乎晚上一步，那传说中的青鸾传承就会被他人捷足先登。

作为破阵的最大功臣，沈瑶反而被这汹涌的、失去理智的人流稍稍挤到了后面。白言冰与陈千羽一左一右护在她身旁，三人随着澎湃的人潮，踏上了殿门前的台阶。

就在陈千羽即将跨过那道古朴青铜门槛的刹那，一声哀鸣，毫无预兆地在她心海最深处响起。

那声音极其微弱，细若游丝，仿佛下一刻就要断绝，却蕴含着如此浓郁的痛苦、无助，以及一种对生命、对这片天地深深的眷恋与不舍。它似乎是直接响彻在灵魂深处，像一根冰冷的针，直接刺中了陈千羽心中最柔软的部分。

她脚步不由得一顿，下意识地转过头，目光越过躁动拥挤的人腿，望向殿门旁侧，一处被高大石雕残骸阴影覆盖的角落。

只见那里，靠近冰冷石基的地面上，躺着一只巴掌大小的鸟儿。

它羽毛本应是华丽的青色，此刻却沾满尘土，凌乱暗淡，失去了所有光泽。

一只翅膀以极不自然的角度弯折着，软软地耷拉在地上，隐约可见红肿。它气息奄奄，小小的胸膛微弱起伏，那双本该灵动剔透的琉璃色眼眸，此刻蒙着一层灰败的阴翳，正哀伤地、近乎绝望地望着眼前蜂拥闯入殿内、对它视而不见的人群。

守旧联盟的弟子们冲在最前，眼里只有殿内的黑暗与可能的珍宝，对脚边这微不足道的小生命连瞥一眼都嫌浪费时间，甚至有人嫌它略略碍路，随意一脚踢开。

青色小鸟孱弱的身躯翻滚了几下，撞在冰冷的石头上，发出的哀鸣更弱了几分，几不可闻。

许多云隐别院的弟子、万仙盟的修士、以及中立的天师府弟子，大多也只是匆匆一瞥，便随着人流涌入殿内。

机缘当前，一只将死的、寻常可见的鸟儿，又算得了什么？

然而，陈千羽的心，却像是被那只小鸟眼中浓得化不开的哀伤，以及那声直达灵魂的哀鸣，给紧紧攥住了，闷闷地发疼。

她是医者，师承云隐别院，见过更惨烈可怖的伤痛，也对生命本身怀有一种根植于道的、本能的悲悯与尊重。这青鸾秘境处处生机盎然，道韵沛然，一只如此模样的小鸟，偏偏出现在这核心殿宇的门口，绝不可能只是巧合。

更重要的是，那哀鸣中传递出的情绪，纯净而强烈，让她无法像其他人那样硬起心肠，转身离开。

“千羽？”白言冰发现她停下，顺着她的目光也看到了阴影中的小鸟，眉宇微微一蹙，显然也察觉到了异常，但并未出声催促，只是安静地等待她的决定。

陈千羽咬了咬下唇，贝齿在唇上留下浅浅的印子。

内心天人交战仅一瞬，那源于医者本心的悲悯终究占据了上风。

她不再犹豫，快步逆着稀疏下来的人流走过去，在少数几个尚未进殿、投来诧异不解目光的弟子注视下，小心翼翼地蹲下身，伸出双手，用最轻柔的力道，如同捧起一滴晨露般，将那只受伤的青色小鸟捧在了温热的掌心。

小鸟在她掌心微微颤抖着，羽毛拂过皮肤带来细微的痒意。它似乎连挣扎的力气都没有了，只是费力地抬起小小的头颅，那双琉璃般的眼眸望向陈千羽。

奇异地，那其中盈满的哀伤，竟在这一望之下稍微减退了些许，取而代之的是一丝微弱却清晰的依赖与祈求，它甚至试图将受伤的脑袋靠近她温热的指腹。

“你们先进去，我马上来。”陈千羽抬头，对白言冰和沈瑶说道，声音轻柔却坚定。

同时，她已运转起体内温润平和的元气，那气息带着她特有的、促进愈合的灵性。

一层淡淡的暖光，轻轻包裹住小鸟折翅的伤处，先稳住它不断流逝的生机与加剧的痛苦。

白言冰与沈瑶对视一眼，彼此眼中都掠过一丝了然。他们深知陈千羽外柔内刚的性情与那份对生命的执着，也相信她的直觉与判断。

“小心。”两人齐声道，不再多言，转身踏入了那深邃的殿门之中，身影迅速被黑暗吞没。

陈千羽怀抱这只奇异的青鸟，成为最后一位踏入大殿的修士。

当她迈过那道高高的青铜门槛时，身后沉重的殿门，仿佛有所感应般，发出一声低沉悠长的叹息，无声无息地缓缓合拢，将外界那片柔和的青光与广阔的广场彻底隔绝在外，只剩下门缝消失前最后一缕微光，映亮她沉静却坚定的侧脸。

\vspace{2em}

殿门之内是一个无比广阔、远超外部所见规模的空间：无数条通道、回廊、拱门纵横交错，向四面八方延伸，构成了一座极其复杂的巨大迷宫。

墙壁的材质不知是何物，触手温润，通体散发着均匀柔和的青光照亮前路，却又不显得刺眼。墙壁之上，镌刻满了巨大的、连贯的古老壁画。壁画的内容恢弘壮丽、栩栩如生，充满了灵动盎然的大道韵律，讲述着青鸾司掌生机、平衡万物的古老神职。

空气中弥漫着的，是浓郁到几乎化为实质的生机灵气，比之外界广场又不知强盛了多少倍。每一次呼吸，都仿佛饮下琼浆玉液，四肢百骸说不出的舒坦。

然而，与此同时，神识在此地如同陷入浓稠的泥沼，勉强能离体数丈便再难延伸；而那错综复杂的通道和看似有规律的壁画，更是在不断混淆、误导着人的方向感，使人难辨东西。

先进来的修士们早已被这巨大的迷宫和躁动的贪欲所分散，如同水滴汇入沙地，消失在各条通道深处。远处，不同方向不时传来惊喜的呼喊声、激烈的打斗声、以及灵力爆鸣或古老禁制被触动后发出的沉闷爆炸声，在这空旷的迷宫中回荡，更添几分混乱与危险的气息。

白言冰、沈瑶很快与吴刚等万仙盟的人，以及云隐别院其余人等在一处宽敞的十字回廊汇合。

众人苦迷路久矣，于是经过商议，选定一个方向，采用最稳妥的方式：每走过一段回廊，遇到门户便开门查探；每探索完一个房间，便在门上留下一个简明的记号，然后便转向下一个目标。

房间内的情形各不相同，有的空空如也，有的暗藏机括陷阱，猝然发动，需得小心应对；但也有的房间，里面存放着一些蒙尘的财富——或是堆积在墙角、灵气盎然的各色灵石；或是玉瓶中药力未曾完全散失的旧日丹药；或是样式古朴、灵光暗淡却依然可用的法器。

不久，周文彬带领的守旧联盟队伍也发现了这种方法的有效性，立刻有样学样，开始分散成数支小队，在迷宫中进行地毯式的搜索，同样在探索过的门户上留下标记。

然而，随着时间推移，无论是新生势力还是守旧联盟的修士，都陆续开始感觉到不对劲。

许多通道看似笔直通向更核心的幽深之处，满怀希望地走上一段后，却往往在拐过一个弯或穿过一道拱门后，愕然发现自己又回到了之前经过的、有着熟悉记号的岔路口，甚至干脆是冰冷的、刻画着壁画的死胡同。

墙上的壁画内容，初看时宏伟壮丽，看久了却觉得它们似乎在缓慢变化，又仿佛亘古如此，总在不经意间将人的思绪引向某种无始无终的循环。

不少弟子开始心浮气躁，如同被困在琉璃罩中的无头苍蝇。

明明宝藏可能就在附近，却不得其门而入，只能徒劳地在相似的廊道间打转，消耗着体力和耐心。

\vspace{2em}

唯独陈千羽，怀抱那只体温正在缓慢回升的青色小鸟，独自走在光影流转的迷宫之中，所经历的感受却与所有人都截然不同。

那干扰神识、混淆方向的无形力量，对她似乎影响甚微。

她心中有一种非常微弱、却真实不虚的模糊意识——似乎是一种若有若无的共鸣与牵引。

似乎……掌心那蜷缩的、依赖着她元气存活的小小生命，正用它微弱的气息，与她脚下这座古老殿宇的某处核心，产生着难以言喻的联系，并在冥冥中为她指明前路。

陈千羽没有抗拒这份感觉。

她自然而然地迈开脚步，不再刻意去记忆复杂的路径或辨识壁画的不同。

她走过刻画青鸾哺育幼雏的回廊，壁画上的神鸟眼眸似乎温和地望向她，轻轻颔首；她穿过一道拱门，门后本应是数条岔路，此刻却唯有其中一条，青光明亮而柔和，道路笔直地向前延伸，仿佛专为她而呈现。

怀中的青鸟，身体不再冰冷颤抖，折翅处在她持续不断输出的温润元气滋养下，那淤肿竟在缓缓消退，破损的皮肉下有微弱却顽强的生机在萌动、凝聚，试图接续断骨。

不知不觉间，她与白言冰、沈瑶等人走散了，也仿佛超脱了这座困扰所有修士的迷宫的固有规则。

她不再属于那些焦躁寻找、重复打转的人群。

\vspace{2em}

周围的景象开始发生微妙而持续的变化。

回廊愈发显得古老、静谧，墙壁散发的青光愈加纯粹，壁画上的青鸾形象愈发鲜活灵动，羽翼上的每一丝纹路都清晰可见，眼眸顾盼生辉，仿佛随时会发出一声清鸣，振翅从那墙壁之上飞腾而出，翱翔于宫殿穹顶。

最终，在迷宫的至深之处，她踏过最后一级泛着微光的台阶，眼前已无岔路，只有一扇孤零零的门户。

这扇门与其他任何房间的石门、木门都截然不同，它通体由凝实的青色光华构成，光晕流转，门扉上隐约有更复杂古老的纹路明灭，静静地矗立在一面刻画着青鸾归巢、敛翼安息的巨大壁画之前。

似是感应到她的到来，尤其是她怀中那只青鸟微弱却清越的一声低鸣，青色光门微微荡漾起来。

随后，它无声地、温柔地向内洞开。

门后流泻出的，是一片温暖、浩瀚、充满无尽生命喜悦的青色光辉。

\vspace{2em}

门后，并非她或任何人想象中的、堆满古老典籍或闪耀宝物的藏宝室，而是一座极为空旷、穹顶高远的圆形殿宇。

殿内空无一物，唯有地面、墙壁、穹顶都流淌着那种温润的青色光辉，将每一寸空间都映照得澄澈而圣洁。这里寂静得能听到自己血液流动的声音，却又仿佛充满了无声的宏大乐章。

而在殿宇的最中央，悬浮着一团无法用言语形容的青色光源。

它是“生机”这个概念最纯粹、最本源的显化，浓缩凝聚到了近乎实质的境地。它如同一颗宏伟的青色心脏，缓缓地、有力地搏动着：每一次舒张，便有无形无质却沛然莫御的生机道韵如水波般荡漾开来，充盈整个殿宇；每一次收缩，又仿佛将万物的喜悦与活力吸纳归一。

站在这团光辉前，陈千羽感到自己浑身上下的每一寸都在欢欣雀跃，多年行医积累的疲惫、先前跋涉时的紧张，乃至灵魂深处细微的尘垢，都被这纯粹的光辉洗涤、抚平。这不是力量的压迫，而是生命本源对同源者的呼唤与抚慰。

就在这时，她掌心传来轻微的悸动。

那只一直安静蜷缩、依赖她元气存活的青色小鸟，忽然抬起了头。它那双原本蒙着阴翳的琉璃色眼眸，此刻变得清澈无比，倒映着中央那团青色光辉，充满了孺慕、欢欣，以及一种归家的释然。

它发出一声清越悦耳、宛如冰玉相击的鸣叫，声音带着穿透灵魂的纯净喜悦。

挣扎着，用那双尚未完全恢复的翅膀，在陈千羽掌心轻轻一蹬，身体歪歪斜斜地飞了起来。仿佛初学飞翔的雏鸟，却又带着义无反顾的决绝，径直投向了殿宇中央那团最为炽盛纯净的青光。

小鸟的身形在接触光团的刹那，便如同水滴融入大海，悄无声息地消失了。

紧接着，那团青色光辉骤然膨胀、盛放！

无穷无尽的光淹没了陈千羽的视野，将她温柔而彻底地包裹。

一个古老、沧桑、却依旧温和慈爱的意念，跨越了漫长到难以想象的岁月长河，如同初春暖阳下解冻的涓涓溪流，自然而然地、涓滴不漏地流入陈千羽敞开心扉的灵台识海。

青鸾残念的声音，充满了时光沉淀的智慧与淡淡的忧伤。

“后来的生灵啊……你终于来了。”那意念的开端，像一声悠长的叹息，“我是此地主人最后残留的一点真灵印记，依托这秘境核心的生机本源而存，等候了太久太久。”

陈千羽心神震撼，在意识中回应：“您是……青鸾前辈？”

“称谓并不重要。重要的是，你看到了它，并选择了看到了它。”

那意念中流露出赞许的情绪，“当贪婪的目光掠过它，投向那扇被认为藏有至宝的门户时，唯有你，为一只看似平凡的受伤小鸟驻足。这种源于本心的、对生命本身的悲悯与尊重，正是我所执掌之道的基石。”

陈千羽眼前仿佛浮现出殿门外那混乱的一幕：匆忙践踏的脚步，冷漠无视的目光，还有那被随意踢开、瑟瑟发抖的渺小生命。她心中微微一疼。

“青鸾主生机，亦掌平衡。”那意念继续流淌，向她展开一幅浩瀚的图景，“生机并非一味滋生繁荣，平衡亦非僵死不变。春日勃发，秋日凋零，皆是循环；滋养万物，亦需抑制过盛，方为长久。”

无数玄奥的意象随之涌入陈千羽的脑海：种子破土时蕴含的爆破之力与柔韧之性，参天古木将阳光雨露转化为滋养的精密过程，伤病躯体中破损生机与修复力量的微妙拉锯，乃至一方天地间灵气的盈亏循环……这些不再是抽象的概念，而是一个个鲜活具体的道之显化。

在此道理之中，自然包含了无数早已失传于世的丹方配伍、药性医理、精粹炼药手法。

陈千羽感到自己熟知的那些药理知识，如同散落的珍珠被一根青色的丝线串联起来，焕发出全新的、更深层次的光芒。

那意念的声音渐渐变得空远，仿佛力量正在消散：“传承与你，并非要你成为另一个青鸾。你的道途与他们乃至我主都不同，没有那种强夺的本性。我只希望你能持此道，行你之路，在你所处的时代，护持那一点对生命的敬畏与平衡之心。”

这最后的印记即将消散。

“我散去后，这秘境也不能再维持多久了，希望你与你的朋友们能尽快离开这里。年轻的医者，善用这份馈赠……”

最后的意念如同温暖的潮水般退去，留下的是一片浩瀚无垠的理解。

不知过了多久，或许是一瞬，又或许是百年。

笼罩殿宇的磅礴青光渐渐收敛、黯淡，最终完全消失。殿宇中央空空如也，唯有空气里残留着令人心旷神怡的清新气息。

陈千羽静静站立在原地，双眸轻阖，长睫在白皙的脸颊上投下淡淡的阴影。

她周身的气息圆融流转，原本温润平和的元气质地未变，其中却多了一股浩瀚、灵动、仿佛能孕育万物的生机道韵，使她整个人看起来更加沉静深邃，宛如一株历经风雨却内含无尽生机的植株。

她眉心之间，一道极淡的、形似展翅青鸾的印记一闪而逝，随即隐没于光洁的皮肤之下。

在这段时间内，她自然从传承中找到了这只青鸾前辈的生前修为：圣魄境不朽期。

这个修为具体有多高，她对此没有什么概念。

不过这不重要。

她缓缓睁开双眼，眸底似有清澈的青光流转了一瞬，旋即恢复为原本的明净；但那明净之中，却多了几分洞彻生命规律的深邃。

传承已毕，几乎无需消化，已然成为她认知与道基的一部分。不过，她并没有试图当场突破境界。

眼下，离开这个迷宫，才是当务之急。

没有丝毫犹豫，心念一动，触动了怀中一枚简易的传讯符。

\vspace{2em}

一道清晰而简短的信息，传向迷宫之中与她气息相连的同伴：“传承已得，速退。”

迷宫深处，白言冰刚以精妙剑势引偏了一具古老铠甲傀儡的重击，沈瑶的刀光趁机切入其关节缝隙，吴刚则咆哮着用一柄刚捡来的重锤砸碎了它的核心。几人配合无间，正准备查看这房间内是否有遗漏之物，怀中传讯符同时传来了微热与波动。

信息入心。

“得了？”沈瑶挥刀震开傀儡碎片，清丽的脸上第一次露出近乎懵然的诧异，她看向白言冰，重复道，“千羽她……找到了核心？这么快？”

这效率远超所有人最乐观的预估。

白言冰的动作也是一顿，眼中瞬间爆发出惊喜的光彩，但那光彩立刻被更强大的决断力所取代。

他对陈千羽毫无保留的信任让他立刻做出反应：“信千羽！所有人，向入口方向撤离，沿途汇合，准备离开秘境！”他的声音沉稳有力，瞬间压过了刚刚战斗的余响。

没有丝毫留恋，他们甚至不去看那傀儡残骸中是否还有灵材，立刻转身，按照记忆中留下的标记和大致方向，朝着青鸾殿大门处疾退。

途中遇到其他仍在迷茫探索或与零星禁制纠缠的新生势力弟子，无论是云隐别院还是万仙盟的，他们立刻言简意赅地告知：“陈师姐已得传承，速退！”

惊愕浮现在每一张脸上，但无论是出于对白言冰权威的信服，还是对陈千羽机缘的震惊与信任，没有人质疑。求生的本能和对大局的判断让这些弟子迅速压下贪念，加入撤退的洪流。

队伍像滚雪球般壮大，行动却愈发迅捷有序。

另一边，守旧联盟的弟子很快察觉到了异常。

原本如同无头苍蝇般在迷宫中偶尔碰面、彼此戒备的云隐别院和万仙盟的人，忽然都朝着同一个方向快速移动，行为坚决，毫不拖泥带水。

几个小队将相同的消息迅速汇总到周文彬处。

周文彬正为迷宫所困而烦躁，闻讯先是一愣，随即脑中闪电般划过进入殿门时那最后一幕——那个云隐别院的女医修，逆着人流，弯腰捧起了一只该死的、碍事的伤鸟！

他的脸色“唰”一下变得铁青，额角青筋跳动，一股混合着极度嫉妒、被戏弄的羞辱以及计划落空的暴怒直冲顶门，烧得他双眼泛红。

他几乎是从紧咬的牙关里，一字一字地挤出充斥着毒恨的话语：“定是那陈千羽！最后抱着鸟进去的那个女人！她得了传承！”

他猛地转向周围同门，声音因激动而尖利，“不能让他们带走！拦住他们！不惜一切代价！”

然而，迷宫的复杂在此刻成了新生势力最好的掩护。众人撤退路线明确，彼此间依靠简单的信号和“三三制”的小组配合互相掩护、交替断后，行动高效得像演练过无数遍。

守旧联盟队伍本就分散，而在他们自己都不熟悉的迷宫之中捉人更是困难。

仓促间的拦截不是扑空，就是被对方默契的配合轻易化解或摆脱。他们只能眼睁睁看着白言冰、沈瑶等人与从核心区域方向匆匆赶来的陈千羽顺利汇合，然后把她簇拥在中心，汇聚成一股洪流，头也不回地冲向迷宫出口。

“废物！追！出去再说！立刻联系长老，告知情况！”周文彬气得浑身发抖，一脚踹在旁边的发光墙壁上，怒吼着带领手下，也疯狂地追了上去。

那扇厚重的青铜殿门，不知何时已无声洞开，仿佛专为离去者让路。

两派人马前一后，如同两道逆向的激流，冲出青鸾殿，掠过那片依旧笼罩在柔和青光下却已感觉不同的广场，向着秘境入口的方向全力疾驰。

沿途，秘境中那些原本稳定的禁制和生机脉络开始出现紊乱，时而凭空爆发出一股混乱的、充满撕裂感的生机冲击，或是地面骤然长出尖锐的荆棘，给这场生死时速般的逃亡与追击增添了无数变数与凶险。

当那熟悉的七彩漩涡入口终于在望，并且开始出现不稳定的、时缩时胀的波动时，双方人马皆是狼狈不堪，不少人身上添了新伤，但速度却丝毫不减。

新生势力众人护着陈千羽，率先一头扎进了漩涡；守旧联盟弟子红着眼，紧随其后冲入。

\vspace{2em}

身体穿过空间通道的短暂晕眩过后，脚下传来了内域熟悉的、带着碎石的土地触感。

然而，等待新生势力众人的，远远不是约定的“界外不得干涉”的平静，更不是师长接应的欣慰，而是一场早有预谋、卑鄙到极致的血腥伏击！

守旧联盟的弟子在冲出秘境七彩漩涡的第一时间，就毫不犹豫地捏碎了早已扣在掌心的紧急传讯符箓。

忽而，峡谷两侧原本看似平静的山岩、树林后方，骤然爆发出数十道强烈的灵力波动！

“陈千羽！交出青鸾传承，饶你同道不死！”

陈炎那熟悉的、充满灼热霸道的怒吼声如同惊雷炸响，伴随着其声浪，漫天炽热狂暴的火球、烈焰长枪如同陨石雨般轰然砸向刚刚立足未稳的新生势力人群！

几乎同时，刘长靖冰冷凌厉的剑气化为一张巨大的光网，封锁了他们侧翼与后方的退路；丹鼎阁门下高手释放出的墨绿色毒瘴如同活物般从地面蔓延开来，更有隐藏的困阵光华亮起；更有四五名同样散发着真符境强大气息、服饰各异的修士，显然是守旧联盟邀约或麾下的客卿、附属势力长老，带着大量眼神凶狠的低阶修士，如同捕猎的群狼，从峡谷两侧嘶吼着冲杀下来！

瞬息之间，新生势力众人便陷入了十面埋伏的绝境。

“卑鄙！你们违背约定！”峡谷上方，传来了吴东浩惊怒交加的吼声。他与叶牧谦本是按约定前来接应，万万没想到对方竟敢如此公然践踏规则。

两道强横的身影立刻从高处扑下，叶牧谦心念领域如一张无形巨网立即张开，直取最为狂暴的陈炎；吴东浩则啪地一声打开折扇，法术波动悍然拦向刘长靖的剑气网络，力图为下方陷入绝境的弟子们分担最大的压力。

随即，最强的陈炎和刘长靖被叶牧谦和吴东浩两人强行缠住，但更多的真符境摹形期、乃至玄脉境界巅峰的修士们开始向白言冰等人攻去！

“结阵！突围！”白言冰目眦欲裂，看着呼啸而来的致命攻击和密集的敌人，一股热血直冲头顶，狂吼出声。

他毫无保留，见独领域全力展开，那圆融澄明、映照己心的无形力场瞬间扩张，将他与身边众多同伴笼罩。他手中长剑发出清越长吟，剑光化作一道坚韧无匹的游龙，主动迎向那最炽热、最狂暴的几道火焰陨石，剑势浑圆连绵，竟能以一己之力卸开真符境长老的含怒一击。

沈瑶的反应同样快到极致，她清叱一声，周身气息陡然变得飘忽诡谲，人随刀走，化为一道捉摸不定的黑色流光游走在侧翼。刀光配合元气激发的幻术，干扰、误导、甚至将数道袭向修为较低同伴的攻击引偏。

吴刚、陈默等散修在最初的震惊后，爆发出亡命之徒般的狠劲与怒吼。他们自然不能像白言冰与沈瑶两人搞出什么精妙的幻术或领域，全靠血勇与平日里刀头舔血的默契。

散修们也三五成群，迎着扑上来的守旧联盟伏兵厮杀在一起，用身体和简陋的武器构筑起一道混乱却顽强的防线。

陈千羽站在众人保护的中心，脸色微微发白。

她刚刚接受完那浩瀚如海的传承，修为也未曾突破。但看着周围瞬间陷入血火的同伴，她毫不犹豫地抬起双手，十指如穿花蝴蝶般结出一个古朴玄奥的手印——这是传承中最基础、却也最核心的生机引导与护持法门。

春——和畅！

柔和的青色光晕以她为中心，如同水波般迅速扩散开来，笼罩住己方众人。

这光晕是她以自身精纯元气为引，强行模拟、构造出的一个临时性的生机领域，比真正的见独领域更耗费心神与灵力；它无法阻挡飞剑火矢，却蕴含着青鸾传承独有的、滋养万物的奇异生机。

光晕笼罩之下，受伤同伴伤口流血的速度肉眼可见地减缓，剧烈消耗带来的疲惫与灵台昏沉为之一清，仿佛干涸的土地得到了春雨的滋润，一股顽强不屈的生命力从心底被点燃。

然而，敌我力量悬殊得令人绝望！

伏击者不仅人数占优，更有多位真符境长老级人物坐镇指挥、出手。对于绝大多数修为仅在玄脉境或养气境的新生势力弟子而言，真符境修士随手一击的余波，都足以让他们骨断筋折，甚至当场殒命。

整个探索队伍中，唯有白言冰凭借见独境的修为能硬抗真符，沈瑶凭借幻术能勉强在真符境的攻击下周旋、抵挡片刻；其他人则如同狂风中的落叶。

顷刻之间，残酷的现实便撕开了希望的口子。

一名万仙盟的散修躲闪不及，被一道逸散的烈焰擦中，半个身子顿时焦黑，惨叫着倒地；一名云隐别院的弟子奋力格开当面之敌的长剑，却被侧面袭来的毒镖射中大腿，麻痹感瞬间蔓延，随即被乱刀砍中；更有甚者，直接被强横的灵力冲击震得口喷鲜血，倒地不起……

猩红的血液，迅速在峡谷灰黑色的地面上蔓延开来，刺鼻的血腥味混杂着焦糊味、毒瘴的甜腥气，构成一幅地狱般的图景。

“走！不要恋战！快走！”叶牧谦充满压抑怒火的声音，如同惊雷般穿透混乱嘈杂的战场，清晰传入每个苦苦支撑的弟子耳中。

他本人虽然见独圆满、修为高绝，但被陈炎与刘长靖两人死死缠住，根本无法分身救援下方。

“冲出去！”白言冰嘶吼道，他浑身浴血，青衣早已被染成深褐色，分不清是自己的还是敌人的。

他的领域在数名真符境修士的冲击下剧烈波动，却依旧坚韧地支撑着。

他一剑荡开面前之敌，剑势陡然变得惨烈决绝，化为一道锐不可当的锋芒，朝着包围圈相对薄弱的一处，悍然突进！

与其说是斩，不如说是刺！

带着全身重量、见独领域，以及不屈不挠的守护意志，全身撞击过去的刺！

“拦住他！”

两侧的两名守旧联盟修士——玄脉境融贯期——急忙一前一后，试图阻挡白言冰的突围。

没想到，这临时变出的一招，竟然蕴含着不可阻挡的威势。

站在前面的那位修士，白言冰的剑尚未及身，便被那无形的领域一撞！

踉跄着后退几步，几乎要撞到后面那人身上。

两人匆忙防御，却发现白言冰竟然不在他身前！

那他哪去了？

下一刻，那位修士忽然觉得后背被另一个人狠狠地撞了一下。

随即，便是剑尖刺穿胸口的感受。

什么！

白言冰自己见到剑上穿成一串的两个死不瞑目的守旧联盟修士，自己也有点愕然。

他刚才刺前人是伪，闪到后人身后背刺之是真。只是没想到，自己见独领域竟然也能用来撞人？

不过眼下并没有什么思考的余裕。

“包围圈已经撕破，快撤！”

沈瑶默契地紧随其后，刀光幻化成一片令人眼花缭乱的死亡光幕，掩护他的侧翼。

吴刚咆哮着，用以伤换伤的打法，将另外两名试图合围的敌人硬生生逼退。陈千羽咬牙维持着生机光晕，同时搀扶起身边腹部受创、血流不止的陈默。

一行人护着几名尚能行动的伤员，靠着“三三制”瞬间爆发的默契配合、陈千羽不断提供的生机支撑，残存的人马如同离弦之箭，从那缺口处冲出，朝着云霞山、云隐别院的方向亡命奔逃。

叶牧谦与吴东浩亦爆发出全力，心念领域中吴东浩的灵气锋刃纵横交错，且战且退，牢牢挡住身后追兵最凶猛的那些攻击，为弟子们的逃亡争取那宝贵的一线生机。

箭矢尖啸，法术轰鸣，剑气纵横，从身后不断袭来。不断有人中招倒地，发出闷哼或惨叫，又立刻被身旁的同伴拼死拽起，拖拽着继续前进。

鲜血，点点滴滴，连成触目惊心的红线，染红了这条用生命铺就的逃亡之路。

每一次回头，都可能看到熟悉的面孔永远倒下；每一次喘息，都带着铁锈般的血腥味。绝望与悲痛如同冰冷的潮水，一次次试图淹没他们，却又被更强烈的求生欲与愤怒点燃的斗志狠狠压回。

当云隐别院那熟悉的亭台轮廓、以及笼罩其上的淡淡防护阵法光晕终于映入眼帘时，出发时三十多名精锐，只剩下十一二人还能勉强站立，还有几人倒地不起、重伤昏迷。

每个人身上都带着或轻或重的伤，衣衫褴褛，血迹斑斑，气息萎靡到了极点，模样凄惨如同从地狱血池中爬出。

身后，喊杀声与追击的灵力波动在靠近别院范围后，终于渐渐减弱、停止。

守旧联盟的伏兵似乎终究对直接攻击云隐别院本体心存顾忌，依然未敢逾越最后的界线。

最后几步路，几乎是用尽最后一丝力气完成的踉跄。

\vspace{2em}

当终于冲入九宫迷踪的范围，确认那致命的攻击不再袭来时，紧绷到极致的神经骤然松弛。

扑通、扑通……幸存者们如同被抽掉了全身骨头，顿时瘫倒一片，剧烈的喘息声、压抑的痛哼声、劫后余生的哽咽声响成一片。

陈千羽自己也到了强弩之末，识海因过度强行模拟见独领域而阵阵抽痛，元气几乎枯竭。

但她只是擦了擦额角的冷汗与血污，便毫不迟疑地跪坐在重伤员之间，苍白的手指迅速而稳定地检查伤口，掏出随身携带的、所剩无几的药瓶与银针，开始紧急施救。

柔和微弱的青光再次从她指尖亮起，虽然远不及之前，却依旧带着生的希望。

白言冰以剑拄地，单膝跪倒，胸膛剧烈起伏，每一次呼吸都牵扯着身上的伤口，带来火辣辣的疼痛。

他脸上溅满血污，俊朗的面容此刻只剩下鏖战后的疲惫与狼藉，但那双向来沉静的眼眸里，劫后余生的庆幸只是一闪而逝，随即被一种深沉到极致、几乎要喷薄而出的熊熊怒焰所取代。

沈瑶靠坐在不远处一段残破的矮墙边，她左边的肩膀，一道深可见骨的剑伤皮肉翻卷，鲜血仍在汩汩外渗，染红了半边黑衣。她清丽的脸上没有太多表情，只是紧抿着失色的唇，额角渗出细密的冷汗。

她咬着牙，用还能活动的右手，颤抖而坚定地给自己洒上金疮药，然后用牙齿配合，撕下一条相对干净的衣襟，试图包扎。

叶牧谦和吴东浩两人的身影无声地落在庭院中央。这两人倒是没有受到看起来过重的伤，却更显凶险。

真符境只是个玻璃大炮。这个阶段的修士之间打斗，要么无伤，要么陨落。

叶牧谦的一袭青衫染满尘埃与暗红色的血点，头发有些散乱。他的面色冰冷如万载玄铁，没有丝毫表情，目光缓缓扫过庭院中横七竖八、伤痕累累的弟子们，扫过地上那来不及擦拭的斑斑血迹，扫过空气中弥漫的浓重血腥与悲愤气息。

他就这样静静地站着，许久，许久，一言未发。然而，任谁都能感受到，在那死寂般的沉默之下，是压抑到极致、即将冲破一切束缚的滔天怒意与冰冷风暴。

陈千羽用最快的速度，将几名最危重伤员的伤势暂时稳住，吊住了他们最后一口气。做完这一切，她已是摇摇欲坠。她强撑着站起身，对叶牧谦、白言冰等人微微颔首，没有多说什么，眼眸中交织着疲惫、悲恸，以及一丝刚刚领悟的、关于生命重量的明悟。

随即，她便转身，径直走向静室，关上了门。她需要时间，需要一个绝对安静的环境，去整理那浩瀚的传承，去冲击那早已松动的瓶颈，去获得足以应对接下来一切的力量。

沈瑶亦寻了一处静室闭关，她也觉得自己的领域快要成型了。

\vspace{2em}

庭院中的沉寂，持续了整整一日。

次日，当晨光再次洒满云隐别院，两扇静室的门，一前一后地被从内轻轻推开。

沈瑶先从她的静室踏出，身上伤口已经用绷带强行止血，但绷带上渗出、干结的血迹依然触目惊心。

她创新性地将自己的领域与幻术相结合，用元气刻画幻阵。从而，在节省布阵珍贵材料的同时，还能让幻阵随着她的移动而移动。

这幻术若与她的灵动步法相结合，对敌人而言将会是一场灭顶之灾。也正因此，沈瑶的脸上罕见地露出了一丝自得。

她的幻术领域，自名为“千幻流光”。

随后就是陈千羽缓步走了出来。

她整个人的气质，已然发生了天翻地覆的蜕变：之前的温婉宁静依然存在，然而在这底色之上，却多了一份源自生命本源的深邃与厚重，一种洞察生灭、触摸平衡的从容。

她自己也成功地将青鸾传承那浩瀚如星海的感悟，与自身坚定不移的求索根基、与那颗医者仁心相互印证、水乳交融；而她的生机领域，更是令人惊诧。

领域悄然展开的范围内，陈千羽能异常清晰地感知到万物内部生机的流转状态、盈虚变化，乃至那些细微的、可能导致崩溃的失衡之点。

此领域竟然包含两重变化：

其一，自然是先前在追逐战中强行模拟出来的“和畅”形态：陈千羽能以自身精纯元气与领悟的生机之道为引，强行调和、重新平衡那些破损的生机脉络，甚至只要对方尚存一息，她便能以损耗自身元气为代价，为其强行续命；

其二，便是尚未出现过的“萧瑟”形态：与“和畅”恰恰相反，这一形态之下，她能主动去破坏敌人体内的生机平衡，从而大大压制敌人的行动能力，甚至敌人死于这萧瑟秋风之中也不无可能。

这已超越了寻常医术的范畴，是在生死边缘行使恩赐与判决的惊世之能，是连丹鼎阁主王丹那种层次的人物都未必触及到的、近乎于“道”的层次，将陈千羽的医道，推上了一个全新的、令人仰望的峰巅。

白言冰倒是没什么好说的，他忙着帮叶牧谦、吴东浩整理秘境中其余所得，并安葬那些因各种原因把生命留在秘境中的弟子的骨灰。

只是把强行突围的那一招起了个名，叫“孤弓”。

\vspace{2em}

然而，境界的突破、领域的领悟、招式的变幻，并未给云隐别院的任何人带来丝毫喜悦，只有更沉默的压抑。

庭院中尚未散尽的药味与血腥味，同伴们苍白憔悴的脸庞和身上缠裹的绷带，都在无声地诉说着昨日的惨烈与卑鄙。

云隐别院方面一直秉持着克制与发展并重的原则，在诸多摩擦中甚至做出了不少让步，只为争取那来之不易的和平发展时间。

但是，守旧联盟一而再、再而三的倒行逆施，从秘境入口的暗算，到秘境中的偷袭，再到最后这公然违背约定、不惜动用长老级力量进行无耻伏击、造成惨重伤亡的行径，已经彻底、完全地践踏了所有的底线与规则。

昨日那场伏击的每一个细节，同伴倒下的每一个身影，鲜血染红大地的每一寸画面，都在幸存者眼中反复灼烧。

是可忍也，孰不可忍也？

那压抑到死寂的庭院之中，风暴，已然在平静的表象下酝酿到了极致。

沉默啊，沉默啊！

不在沉默中爆发，就在沉默中灭亡！
